\section{Reciprocity kernels of rows}\label{sec:pf:reckernel}

\providecommand{\reckKur}{\widehat{\mathbb{Q}}{}_\ell^{\mathrm{ur}}}
\providecommand{\reckKurp}{\widehat{\mathbb{Q}}{}_p^{\mathrm{ur}}}
\providecommand{\reckCl}{\mathbb{C}_\ell}
\providecommand{\reckRow}{\sfname{Row}}
\providecommand{\reckRowE}{\mathrm{Ev}_{\mathrm{lk}}}
\providecommand{\reckAtAll}{\mathrm{All}}
\providecommand{\reckClosed}{\sfname{Closed}}
\providecommand{\reckEvP}{\mathrm{Ev}_{\mathrm{par}}}
\providecommand{\reckServH}{\mathrm{SK}}
\providecommand{\reckLk}{\mathsf{Lk}}
\providecommand{\reckBd}{\mathsf{Bd}}
\providecommand{\reckm}{\mathfrak{m}}
\providecommand{\reckart}{\operatorname{art}}

\subsection{Introduction}\label{pf:reckernel:sec-intro}

The row package of Part~I (\cref{prop:families:allrows}) contains a row $\mathsf{hServ}$, the \emph{service
row}, which asserts that the selected bundles of every class ``serve'' the designed subclass: the kernel
entries carried by their reciprocity records satisfy a list of local conditions at the inert and level-free
places, and the twisted local systems of the members satisfy the twisted system bound
(\cref{tab:families:fields}). This section proves that row from global Artin reciprocity. The main result is
\cref{pf:reckernel:main}: for every parameter set $p$ and every linkage $\reckLk$ whose linked bundles carry
reciprocity-linked records, the service row holds for the good positions of $p$.

\paragraph{Where the mathematics is.}
The kernel entries of the construction are built from theta functions attached to elliptic curves with complex
multiplication: the lattice $\delta\cdot\mathfrak c\subset\C$ of a non-zero ideal $\mathfrak c$ of the ring of
integers of an imaginary quadratic field $K=\Q(\sqrt{-d})$, scaled so that its discriminant is $1$, with a
marked $2$-torsion point. All local information about the kernels therefore comes from the arithmetic of these
curves at primes $\ell\nmid 6d$. The formal proof derives, from global Artin reciprocity alone (together with
the Chebotarev density theorem, which is proved in the development), the following classical facts:
\begin{enumerate}[label=(\roman*)]
\item (\emph{good models}) a lattice with complex multiplication by $\OO_K$ acquires good reduction over the
  completion $\reckKur$ of the maximal unramified extension of $\Q_\ell$ for every prime $\ell\ge5$ with
  $\ell\nmid d$ (\cref{pf:reckernel:cm2});
\item (\emph{torsion}) prime-to-$\ell$ torsion points of such a good model have integral coordinates in
  $\reckKur$ (\cref{pf:reckernel:cm4});
\item (\emph{Honda}) at primes $\ell$ inert in $K$ the formal group of the good model is isomorphic over the
  ring of integers of $\reckKur$, after a unit twist, to the Lubin--Tate group with logarithm
  $\sum_j X^{\ell^{2j}}/\ell^j$ (\cref{pf:reckernel:cm3}).
\end{enumerate}
From these the development deduces six statements about the CM theta data (algebraicity of Taylor coefficients,
existence of models over the Witt vectors, a curve dictionary, integrality of kernels at inert primes,
admissibility of theta-series data and existence of curves in the required position), and from the six
statements the two components of the service row.

The proof of (i) is complete and classical in spirit, and we give it in full in
\cref{pf:reckernel:sec-cm2}: an ideal-class Artin symbol is extracted from global Artin reciprocity; Kronecker's
congruence for the modular function $j$ is proved from $q$-expansions; together with Chebotarev's theorem it
shows that every inertia element fixes $j(L)$, so that $j(L)\in\reckKur$; and an analysis of kernel polynomials
of CM endomorphisms (via V\'elu's formula) and of torsion points of order prime to $\ell$ shows that the
valuations of $j(L)$ and of $j(L)-1728$ are divisible by $3$ and $2$ respectively, which is exactly what is needed
for a model with unit discriminant over $\reckKur$. The proof of (ii) is given in \cref{pf:reckernel:sec-cm4},
including the identification of the chord-tangent law on the points of the formal group with the formal group law
(\cref{pf:reckernel:fgHom}). For (iii) we give the structure of the proof. For the six statements on theta data and the
two rows we give complete proofs at the level of their main steps: the reduction of the integrality at inert primes to
weak approximation and orthonormal frames (\cref{pf:reckernel:intRed}), the first kernel law and admissibility
(\cref{pf:reckernel:firstLaw,pf:reckernel:admLaw}), the construction of joints and of admitted systems at the sites
(\cref{pf:reckernel:jointsStd,pf:reckernel:jointSite}), the lower bounds at the second real place
(\cref{pf:reckernel:secondPlace}) and the kernel conditions at true columns (\cref{pf:reckernel:atTrue}). The deeper
lemmas used there, which concern the specific theta kernels of the construction, are cited with their statements.

\paragraph{Organisation.}
\Cref{pf:reckernel:sec-row} recalls the row formalism, defines the service condition and reduces the main
theorem to the statement ``served kernels from Artin reciprocity'' (\cref{pf:reckernel:servedArtin}), which is
itself reduced to the three CM statements and six statements on CM theta data. \Cref{pf:reckernel:sec-cmdefs} fixes
the notation for lattices with complex multiplication. \Cref{pf:reckernel:sec-cm2} proves (i),
\cref{pf:reckernel:sec-cm4} proves (ii) and discusses (iii). \Cref{pf:reckernel:sec-theta} treats the six statements
on theta data, including integrality at inert primes (\cref{pf:reckernel:sec-intred}) and admissibility
(\cref{pf:reckernel:sec-admproof}), and \cref{pf:reckernel:sec-rows} the two rows, including joints
(\cref{pf:reckernel:sec-joints}), the second real place (\cref{pf:reckernel:sec-second}) and the kernel conditions at
true columns (\cref{pf:reckernel:sec-truecols}). \Cref{pf:reckernel:sec-remarks} collects remarks on the
formal statements.

\paragraph{Conventions.}
For a prime $\ell$ we write $\reckCl$ for the completion of an algebraic closure of $\Q_\ell$ and $|\cdot|$ for
its absolute value, normalised by $|\ell|=\ell^{-1}$. We write
\[
\reckKur\ :=\ \overline{\Q(\mu_{\ell'})}\ \subset\ \reckCl,
\]
the topological closure in $\reckCl$ of the subfield generated by the roots of unity of order prime to $\ell$;
this is the completion of the maximal unramified extension of $\Q_\ell$, and the formal statements below use
exactly this field (\cref{pf:reckernel:rem-kur}). A \emph{valuation ring} of a field $\Omega$ is a subring
$A\subseteq\Omega$ with $x\in A$ or $x^{-1}\in A$ for every $x\ne0$; $\reckm_A$ is its maximal ideal.

% ---------------------------------------------------------------------------
\subsection{The service row}\label{pf:reckernel:sec-row}

\subsubsection{Rows}
We use the formalism of rows of \cref{def:families:closed,def:families:rows}. Recall the ingredients: contexts $\mathcal K$;
bundle points $\beta$ of $\mathcal K$, each with a datum $D_\beta$ and an object $O_\beta$ over $D_\beta$; slots $Q$
(predicates on the cores of bundle points; we write $\beta\models Q$); parameters $\kappa$, which may \emph{govern} a
context ($\kappa\rhd\mathcal K$) and may be \emph{linked} to a family; bundle types $\reckBd$ (a type $\reckBd(D,O,c)$ for every
datum $D$, object $O$ over $D$ and $c\in\R$); and linkages $\reckLk$ (relations $\reckLk(\mathcal K,\beta,c,b)$ between bundle
points and bundles $b\in\reckBd(D_\beta,O_\beta,c)$). A \emph{selection} $\sigma$ over a family class $\mathcal A$
chooses a bundle $\sigma(F,x,i)$ over every member $(x,i)$, $x>x_1(F)$, of every $F\in\mathcal A$. We write $\reckEvP$
for the filter ``for small exponent, large staged constants and then large scale'' on pairs $(\kappa,\mathcal K)$. For a
property $E(\kappa,D,O,c,b)$ of bundles:
\begin{itemize}
\item $E$ holds \emph{eventually at linked bundles}, written $\reckRowE(Q,\reckLk,E)$, if
  \[
  \reckEvP\bigl(\kappa,\mathcal K\mapsto\forall\beta\models Q\ \forall c\ \forall b:\
  \reckLk(\mathcal K,\beta,c,b)\Rightarrow E(\kappa,D_\beta,O_\beta,c,b)\bigr);
  \]
\item a bundle $b$ over $(D,O)$ is \emph{built} for $\kappa$ if $(D,O,b)=(D_\beta,O_\beta,b')$ for a bundle point
  $\beta$ of a context governed by $\kappa$ with $\beta\models Q$ and $\reckLk(\mathcal K,\beta,c,b')$;
\item $\sigma$ is an \emph{$E$-selection} for $\mathcal A$ if every $F\in\mathcal A$ has a parameter $\kappa$ linked to
  it such that $F$ satisfies the duties of the construction and every selected bundle $\sigma(F,x,i)$ is built for
  $\kappa$ and satisfies $E(\kappa,\cdot)$;
\item $\reckRow(Q,\reckLk,H)$ means: there is $E$ with $\reckRowE(Q,\reckLk,E)$ such that $H(\mathcal A,\sigma)$ holds
  for every class $\mathcal A$ and every $E$-selection $\sigma$ for $\mathcal A$.
\end{itemize}
For a property $R$ of bundles, $\reckAtAll(R)(\mathcal A,\sigma)$ means that $R(\sigma(F,x,i))$ holds for all
$F\in\mathcal A$, $x>x_1(F)$ and indices $i$; $\reckClosed(\mathcal C,P)$ is the closedness of
\cref{def:families:closed}; and $\mathcal D_3(\mathcal A)\subseteq\mathcal A$ is the designed subclass of
\cref{def:families:closed} (\lref{GlobalTheory.Cp.d3_a}: every member of $\mathcal D_3(\mathcal A)$ is a member of
$\mathcal A$).

We use four formal properties of rows. The first is part of this section; the other three are shared lemmas with
proofs of a few lines, which we include for the reader's convenience.

\begin{lemma}[Calculus of rows]\label{pf:reckernel:calculus}
Let $Q,Q'$ be slots and $\reckLk,\reckLk'$ linkages for bundle types $\reckBd,\reckBd'$.
\begin{enumerate}
\item If $\reckRow(Q,\reckLk,H)$ and $\reckRow(Q,\reckLk,H')$, then $\reckRow(Q,\reckLk,H\wedge H')$.
\item If $\reckRow(Q,\reckLk,H)$, and $H(\mathcal A,\sigma)\Rightarrow H'(\mathcal A,\sigma)$ whenever $\sigma$ is a
  selection for the trivial property, then $\reckRow(Q,\reckLk,H')$.
\item If $E(\kappa,D_\beta,O_\beta,c,b)$ holds whenever $\kappa\rhd\mathcal K$, $\beta\models Q$ and $\reckLk(\mathcal K,\beta,c,b)$, then
  $\reckRowE(Q,\reckLk,E)$.
\item \textup{(Transport.)} Let $\pi$ map $\reckBd(D,O,c)$ to $\reckBd'(D,O,c)$, assume $Q\subseteq Q'$ and that
  $\beta\models Q$ and $\reckLk(\mathcal K,\beta,c,b)$ imply $\reckLk'(\mathcal K,\beta,c,\pi(b))$. Then $\reckRow(Q',\reckLk',H)$
  implies $\reckRow\bigl(Q,\reckLk,(\mathcal A,\sigma)\mapsto H(\mathcal A,\pi\circ\sigma)\bigr)$.
\end{enumerate}
\formallabel{GlobalTheory.Fc.Row.and}
\end{lemma}

\begin{proof}
(1) If $E,E'$ witness the two rows, $E\wedge E'$ witnesses the conjunction: the filter $\reckEvP$ is closed under
finite intersections, and an $(E\wedge E')$-selection is an $E$-selection and an $E'$-selection. (2) The same $E$
works, since an $E$-selection is a selection for the trivial property (\lref{GlobalTheory.Fc.Row.imp}: rows
are monotone in this sense). (3) The filter statement holds for all parameters and contexts
(\lref{GlobalTheory.Fc.RowE.of_all}: statement (3)). (4) If $E$ witnesses $\reckRow(Q',\reckLk',H)$, then
$E\circ\pi$ witnesses the transported row: its filter statement follows from that of $E$, because linked bundles at
$Q$-points map to $\reckLk'$-linked bundles at $Q'$-points; and a bundle built for $(Q,\reckLk,\kappa)$ maps to one
built for $(Q',\reckLk',\kappa)$, so that $\pi\circ\sigma$ is an $E$-selection whenever $\sigma$ is an
$(E\circ\pi)$-selection (\lref{GlobalTheory.Fc.Row.map}: statement (4)).
\end{proof}

\subsubsection{Served kernels}
Every bundle $b$ of the bundle type of the package (\cref{def:families:rows}) carries a reciprocity record $R_b$, a
record over the extra data $\mathrm{ex}_O$ of the object $O$ over which $b$ lies. A record $R$ contains tied theta data
and in particular a \emph{tied kernel entry} $t_R$. A tied kernel entry $t$ over $O$ consists of an ideal
$\mathfrak c$ of $\OO_K$, a number $\delta$, a CM base (\cref{pf:reckernel:sec-theta}), a character $\chi$ and a
\emph{kernel entry} $E_t$, together with many compatibilities. A kernel entry $E$ consists, for every word $u$, tuple
$\mu$, place $w$ of the coefficient field over a prime $p_w$ and level $i$, of multivariable power series
$\mathrm{Th}_E(u,\mu,w,i,c,c')$ with coefficients in $\C_{p_w}$ (the \emph{theta kernels}), spaces with evaluation
maps, local projections and non-zero \emph{pair scalars} $\mathrm{un}_E(u,\mu,w,i,c,c')\in\C_{p_w}$, related by a
kernel formula; we shall not need their details.

\begin{definition}[Served kernels]\label{pf:reckernel:def-servH}
Let $\reckBd$ be a bundle type and $t$ an assignment of a tied kernel entry $t(b)$ over $O$ to every bundle
$b\in\reckBd(D,O,c)$.
\begin{enumerate}
\item A family $F$ satisfies the \emph{twisted-system condition} $\mathrm{TwU}(F)$ if there is $x_0$ such that for all
  $x\ge x_0$ with $x>x_1(F)$ and all indices $i$ there is a choice of lattice components for which the
  Weil--Jacquet--Shalika compatibility and the twisted system bound (with constant $1$) hold at the member $(x,i)$.
\item For extra data $\mathrm{ex}$ and a kernel entry $E$, the \emph{kernel conditions} $\mathrm{KC}(\mathrm{ex},E)$ are:
  the pair scalars of $E$ are units at the level-free places; the transition maps of the twisted chains carry the
  rational spaces of $E$ at level $i$ into those at level $i+1$; the chain maps of the members and of the twisted
  chains preserve the filtrations of the forms; the stabilisers at level $0$ are trivial; and, for suitable auxiliary
  data (a scalar, power series and scalars at each place), the derived entry satisfies the even conditions at the
  inert places and at the places of two further residue types, a lower bound for its orbit valuations at the places of
  residue class zero, jet integrality (integrality of the theta kernels at level-free places away from $d$ when
  $N_iN_{i+1}$ is even), the shape condition for the theta kernels at the ramified prime $d$, and the shape
  inequalities between the sizes, rows and weights of the entry.
\item $\reckServH_t(\mathcal A,\sigma)$ is the conjunction of the two statements
  \[
  \reckClosed(\mathcal A,\mathrm{TwU})\qquad\text{and}\qquad
  \reckAtAll\bigl(b\mapsto\mathrm{KC}(\mathrm{ex}_O,E_{t(b)})\bigr)(\mathcal A,\sigma).
  \]
\end{enumerate}
\formallabel{GlobalTheory.TwU.RowTwU}
\formallabel{GlobalTheory.Pdv.At}
\formallabel{GlobalTheory.Svd.ServH}
\end{definition}

The service row of the package is phrased with the designed subclass. For the bundle type of the package, a selection
$\sigma$ \emph{serves} $\mathcal A$, written $\mathrm{Serves}(\mathcal A,\sigma)$, if the two statements of
\cref{pf:reckernel:def-servH}(3), with $t(b)=t_{R_b}$, hold for $\mathcal D_3(\mathcal A)$ in place of $\mathcal A$,
the second one in the closed form ``closedly, for every member and every $x>x_1$''
(\lref{GlobalTheory.Cp.SelectionServed}: the formal definition of serving). The \emph{service row} of a slot $Q$
and a linkage $\reckLk$ is $\mathrm{Serv}(Q,\reckLk):=\reckRow(Q,\reckLk,\mathrm{Serves})$
(\lref{GlobalTheory.Fc.R1.SlotServiceCondition}: the row $\mathsf{hServ}$ of the package).

\begin{lemma}\label{pf:reckernel:kerServ}
For every slot $Q$ and every linkage $\reckLk$ for the bundle type of the package,
$\reckRow\bigl(Q,\reckLk,\reckServH_{b\mapsto t_{R_b}}\bigr)$ implies $\mathrm{Serv}(Q,\reckLk)$.
\formallabel{GlobalTheory.Fc.R1.row_hServ_of_kerServ_alt}
\end{lemma}

\begin{proof}
By \cref{pf:reckernel:calculus}(2) it suffices to show that $\reckServH_{b\mapsto t_{R_b}}(\mathcal A,\sigma)$ implies
$\mathrm{Serves}(\mathcal A,\sigma)$. Let $\eps_1$ and the thresholds $\mathsf T(\eps)$ witness the closedness of
$\mathrm{TwU}$ on $\mathcal A$. Since $\mathcal D_3(\mathcal A)\subseteq\mathcal A$, the same $\eps_1$ and $\mathsf T(\eps)$
witness its closedness on $\mathcal D_3(\mathcal A)$, which is the first component. For the second component take the
same $\eps_1$ and $\mathsf T(\eps)$: the kernel conditions hold at all selected bundles of all members of $\mathcal A$, hence
of $\mathcal D_3(\mathcal A)$, without any largeness assumption.
\end{proof}

\subsubsection{The reciprocity linkage and the standard tie}
Let $p$ be a parameter set (\cref{def:families:rows}); its type is formed using the statement $\sfname{RecTExists}$ of
Part~I, which plays no other role here. It determines a slot system and three slots
$\mathcal G_p\subseteq\mathcal B_p\subseteq\mathcal B^\circ_p$: good positions, built positions with explicit kernel
entries, and built positions. A core is a \emph{position with a part} if it arises from a column choice and a later
stage which admit a \emph{part}: a column pin, a column datum, and kernel data matching those of the column choice; we
write $\mathcal P$ for this slot. Let $\mathcal R$ be the bundle type of reciprocity records: $\mathcal R(D,O,c)$ is
the type of records over $\mathrm{ex}_O$. The \emph{reciprocity linkage} $\reckLk^{\mathrm{rec}}_p$ for $\mathcal R$
relates a bundle point $\beta$ and a record $R$ iff $(D_\beta,O_\beta,R)=(D_q,O_q,R')$ for a witnessed point $q$ of the
slot system of $p$ and a record $R'$ over $q$ which is compatible with $q$ in the extended sense
(\lref{GlobalTheory.Mem.Top.Bl.RecordCompatibleExtended}: among other conditions, the tied kernel entry of $R'$ is
the standard one at the column choice of $q$, in the sense of \lref{GlobalTheory.Svd.StdAt}).

\begin{definition}[Standard tie]\label{pf:reckernel:def-stdTie}
A tied kernel entry $t$ over an object $O$ over a datum $D$ is \emph{standard-tied}, $\mathrm{Std}(D,O,t)$, if
$(D,O,t)=(D',O',t')$ where $D'$ is the datum of a column construction with laws, initial law and later stage, $O'$
an object over $D'$, and $t'$ the standard kernel entry at these data.
\formallabel{GlobalTheory.Svd.StdTie}
\end{definition}

\begin{lemma}\label{pf:reckernel:linkStd}
For every parameter set $p$ and every slot $Q$,
\[
\reckRowE\bigl(Q,\ \reckLk^{\mathrm{rec}}_p,\ (D,O,c,R)\mapsto\mathrm{Std}(D,O,t_R)\bigr).
\]
\formallabel{GlobalTheory.Svd.LinkStd.linkR_stdTie_entry}
\formallabel{GlobalTheory.Svd.LinkStd.stdTie_of_eq_alt}
\end{lemma}

\begin{proof}
By \cref{pf:reckernel:calculus}(3) it suffices to show $\mathrm{Std}(D_\beta,O_\beta,t_R)$ for every linked pair
$(\beta,R)$. Unpacking the linkage gives a witnessed point $q$ and a record $R'$ compatible with $q$ in the extended
sense with $(D_\beta,O_\beta,R)=(D_q,O_q,R')$. The compatibility says that $t_{R'}$ is the standard kernel entry at the
column choice of $q$, whose datum is $D_q$; this is $\mathrm{Std}(D_q,O_q,t_{R'})$, and the claim follows by transport
along the equality of triples.
\end{proof}

\begin{lemma}\label{pf:reckernel:positions}
For every parameter set $p$, $\mathcal G_p\subseteq\mathcal B_p\subseteq\mathcal B^\circ_p\subseteq\mathcal P$.
\formallabel{GlobalTheory.Mem.Top.Bl.IsGoodE.isBuiltE}
\formallabel{GlobalTheory.Mem.Top.Bl.IsBuiltE.isBuilt}
\end{lemma}

\begin{proof}
A good position is a built position with explicit kernel entries satisfying an additional goodness condition, and the
latter is in particular a built position. Every built position is a position with a part
(\lref{GlobalTheory.Svd.builtC_hasPart}: the built record provides the column pin, the column datum and the
matching kernel data).
\end{proof}

\subsubsection{The main theorem}

\begin{theorem}[Service row]\label{pf:reckernel:main}
Assume global Artin reciprocity (\cref{pf:reckernel:artinLit}) and $\sfname{RecTExists}$. Let $p$ be a parameter set
and $\reckLk$ a linkage for the bundle type of the package such that $\reckLk(\mathcal K,\beta,c,b)$ implies
$\reckLk^{\mathrm{rec}}_p(\mathcal K,\beta,c,R_b)$ for all $\mathcal K,\beta,c,b$. Then the service row $\mathrm{Serv}(\mathcal G_p,\reckLk)$
holds.
\formallabel{GlobalTheory.T.hServ_of_linkR}
\formallabel{GlobalTheory.row_hServ_of_link_alt}
\end{theorem}

The theorem is used in the proof of the row package (\cref{prop:families:allrows}), where the hypothesis on $\reckLk$
is supplied for the bundle linkage of $p$
(\lref{GlobalTheory.Fc.AllAtMemberStage0.linkR_of_isLinked'}: bundle-linked bundles have reciprocity-linked
records).

\begin{theorem}[Reciprocity kernels]\label{pf:reckernel:recKernel}
Assume global Artin reciprocity and $\sfname{RecTExists}$, and let $p$ be a parameter set. Then
$\reckRow\bigl(\mathcal G_p,\ \reckLk^{\mathrm{rec}}_p,\ \reckServH_{R\mapsto t_R}\bigr)$ for the bundle type $\mathcal R$.
\formallabel{GlobalTheory.T.reciprocity_kernel_of_artin}
\formallabel{GlobalTheory.Svd.servH_row_linkR_alt}
\end{theorem}

\begin{theorem}[Served kernels from Artin reciprocity]\label{pf:reckernel:servedArtin}
Assume global Artin reciprocity. Let $\reckBd$ be a bundle type, $Q\subseteq\mathcal P$ a slot, $\reckLk$ a linkage,
and $t$ an assignment of tied kernel entries to bundles such that
$\reckRowE\bigl(Q,\reckLk,(D,O,c,b)\mapsto\mathrm{Std}(D,O,t(b))\bigr)$. Then $\reckRow(Q,\reckLk,\reckServH_t)$.
\formallabel{GlobalTheory.T.bundle_bothParts_of_standardKernel_of_globalArtinReciprocity}
\formallabel{GlobalTheory.Svd.servH_row_of_artin}
\end{theorem}

\begin{proof}[Proof of \cref{pf:reckernel:recKernel} from \cref{pf:reckernel:servedArtin}]
Apply \cref{pf:reckernel:servedArtin} to $\reckBd=\mathcal R$, $Q=\mathcal G_p$, $\reckLk=\reckLk^{\mathrm{rec}}_p$ and
$t(R)=t_R$. The inclusion $\mathcal G_p\subseteq\mathcal P$ is \cref{pf:reckernel:positions}, and the hypothesis on
$\reckRowE$ is \cref{pf:reckernel:linkStd}.
\end{proof}

\begin{proof}[Proof of \cref{pf:reckernel:main}]
By \cref{pf:reckernel:recKernel} we have the row of reciprocity kernels for records. Apply
\cref{pf:reckernel:calculus}(4) with $Q=Q'=\mathcal G_p$, $\reckLk'=\reckLk^{\mathrm{rec}}_p$ and $\pi(b)=R_b$; the
hypothesis of (4) is the hypothesis on $\reckLk$. Since $\reckServH_{R\mapsto t_R}(\mathcal A,\pi\circ\sigma)$ is
literally $\reckServH_{b\mapsto t_{R_b}}(\mathcal A,\sigma)$, we obtain
$\reckRow\bigl(\mathcal G_p,\reckLk,\reckServH_{b\mapsto t_{R_b}}\bigr)$, and \cref{pf:reckernel:kerServ} concludes.
\end{proof}

\subsubsection{Reduction to statements on CM data}
\Cref{pf:reckernel:servedArtin} is deduced from three statements on elliptic curves with complex multiplication, the
good-model statement (GM), the Lubin--Tate statement (LT) and the torsion statement (TI) of
\cref{pf:reckernel:def-CM}, and from six statements on the CM theta data of the construction: algebraicity (Alg),
models over the Witt vectors (Mod), the curve dictionary (Dict), integrality at inert primes (Int), admissibility
below maximal level (Adm) and existence of curves in position (Ex) (\cref{pf:reckernel:sec-theta}).

\begin{theorem}[Served kernels from six statements]\label{pf:reckernel:six}
Assume \textup{(Alg)}, \textup{(Mod)}, \textup{(Dict)}, \textup{(Int)}, \textup{(Adm)} and \textup{(Ex)}. Then the
conclusion of \cref{pf:reckernel:servedArtin} holds: for $\reckBd,Q,\reckLk,t$ as there,
$\reckRow(Q,\reckLk,\reckServH_t)$.
\formallabel{GlobalTheory.T.bundle_bothParts_of_standardKernel_and_six_hypotheses}
\formallabel{GlobalTheory.Svd.servH_row}
\end{theorem}

\begin{proof}
Let $\mathrm{LM}(D,O)$ be the local-model property of \cref{pf:reckernel:def-modelsAt}. By
\cref{pf:reckernel:localModels} (which uses (Adm) and (Ex)), $\mathrm{LM}(D_\beta,O_\beta)$ holds at every bundle point
$\beta$ in a position with a part, in particular at every $\beta\models Q$; by \cref{pf:reckernel:calculus}(3),
$\reckRowE\bigl(Q,\reckLk,(D,O,c,b)\mapsto\mathrm{LM}(D,O)\bigr)$. Then \cref{pf:reckernel:rowTwU} gives
$\reckRow\bigl(Q,\reckLk,(\mathcal A,\sigma)\mapsto\reckClosed(\mathcal A,\mathrm{TwU})\bigr)$, and
\cref{pf:reckernel:pdv} (which uses (Alg), (Mod), (Dict), (Int) and the standard tie) gives
$\reckRow\bigl(Q,\reckLk,\reckAtAll(b\mapsto\mathrm{KC}(\mathrm{ex}_O,E_{t(b)}))\bigr)$. Their conjunction
(\cref{pf:reckernel:calculus}(1)) is $\reckRow(Q,\reckLk,\reckServH_t)$.
\end{proof}

\begin{proof}[Proof of \cref{pf:reckernel:servedArtin}]
Global Artin reciprocity implies (GM) (\cref{pf:reckernel:cm2}); (LT) and (TI) hold unconditionally
(\cref{pf:reckernel:cm3,pf:reckernel:cm4}). From these, \cref{pf:reckernel:threeFromCM} gives (Alg), (Mod) and (Dict),
\cref{pf:reckernel:intInert} gives (Int), \cref{pf:reckernel:adm} gives (Adm), and \cref{pf:reckernel:curveEx} gives
(Ex). Now apply \cref{pf:reckernel:six}.
\end{proof}

% ---------------------------------------------------------------------------
\subsection{Lattices with complex multiplication}\label{pf:reckernel:sec-cmdefs}

\paragraph{Lattices.}
A \emph{lattice} $L\subset\C$ is given by a pair of $\R$-linearly independent periods. We write
$g_2(L)=60\sum'_{w\in L}w^{-4}$, $g_3(L)=140\sum'_{w\in L}w^{-6}$, $\Delta(L)=g_2(L)^3-27g_3(L)^2\ne0$,
$\wp=\wp_L$ for the Weierstrass function and
\[
j(L)=\frac{1728\,g_2(L)^3}{g_2(L)^3-27g_3(L)^2}.
\]
For $c\in\C^\times$ one has $g_2(cL)=c^{-4}g_2(L)$, $g_3(cL)=c^{-6}g_3(L)$, $\wp_{cL}(cz)=c^{-2}\wp_L(z)$ and
$j(cL)=j(L)$. Over any field $F$ of characteristic $\ne2$ and $g_2,g_3\in F$ we write $W_{g_2,g_3}$ for the
Weierstrass curve $y^2=x^3-\tfrac{g_2}{4}x-\tfrac{g_3}{4}$; its discriminant is $g_2^3-27g_3^2$. For a lattice
$L$ and $z\notin L$ the point $(\wp(z),\tfrac12\wp'(z))$ lies on $W_{g_2(L),g_3(L)}$.

\paragraph{Complex multiplication.}
Throughout, $d$ is a squarefree natural number with $d\equiv3\pmod 4$, $K_d=\Q(\sqrt{-d})\subset\C$ with
$\sqrt{-d}=i\sqrt d$, and $\omega_d=(1+\sqrt{-d})/2$, so that $\OO_{K_d}=\Z[\omega_d]$ and
$\omega_d^2-\omega_d+\tfrac{d+1}{4}=0$. A lattice $L$ \emph{has complex multiplication by $\omega_d$} if
$\omega_dL\subseteq L$. Then $\alpha L\subseteq L$ for every $\alpha\in\OO_{K_d}$, and $cL$ has complex
multiplication by $\omega_d$ for every $c\ne0$. We write $\Omega_d$ for the algebraic closure of $K_d$ in $\C$.
We use the following shared facts: for such $L$, $j(L)$ is an algebraic integer
(\lref{ReciprocityTheory.ChartRestriction.L1.cm1}; classically \cite[Ch.~II, \S6]{Silverman1994}), hence
$|\iota(j(L))|\le1$ for every prime $\ell$ and every field isomorphism $\iota\colon\C\to\reckCl$
(\lref{ReciprocityTheory.ChartRestriction.L1.jInt}).

\begin{definition}[Good models; the statements (GM), (LT), (TI)]\label{pf:reckernel:def-CM}
Let $\ell$ be a prime and $\iota\colon\C\to\reckCl$ a field isomorphism.
\begin{enumerate}
\item A pair $(g_2',g_3')\in\reckCl^2$ is a \emph{good model} of $L$ at $(\ell,\iota)$ if
  $g_2',g_3'\in\reckKur$, $|g_2'|\le1$, $|g_3'|\le1$, $|g_2'^3-27g_3'^2|=1$, and there is $\lambda\in\reckCl^\times$
  with $g_2'=\lambda^4\iota(g_2(L))$ and $g_3'=\lambda^6\iota(g_3(L))$.
\item \textup{(GM)} (\emph{good models}): for all $d$ as above, all $L$ with complex multiplication by $\omega_d$, all primes
  $\ell\ge5$ with $\ell\nmid d$ and all $\iota$, $L$ has a good model at $(\ell,\iota)$.
\item \textup{(LT)} (\emph{Lubin--Tate twist}): for $d,L,\ell,\iota$ as in (2) such that moreover $-d$ is not a square modulo $\ell$, and for
  every good model $(g_2',g_3')$ of $L$ at $(\ell,\iota)$, there is $\eps\in\reckKur$ with $|\eps|=1$ such that all
  coefficients of the power series
  \[
  u_{g_2',g_3'}\bigl(\eps\,\lambda_\ell(X)\bigr),\qquad \lambda_\ell(X)=\sum_{j\ge0}\frac{X^{\ell^{2j}}}{\ell^j},
  \]
  lie in $\reckKur$ and have absolute value $\le1$. Here $u_{g_2,g_3}(z)\in F[[z]]$ is the power series
  $z\cdot\frac{z^2\wp(z)}{-z^3\wp'(z)/2}$, where $z^2\wp(z)=1+\sum_{k\ge2}c_kz^{2k}$ and
  $-z^3\wp'(z)/2=1-\sum_{k\ge2}(k-1)c_kz^{2k}$ are computed from $c_2=g_2/20$, $c_3=g_3/28$ and the recursion
  $c_k=\frac{3}{(2k+1)(k-3)}\sum_{m=2}^{k-2}c_mc_{k-m}$; over $\C$, $u(z)=-2\wp(z)/\wp'(z)$.
\item \textup{(TI)} (\emph{torsion integrality}): for $d,L,\ell,\iota$ as in (2), every $\lambda\in\reckCl^\times$ such that
  $(\lambda^4\iota(g_2(L)),\lambda^6\iota(g_3(L)))$ is a good model, and all $z\in\C$, $N\ge1$ with $\ell\nmid N$,
  $Nz\in L$ and $z\notin L$: the numbers $\lambda^2\iota(\wp(z))$ and $\lambda^3\iota(\wp'(z))$ lie in $\reckKur$ and
  have absolute value $\le1$.
\end{enumerate}
\formallabel{Formalism.L2.Th1.CM.IsGoodModel}
\formallabel{Formalism.L2.Th1.CM.CM2}
\formallabel{Formalism.L2.Th1.CM.CM3}
\formallabel{Formalism.L2.Th1.CM.CM4}
\formallabel{Formalism.L2.Th1.CM.uSer}
\formallabel{Formalism.L2.Th1.CM.ltLog}
\end{definition}

The series $u_{g_2,g_3}$ expresses the parameter $t=-x/y$ of $W_{g_2,g_3}$ at the origin as a power series in the
formal logarithm $z$ (normalised by $dz=dx/2y$), and $\lambda_\ell$ is the logarithm of the Lubin--Tate formal
group of the unramified quadratic extension of $\Q_\ell$ with uniformiser $\ell$ \cite{LubinTate1965}. Thus
(LT) says that, after a unit twist, the formal group of the good model is isomorphic over the ring of
integers of $\reckKur$ to that Lubin--Tate group, in accordance with Honda's theory \cite{Honda1970}.

\begin{remark}\label{pf:reckernel:rem-kur}
The field called ``maximal unramified extension'' in these statements is the closure $\reckKur$ of
$\Q(\mu_{\ell'})$ in $\reckCl$, i.e.\ the \emph{completion} of $\Q_\ell^{\mathrm{ur}}$, not $\Q_\ell^{\mathrm{ur}}$
itself. Its value group is $\ell^\Z$ (\lref{ReciprocityTheory.ChartRestriction.L1.kur_val}), it is complete,
and its residue field is $\overline{\F}_\ell$. The classical statement \cite{Deuring1941,SerreTate1968} gives good
reduction over $\Q_\ell^{\mathrm{ur}}$ itself; the formal statement (GM) is the corresponding statement
over the completion.
\end{remark}

\begin{littheorem}[Global Artin reciprocity {\cite{Tate1967gcft,Neukirch1999}}]\label{pf:reckernel:artinLit}
Let $E$ be a number field and $\Omega$ an algebraic closure of $E$. There is a homomorphism
$\reckart\colon\A_E^\times\to\Gal(\Omega/E)^{\mathrm{ab}}$ (topological abelianisation) which is an open
quotient map, whose kernel is the closure of $E^\times\cdot(\A_{E,\infty}^\times)^\circ$, and such that for every
finite place $v$ of $E$ and every valuation ring $A$ of $\Omega$ over $v$: the image of the local units
$\OO_v^\times\subset E_v^\times\subset\A_E^\times$ is the image of the inertia group of $A$; the images of the
uniformisers of $E_v$ are the images of the inverses of the arithmetic Frobenius elements at $A$; and the closure
of the image of $E_v^\times$ is the image of the decomposition group of $A$.
\formal{GlobalTheory.Published.GlobalArtinReciprocity}
\end{littheorem}

% ---------------------------------------------------------------------------
\subsection{Good models from global Artin reciprocity}\label{pf:reckernel:sec-cm2}

The aim of this subsection is:

\begin{theorem}[Good models]\label{pf:reckernel:cm2}
Global Artin reciprocity implies \textup{(GM)}.
\formallabel{GlobalTheory.Corev.CmTwoClosing.cmGoodModel_of_globalArtinReciprocity}
\end{theorem}

The proof is assembled at the end of the subsection. For a finite place $v$ of a number field $E$, an algebraic
closure $\Omega$ of $E$ and a valuation ring $A$ of $\Omega$, we say that $A$ \emph{lies above} $v$ if
$\OO_E\cap\reckm_A=v$; an element $g\in\Gal(\Omega/E)$ is an \emph{inertia element at $A$} if $gA=A$ and
$gx-x\in\reckm_A$ for all $x\in A$; it is an \emph{arithmetic Frobenius at $A$} (over $v$) if $gA=A$ and
$gx-x^{\Nm v}\in\reckm_A$ for all $x\in A$, where $\Nm v=|\OO_E/v|$.

\subsubsection{The Artin symbol on ideal classes}

\begin{lemma}\label{pf:reckernel:classHom}
For every number field $E$ there is a homomorphism $\mathrm{cl}\colon\A_E^\times\to\operatorname{Cl}(\OO_E)$ with open kernel
containing $E^\times\cdot(\A_{E,\infty}^\times)^\circ$, trivial on the units of $E_v$ for every finite $v$, and
mapping every uniformiser of $E_v$ to the class of $v$.
\formallabel{GlobalTheory.T.exists_ideleClassGroupHom_open_kernel}
\formallabel{GlobalTheory.exists_ideleClassHom}
\end{lemma}

\begin{proof}
Take $\mathrm{cl}(x)=$ the class of the fractional ideal $\prod_{v<\infty}v^{\ord_v(x_v)}$
(\lref{GlobalTheory.Mem.UnitIndex.classHom}). Its kernel contains the open subgroup
$\prod_{v\mid\infty}E_v^\times\times\prod_{v<\infty}\OO_v^\times$ and the principal ideles, and a uniformiser at
$v$ maps to $[v]$.
\end{proof}

\begin{lemma}[Ideal-class Artin symbol]\label{pf:reckernel:artinSymbol}
Assume global Artin reciprocity. For every number field $E$ and algebraic closure $\Omega$ of $E$ there is a
homomorphism $a_E\colon\Gal(\Omega/E)\to\operatorname{Cl}(\OO_E)$ with open kernel such that $a_E(g)=1$ for every inertia element
$g$ at any valuation ring of $\Omega$, and $a_E(\varphi)=[v]$ for every finite place $v$, every valuation ring $A$
above $v$ and every arithmetic Frobenius $\varphi$ at $A$.
\formallabel{GlobalTheory.T.exists_idealClass_artinSymbol_of_globalArtin}
\formallabel{GlobalTheory.IdealArtinSymbol.exists_artinSymbol}
\formallabel{GlobalTheory.T.inertiaElement_eq_one_or_liesAbove_finitePlace}
\end{lemma}

\begin{proof}
Let $\reckart$ be as in \cref{pf:reckernel:artinLit} and $\mathrm{cl}$ as in \cref{pf:reckernel:classHom}. Since
$\ker\mathrm{cl}$ is open, it is closed, so it contains the closure $\ker\reckart$ of
$E^\times(\A_{E,\infty}^\times)^\circ$. As $\reckart$ is a quotient map, $\mathrm{cl}=\rho\circ\reckart$ for a
homomorphism $\rho\colon\Gal(\Omega/E)^{\mathrm{ab}}\to\operatorname{Cl}(\OO_E)$ with open kernel. Put $\psi(g)=\rho(\bar g)$ for
the image $\bar g$ of $g$ in the abelianisation, and $a_E=\psi^{-1}$ (pointwise inverse); its kernel is the
preimage of $\ker\rho$, which is open. If $g$ is an inertia element at $A$, then either $g=1$ or $A$ lies above a
finite place $v$ (\lref{GlobalTheory.inertiaEl_eq_one_or_liesAbove}: a non-trivial inertia element
stabilises a valuation ring lying above a finite place); in the second case
$\bar g=\reckart(x)$ for a unit $x$ of $E_v$, so $\psi(g)=\mathrm{cl}(x)=1$. If $\varphi$ is an arithmetic
Frobenius at $A$ above $v$, then $\overline{\varphi^{-1}}=\reckart(x)$ for a uniformiser $x$ of $E_v$, so
$a_E(\varphi)=\psi(\varphi^{-1})=\mathrm{cl}(x)=[v]$.
\end{proof}

\subsubsection{Chebotarev at places of degree one}
We use the Chebotarev density theorem \cite{Chebotarev1926} in the following form, which is proved in the
development (\lref{GlobalTheory.Published.cheb_holds}): \emph{for a number field $F$, an algebraic closure
$\Omega$, an open normal subgroup $U\subseteq\Gal(\Omega/F)$, $g\in\Gal(\Omega/F)$ and a finite set $S$ of finite
places, there are a place $v\notin S$ of inertia degree $1$ and ramification index $1$ over $\Q$, a valuation ring
$B$ above $v$ whose inertia elements lie in $U$, and an arithmetic Frobenius $\sigma$ at $B$ with $\sigma
g^{-1}\in U$.}

\begin{lemma}\label{pf:reckernel:cheb}
Let $E$ be a number field, $\Omega$ an algebraic closure, $S$ a finite set of finite places, $H$ an open normal
subgroup of $\Gal(\Omega/E)$ containing every inertia element at every valuation ring of $\Omega$ above a place outside
$S$, and $\sigma\in\Gal(\Omega/E)$. Then there are a finite place $v\notin S$ whose
absolute norm $\Nm v$ is a prime number, a valuation ring $A$ of $\Omega$ above $v$ and an arithmetic Frobenius
$\varphi$ at $A$ with $\sigma^{-1}\varphi\in H$.
\formallabel{GlobalTheory.T.chebotarev_degreeOne_valuationSubrings}
\formallabel{GlobalTheory.chebOmSt_holds}
\end{lemma}

\begin{proof}
Apply the Chebotarev statement with $U=H$, $g=\sigma$ and $S$; it gives $v$, $B=A$ and $\varphi$ with
$\varphi\sigma^{-1}\in H$. A place of inertia degree one has prime absolute norm
(\lref{GlobalTheory.absNorm_prime_iff_inertiaDeg_eq_one_and_finite_ramified}), and
$\sigma^{-1}\varphi=\sigma^{-1}(\varphi\sigma^{-1})\sigma\in H$ by normality.
\end{proof}

\begin{remark}
The hypothesis that $H$ contains the inertia elements above the places outside $S$ is not used in the proof.
\end{remark}

\subsubsection{Kronecker's congruence}
Let $E_4,E_6$ be the normalised Eisenstein series of weights $4$ and $6$ for $\SL_2(\Z)$, $A=E_4^3$,
$\Delta=(E_4^3-E_6^2)/1728$ and $j=A/\Delta$ on the upper half plane $\mathfrak H$. For $m\ge1$ let
$T_m=\{(a,b,e)\in\N^3: ae=m,\ 0\le b<e\}$, for $t=(a,b,e)\in T_m$ write $t\tau=(a\tau+b)/e$, and put
\[
G_m(\tau)=\prod_{t\in T_m}\bigl(A(\tau)\Delta(t\tau)-A(t\tau)\Delta(\tau)\bigr)
=\Delta(\tau)^{|T_m|}\prod_{t\in T_m}\Delta(t\tau)\,\bigl(j(\tau)-j(t\tau)\bigr).
\]
Let $a(q),d(q)\in\Z[[q]]$ be the $q$-expansions of $A$ and $\Delta$, let $\zeta_e=e^{2\pi i/e}$ and
$R_m=\Z[\zeta_m]$. We use the following shared facts on modular forms, all classical
\cite[\S11]{Cox1989}:
\begin{enumerate}[label=(M\arabic*)]
\item\label{pf:reckernel:M1} $G_m$ is a modular form of weight $24|T_m|$ for $\SL_2(\Z)$
  (\lref{ReciprocityTheory.ChartRestriction.L1.gmForm}).
\item\label{pf:reckernel:M2} If $g_m(q)$ is the $q$-expansion of $G_m$, then in $\C[[Q]]$
  \[
  g_m(Q^m)=\prod_{(a,b,e)\in T_m}\Bigl(a(Q^m)\,d(\zeta_e^bQ^{a^2})-a(\zeta_e^bQ^{a^2})\,d(Q^m)\Bigr)
  \]
  (\lref{ReciprocityTheory.ChartRestriction.L1.gmQ_xp}; note $e^{2\pi i t\tau}=\zeta_e^bQ^{a^2}$ for
  $Q=e^{2\pi i\tau/m}$).
\item\label{pf:reckernel:M3} A modular form of weight $12N$ for $\SL_2(\Z)$ whose $q$-expansion has coefficients in
  a subring $R\subseteq\C$ equals $\Delta^NP(j)$ for a polynomial $P\in R[X]$ of degree $\le N$
  (\lref{ReciprocityTheory.ChartRestriction.L1.qPrin}).
\item\label{pf:reckernel:M4} If $m$ is not a square and $G_m(\tau)=0$, then $j(\tau)$ is an algebraic integer
  (\lref{ReciprocityTheory.ChartRestriction.L1.jT_isIntegral}).
\item\label{pf:reckernel:M5} If $A',B',C',D'\in\Z$ with $A'D'-B'C'=m>0$ and $A'\tau+B'=\tau(C'\tau+D')$, then
  $G_m(\tau)=0$ (\lref{ReciprocityTheory.ChartRestriction.L1.Gm_vanish}).
\item\label{pf:reckernel:M6} Every lattice is $L=c(\Z\tau+\Z)$ with $c\ne0$, $\tau\in\mathfrak H$, and then
  $j(L)=j(\tau)$; if $L$ has complex multiplication by $\omega_d$, then for all $x,y\in\Z$ there are
  $A',B',C',D'\in\Z$ with $(x+y\omega_d)\tau=A'\tau+B'$, $x+y\omega_d=C'\tau+D'$ and
  $4(A'D'-B'C')=4x^2+4xy+(1+d)y^2$ (\lref{ReciprocityTheory.ChartRestriction.L1.latTau}, \lref{ReciprocityTheory.ChartRestriction.L1.jInv_eq_jT}, and \lref{ReciprocityTheory.ChartRestriction.L1.cmPoint}).
\end{enumerate}

For power series $f,g\in\C[[Q]]$ and a prime $p$ write $f\equiv g$ if $g\in R_p[[Q]]$ and
$f-g\in(1-\zeta_p)R_p[[Q]]$. This relation is compatible with sums and products of series with coefficients in
$R_p$, and with the substitution $Q\mapsto Q^p$.

\begin{lemma}[$q$-expansion congruence]\label{pf:reckernel:kronSer}
Let $p$ be a prime and $P_p(Q)$ the right-hand side of \ref{pf:reckernel:M2} for $m=p$. Then
$P_p(Q)+\bigl((ad^p-a^pd)(Q^p)\bigr)^2\in(1-\zeta_p)R_p[[Q]]$.
\formallabel{ReciprocityTheory.ChartRestriction.L1.kronSer}
\end{lemma}

\begin{proof}
We have $T_p=\{(p,0,1)\}\cup\{(1,b,p):0\le b<p\}$. We use three elementary congruences for $f\in\Z[[Q]]$:
(i) $f(\zeta_p^bQ)\equiv f(Q)$, since $\zeta_p\equiv1$ modulo $1-\zeta_p$;
(ii) $f(Q^p)-f(Q)^p\in p\,\Z[[Q]]$, since $f(Q)^p\equiv f(Q^p)$ modulo $p$ (Frobenius on $\F_p[[Q]]$), and
$p\in(1-\zeta_p)R_p$; (iii) $(F-G)^p-F^p+G^p\in p\,\Z[[Q]]$ for $F,G\in\Z[[Q]]$ (the binomial coefficients
$\binom pk$, $0<k<p$, are divisible by $p$; for $p=2$ use $2G^2\in2\Z[[Q]]$).
The factor of $t=(p,0,1)$ is $a(Q^p)d(Q^{p^2})-a(Q^{p^2})d(Q^p)$; by (ii) applied at the variable $Q^p$ it is
$\equiv\bigl(ad^p-a^pd\bigr)(Q^p)$. The factor of $t=(1,b,p)$ is $a(Q^p)d(\zeta_p^bQ)-a(\zeta_p^bQ)d(Q^p)$, which is
$\equiv a(Q^p)d(Q)-a(Q)d(Q^p)$ by (i). The product of the $p$ latter factors is therefore
\[
\equiv\bigl(a(Q^p)d(Q)-a(Q)d(Q^p)\bigr)^p\equiv a(Q^p)^pd(Q)^p-a(Q)^pd(Q^p)^p
\equiv a(Q^p)^pd(Q^p)-a(Q^p)d(Q^p)^p,
\]
by (iii) and then (ii) for $a$ and $d$; the right-hand side is $-(ad^p-a^pd)(Q^p)$. Multiplying the two
congruences, $P_p\equiv-\bigl((ad^p-a^pd)(Q^p)\bigr)^2$, which is the claim.
\end{proof}

\begin{lemma}[Kronecker's congruence at a point]\label{pf:reckernel:kronPt}
Let $p$ be a prime and $\tau\in\mathfrak H$ with $G_p(\tau)=0$. Then there is an algebraic integer $r$ with
\[
\bigl(j(\tau)^p-j(\tau)\bigr)^2=(1-\zeta_p)\,r .
\]
\formallabel{ReciprocityTheory.ChartRestriction.L1.kronPt}
\formallabel{ReciprocityTheory.ChartRestriction.L1.kronFun_poly}
\formallabel{ReciprocityTheory.ChartRestriction.L1.kronFun_div}
\end{lemma}

\begin{proof}
Put $\mathcal K_p=G_p+(A\Delta^p-A^p\Delta)^2$. Since $|T_p|=p+1$, both summands are modular forms of weight
$24(p+1)$ \ref{pf:reckernel:M1}. Let $k(q)$ be the $q$-expansion of $\mathcal K_p$. By \ref{pf:reckernel:M2} and
\cref{pf:reckernel:kronSer}, $k(Q^p)=(1-\zeta_p)w(Q)$ with $w\in R_p[[Q]]$; comparing coefficients,
$k(q)=(1-\zeta_p)\sum_nw_{pn}q^n$. Hence $(1-\zeta_p)^{-1}\mathcal K_p$ is a modular form of weight $12\cdot2(p+1)$
with $q$-expansion in $R_p[[q]]$, and by \ref{pf:reckernel:M3}
$\mathcal K_p=(1-\zeta_p)\Delta^{2(p+1)}P(j)$ with $P\in R_p[X]$. At $\tau$ we have $G_p(\tau)=0$ and $A=j\Delta$,
so $A\Delta^p-A^p\Delta=\Delta^{p+1}(j-j^p)$ and
$\Delta(\tau)^{2(p+1)}(j(\tau)^p-j(\tau))^2=(1-\zeta_p)\Delta(\tau)^{2(p+1)}P(j(\tau))$. As $\Delta(\tau)\ne0$,
the claim holds with $r=P(j(\tau))$, which is an algebraic integer because $j(\tau)$ is one
\ref{pf:reckernel:M4} ($p$ is not a square) and $R_p$ consists of algebraic integers.
\end{proof}

\begin{lemma}\label{pf:reckernel:zetaPow}
For a prime $p$ there is $r_0\in R_p$ with $(1-\zeta_p)^p=p\,r_0$.
\formallabel{ReciprocityTheory.ChartRestriction.L1.one_sub_zeta_pow}
\end{lemma}

\begin{proof}
By the binomial theorem $(1-\zeta_p)^p=1+(-\zeta_p)^p+p\,s$ with $s\in R_p$. For odd $p$, $(-\zeta_p)^p=-1$; for
$p=2$, $(-\zeta_2)^2=1$ and $1+1=2$.
\end{proof}

\begin{lemma}[Kronecker's congruence for CM values]\label{pf:reckernel:frobCong}
Let $L$ have complex multiplication by $\omega_d$, let $x=j(L)\in\Omega_d$, let $v$ be a finite place of $K_d$ of
prime absolute norm $N=\Nm v$ whose ideal class is trivial, and let $A$ be a valuation ring of $\Omega_d$ above
$v$. Then $x^N-x\in\reckm_A$.
\formallabel{ReciprocityTheory.ChartRestriction.L1.frobCong}
\formallabel{ReciprocityTheory.ChartRestriction.L1.princ_form}
\end{lemma}

\begin{proof}
Since $v$ is principal and $\OO_{K_d}=\Z[\omega_d]$, $v=(a+b\omega_d)$ with $a,b\in\Z$, and
$N=|\Nm(a+b\omega_d)|=a^2+ab+\tfrac{1+d}{4}b^2$, i.e.\ $4N=4a^2+4ab+(1+d)b^2$. Write $L=c(\Z\tau+\Z)$ as in
\ref{pf:reckernel:M6}; the integers $A',B',C',D'$ attached to $(a,b)$ satisfy $A'D'-B'C'=N$ and
$A'\tau+B'=(a+b\omega_d)\tau=\tau(C'\tau+D')$, so $G_N(\tau)=0$ by \ref{pf:reckernel:M5}. By
\cref{pf:reckernel:kronPt,pf:reckernel:zetaPow},
\[
(x^N-x)^{2N}=(1-\zeta_N)^Nr^N=N\,r_0\,r^N
\]
with algebraic integers $r,r_0$. Every valuation ring of $\Omega_d$ contains the algebraic integers, and
$N\in\reckm_A$ because $N\in v=\OO_{K_d}\cap\reckm_A$. Hence $(x^N-x)^{2N}\in\reckm_A$, and so $x^N-x\in\reckm_A$.
\end{proof}

\subsubsection{Inertia fixes the singular moduli}

\begin{lemma}\label{pf:reckernel:sepInt}
Let $\Omega$ be a field of characteristic $0$ and $s\subset\Omega$ a finite set of algebraic integers. There is
$n\in\Z\smallsetminus\{0\}$ such that for all $r\ne r'$ in $s$ and every valuation ring $A$ of $\Omega$ with
$r-r'\in\reckm_A$ one has $n\in\reckm_A$.
\formallabel{ReciprocityTheory.ChartRestriction.L1.sepInt}
\end{lemma}

\begin{proof}
For each pair $r\ne r'$ the non-zero algebraic integer $r-r'$ divides a non-zero rational integer $m_{r,r'}$ in the
ring of algebraic integers (e.g.\ its norm): $m_{r,r'}=(r-r')\gamma_{r,r'}$ with $\gamma_{r,r'}$ integral. Let $n$ be
the product of all $m_{r,r'}$. If $r-r'\in\reckm_A$, then $n=(r-r')\gamma_{r,r'}\cdot n/m_{r,r'}\in\reckm_A$, because
$\gamma_{r,r'}$ and the integer $n/m_{r,r'}$ lie in $A$.
\end{proof}

\begin{theorem}[Inertia fixes $j$]\label{pf:reckernel:inertiaFixes}
Assume the conclusions of \cref{pf:reckernel:artinSymbol} and of \cref{pf:reckernel:cheb} (for all number fields).
Let $L$ have complex multiplication by $\omega_d$ and $x=j(L)\in\Omega_d$. Then $\sigma x=x$ for every valuation
ring $B$ of $\Omega_d$ and every inertia element $\sigma\in\Gal(\Omega_d/K_d)$ at $B$.
\formallabel{GlobalTheory.ArtinClassCM.inertia_fixes}
\end{theorem}

\begin{proof}
Let $a=a_{K_d}$ be the Artin symbol of \cref{pf:reckernel:artinSymbol}. Since $x$ is an algebraic integer, it is a
root of a monic $P\in\Z[X]$; let $s\subset\Omega_d$ be the (finite, Galois-stable) set of roots of $P$, let $n$ be
as in \cref{pf:reckernel:sepInt} and $S$ the finite set of places $v$ of $K_d$ with $n\in v$. Then:
\begin{itemize}
\item[($\ast$)] if $A$ lies above $v\notin S$, $r,r'\in s$ and $r-r'\in\reckm_A$, then $r=r'$.
\end{itemize}
Let $H=\ker a\cap\{g: g|_s=\mathrm{id}\}$, an open normal subgroup. It contains every inertia element $g$ at a
valuation ring $A$ above a place $v\notin S$: $a(g)=1$, and for $r\in s$ we have $gr\in s$ and $gr-r\in\reckm_A$, so
$gr=r$ by ($\ast$). By \cref{pf:reckernel:cheb} (with this $S$, $H$ and $\sigma$) there are a place $v\notin S$ of prime norm $N$, a valuation ring $A$ above $v$ and an arithmetic
Frobenius $\varphi$ at $A$ with $\sigma^{-1}\varphi\in H$. Now $a(\sigma)=1$ because $\sigma$ is an inertia element,
so $a(\varphi)=a(\sigma)\,a(\sigma^{-1}\varphi)=1$, while $a(\varphi)=[v]$; thus $v$ has trivial class. By
\cref{pf:reckernel:frobCong}, $x^N-x\in\reckm_A$, and $\varphi x-x^N\in\reckm_A$ since $\varphi$ is a Frobenius
and $x\in A$. Hence $\varphi x-x\in\reckm_A$ with $\varphi x,x\in s$, and ($\ast$) gives $\varphi x=x$. Finally
$\sigma^{-1}\varphi$ fixes $s$, so $\sigma^{-1}x=\sigma^{-1}\varphi x=x$, i.e.\ $\sigma x=x$.
\end{proof}

\subsubsection{Singular moduli are unramified}

\begin{lemma}[Hensel's lemma in $\reckKur$]\label{pf:reckernel:hensel}
Let $\ell$ be a prime, $P\in\reckKur[X]$ with all coefficients of absolute value $\le1$, and $x\in\reckCl$ with
$|x|\le1$, $P(x)=0$ and $|P'(x)|=1$. Then $x\in\reckKur$. In particular, if $\ell\ne2$, $c\in\reckKur$, $y^2=c$ and
$|y|=1$, then $y\in\reckKur$; and if $\ell\nmid 2d$, then every $s\in\reckCl$ with $s^2=-d$, as well as
$\iota(\omega_d)$ for every field isomorphism $\iota\colon\C\to\reckCl$, lies in $\reckKur$.
\formallabel{ReciprocityTheory.ChartRestriction.L1.hensel_mem}
\formallabel{ReciprocityTheory.ChartRestriction.L1.sqrt_mem}
\formallabel{ReciprocityTheory.ChartRestriction.L1.sqrt_mem_Kur}
\formallabel{ReciprocityTheory.ChartRestriction.L1.omega_mem_kur}
\formallabel{ReciprocityTheory.ChartRestriction.L1.res_full}
\end{lemma}

\begin{proof}
Every element of $\reckCl$ of absolute value $\le1$ is congruent modulo the maximal ideal to $0$ or to a root of
unity of order prime to $\ell$ (the residue field of $\reckCl$ is $\overline{\F}_\ell$, whose non-zero elements
are roots of unity of order prime to $\ell$, and these lift to roots of unity in $\reckCl$). Choose such a $\xi\in\reckKur$ with $|x-\xi|<1$.
Since the coefficients of $P$ and $P'$ are integral, $|P(\xi)|=|P(\xi)-P(x)|<1$ and $|P'(\xi)|=1$. Newton's
iteration starting at $\xi$ stays in the complete field $\reckKur$ and converges to a root $y$ with $|y-\xi|<1$;
since $|P'|=1$ on the disc $|z-\xi|<1$, $P$ has at most one root there, so $x=y\in\reckKur$. For the square root
take $P=X^2-c$, $P'(y)=2y$, $|2|=1$. For $s$ note $|s|^2=|d|=1$. For $w=\iota(\omega_d)$ take
$P=X^2-X+\tfrac{d+1}{4}$: if $|w|>1$ then $|w(w-1)|=|w|^2>1\ge|\tfrac{d+1}4|$, impossible, so $|w|\le1$; and
$|P'(w)|=|2w-1|=|\iota(\sqrt{-d})|=1$.
\end{proof}

For a prime $\ell$ and a field isomorphism $\iota\colon\C\to\reckCl$ let $A_\iota=\{y\in\Omega_d:|\iota(y)|\le1\}$, a
valuation ring of $\Omega_d$; its maximal ideal consists of the $y$ with $|\iota(y)|<1$.

\begin{lemma}[Inertia-invariant elements]\label{pf:reckernel:inertKur}
Let $\ell\ne2$ be a prime with $\ell\nmid d$, $\iota\colon\C\to\reckCl$ a field isomorphism, and $x\in\Omega_d$ with
$\sigma x=x$ for every inertia element $\sigma\in\Gal(\Omega_d/K_d)$ at $A_\iota$. Then $\iota(x)\in\reckKur$.
\formallabel{ReciprocityTheory.ChartRestriction.L1.inertKur}
\formallabel{ReciprocityTheory.ChartRestriction.L1.locLift}
\formallabel{ReciprocityTheory.ChartRestriction.L1.kur_le}
\formallabel{ReciprocityTheory.ChartRestriction.L1.locEmb}
\end{lemma}

\begin{proof}
Let $\overline{\Q}_\ell\subset\reckCl$ be the algebraic closure of $\Q_\ell$. Since $\Omega_d$ is algebraic over
$\Q$, $\iota(\Omega_d)\subseteq\overline\Q_\ell$, so $\iota|_{\Omega_d}$ is a ring homomorphism
$e\colon\Omega_d\to\overline\Q_\ell$. Let
$s=e(\sqrt{-d})$ and $F=\Q_\ell(s)$, a finite extension of $\Q_\ell$.

\emph{Step 1: local inertia restricts to global inertia.} Let $\rho\in\Gal(\overline\Q_\ell/F)$ with
$|\rho z-z|<1$ for all $|z|\le1$. For $y\in\Omega_d$, $\rho(e(y))$ is a root of the minimal polynomial of $y$
over $\Q$, which splits in $\Omega_d$; hence $\rho(e(\Omega_d))=e(\Omega_d)$ and $\sigma=e^{-1}\circ\rho\circ e$ is an
automorphism of $\Omega_d$; it fixes $K_d=\Q(\sqrt{-d})$ because $\rho$ fixes $s$. Since $|\rho z|\le1$ iff
$|z|\le1$, $\sigma$ preserves $A_\iota$, and $|\iota(\sigma y)-\iota(y)|=|\rho(e(y))-e(y)|<1$ for $y\in A_\iota$. So
$\sigma$ is an inertia element at $A_\iota$, and $\rho(e(x))=e(\sigma x)=e(x)$.

\emph{Step 2.} By Step 1, $e(x)$ is fixed by the inertia subgroup of $\Gal(\overline\Q_\ell/F)$, hence lies in the
maximal unramified extension $F^{\mathrm{ur}}=\bigcup_n F(\mu_{q^n-1})$ of $F$, where $q$ is the cardinality of
the residue field of $F$ (local field theory; \lref{ReciprocityTheory.DiscreteGroups.LocalClassField.mem_subG_Kur_iff} and \lref{ReciprocityTheory.DiscreteGroups.LocalClassField.mem_iff_fixed}).

\emph{Step 3: $F^{\mathrm{ur}}\subseteq\reckKur$.} We have $s\in\reckKur$ by \cref{pf:reckernel:hensel}. The residue
characteristic of $F$ is $\ell$, so $\ell\mid q$ and $q^n-1$ is prime to $\ell$; hence the roots of unity of order
$q^n-1$ lie in $\Q(\mu_{\ell'})\subseteq\reckKur$. Since $\reckKur$ is a field containing $\Q_\ell$ (it is closed
and contains $\Q$), it contains every $F(\mu_{q^n-1})$. Therefore $\iota(x)=e(x)\in\reckKur$.
\end{proof}

\begin{corollary}[Singular moduli in $\reckKur$]\label{pf:reckernel:jUnram}
Assume global Artin reciprocity. For $d$ and $L$ with complex multiplication by $\omega_d$, every prime $\ell\ge5$
with $\ell\nmid d$ and every field isomorphism $\iota\colon\C\to\reckCl$, $\iota(j(L))\in\reckKur$ and
$|\iota(j(L))|\le1$.
\end{corollary}

\begin{proof}
$j(L)$ is an algebraic integer, so $j(L)\in\Omega_d$ and $|\iota(j(L))|\le1$. By
\cref{pf:reckernel:artinSymbol,pf:reckernel:cheb,pf:reckernel:inertiaFixes}, $j(L)$ is fixed by every inertia
element at $A_\iota$, and \cref{pf:reckernel:inertKur} applies.
\end{proof}

\subsubsection{Kernel polynomials}
Let $L$ be a lattice and $\alpha\in\C^\times$ with $\alpha L\subseteq L$ and $N_\alpha:=|\alpha|^2$ an odd integer. A
finite set $T\subset\C$ is a set of \emph{half representatives} for $\alpha$ if $\alpha P\in L$ and $2P\notin L$ for
$P\in T$; $P\pm P'\in L$ with $P,P'\in T$ implies $P=P'$; and every $z\notin L$ with $\alpha z\in L$ is congruent to
$\pm P$ modulo $L$ for some $P\in T$.

\begin{lemma}\label{pf:reckernel:halfReps}
For $\alpha$ as above, sets of half representatives exist; for each of them $2|T|+1=N_\alpha$, $\wp$ is injective on
$T$, and $N_\alpha z\in L$ whenever $\alpha z\in L$.
\formallabel{ReciprocityTheory.ChartRestriction.L1.halfReps_exist}
\formallabel{ReciprocityTheory.ChartRestriction.L1.IsHalfReps.card}
\formallabel{ReciprocityTheory.ChartRestriction.L1.ker_norm}
\formallabel{ReciprocityTheory.ChartRestriction.L1.IsHalfReps}
\formallabel{Formalism.L2.Th1.F.mulStableIntegral}
\end{lemma}

\begin{proof}
The kernel of multiplication by $\alpha$ on $\C/L$ has $N_\alpha$ elements
(\lref{Formalism.L2.Th1.F.kerRepsCard}); as $N_\alpha$ is odd, it contains no point of order $2$, so its
non-zero elements fall into $(N_\alpha-1)/2$ pairs $\{P,-P\}$, and $T$ is obtained by choosing one element of each
pair. Since $\wp(z)=\wp(z')$ iff $z\equiv\pm z'$, $\wp$ is injective on $T$
(\lref{Formalism.L2.Th1.F.pInjective}). Finally $\alpha$ is an algebraic integer of degree $\le2$ acting on
$L\cong\Z^2$, so $\alpha^2-t\alpha+N_\alpha=0$ with $t\in\Z$, and $N_\alpha z=t(\alpha z)-\alpha(\alpha z)\in L$ if $\alpha z\in
L$.
\end{proof}

Let $f(X)=4X^3-g_2X-g_3$ with $g_i=g_i(L)$. For a set $T$ of half representatives put $x_P=\wp(P)$,
\begin{align*}
R_T(X)&=\prod_{P\in T}(X-x_P),\\
V_T(X)&=X\,R_T(X)^2+\sum_{P\in T}\Bigl(f(x_P)+\tfrac12f'(x_P)(X-x_P)\Bigr)\prod_{P'\ne P}(X-x_{P'})^2 .
\end{align*}

\begin{lemma}[V\'elu's formula]\label{pf:reckernel:velu}
For $z$ with $\alpha z\notin L$,
\[
\alpha^2\wp(\alpha z)=\wp(z)+\sum_{P\in T}\Bigl(\frac{f(x_P)}{(\wp(z)-x_P)^2}+\frac{f'(x_P)}{2(\wp(z)-x_P)}\Bigr)
=\frac{V_T(\wp(z))}{R_T(\wp(z))^2}.
\]
\formallabel{ReciprocityTheory.ChartRestriction.L1.velu_fun}
\formallabel{ReciprocityTheory.ChartRestriction.L1.velu_pair}
\end{lemma}

\begin{proof}
By V\'elu's formula \cite{Velu1971} (\lref{Formalism.L2.Th1.F.isogenyP}),
$\alpha^2\wp(\alpha z)=\wp(z)+\sum_{P\in S}(\wp(z+P)-\wp(P))$ for a set $S$ of representatives of the non-zero
elements of $\ker\alpha$; take $S=T\cup(-T)$. For $P\in T$, $X=\wp(z)$ and $x=x_P$, the addition theorem gives
$\wp(z+P)+\wp(z-P)=\frac{(X+x)(2Xx-g_2/2)-g_3}{(X-x)^2}$, and a direct computation using $\wp'(P)^2=f(x)$ shows
$\wp(z+P)+\wp(z-P)-2x=\frac{f(x)}{(X-x)^2}+\frac{f'(x)}{2(X-x)}$. The second equality is the definition of
$V_T$.
\end{proof}

For a field $E$, a power series $U\in E[[z]]$ with $U(0)=1$ and a polynomial $P\in E[X]$ of degree $\le m$, put
$H_m(P)=\sum_{i\le m}P_i\,U^i\,z^{2(m-i)}\in E[[z]]$. Then $H_m(P)H_n(P')=H_{m+n}(PP')$, $H_m$ is additive, and
$H_m(P)=0$ implies $P=0$: the coefficient of $z^{2(m-\deg P)}$ in $H_m(P)$ is the leading coefficient of $P$,
since $U(0)=1$ and the terms with $i<\deg P$ only contribute to higher powers of $z$. We take $U(z)=z^2\wp(z)$, i.e.\
$U=1+\sum_{k\ge2}c_kz^{2k}$ with the $c_k$ of \cref{pf:reckernel:def-CM}; its coefficients lie in $\Q(g_2,g_3)$.
For $z$ near $0$, $H_m(P)$ is the Taylor series of $z^{2m}P(\wp(z))$.

\begin{lemma}\label{pf:reckernel:veluSeries}
With $n=|T|$, $U(\alpha z)\,H_{2n}(R_T^2)=H_{2n+1}(V_T)$ in $\C[[z]]$.
\formallabel{ReciprocityTheory.ChartRestriction.L1.velu_series}
\end{lemma}

\begin{proof}
Both sides are Taylor series at $0$ of functions holomorphic near $0$, namely of
$(\alpha z)^2\wp(\alpha z)\cdot z^{4n}R_T(\wp(z))^2$ and $z^{4n+2}V_T(\wp(z))$. By \cref{pf:reckernel:velu} these
functions agree on a punctured neighbourhood of $0$ (where $\alpha z\notin L$ and $z\notin L$), hence the series are
equal.
\end{proof}

\begin{lemma}[Descent of a denominator]\label{pf:reckernel:proj}
Let $E$ be a field, $k\subseteq E$ a subfield, $U,V\in k[[z]]$ with $U(0)=1$, and $A,B\in E[X]$ coprime with $B$
monic of degree $N$ and $\deg A\le N+1$, such that $V\,H_N(B)=H_{N+1}(A)$. Then $B\in k[X]$.
\formallabel{ReciprocityTheory.ChartRestriction.L1.proj_trick}
\formallabel{ReciprocityTheory.ChartRestriction.L1.exists_proj}
\formallabel{ReciprocityTheory.ChartRestriction.L1.homS_eq_zero}
\end{lemma}

\begin{proof}
Choose a $k$-linear map $\psi\colon E\to k$ with $\psi|_k=\mathrm{id}$ (a left inverse of the
inclusion of $k$-vector spaces $k\subseteq E$), and let $A',B'$ be obtained by applying $\psi$ to
the coefficients of $A$ (in degrees $\le N+1$) and of $B$ (in degrees $\le N$). Since $U,V$ have coefficients in $k$
and the identity $VH_N(B)=H_{N+1}(A)$ is $k$-linear in the coefficients of $A$ and $B$, also
$VH_N(B')=H_{N+1}(A')$. Hence
$H_{2N+1}(AB'-A'B)=H_{N+1}(A)H_N(B')-H_{N+1}(A')H_N(B)=VH_N(B)H_N(B')-VH_N(B')H_N(B)=0$, so $AB'=A'B$. As $A$ and $B$
are coprime, $B\mid B'$. Now $B'$ has degree $\le N$ and coefficient $\psi(1)=1$ in degree $N$, so $B'$ is monic of
degree $N$, and $B=B'\in k[X]$.
\end{proof}

\begin{lemma}\label{pf:reckernel:sqrtPoly}
Let $E$ be a field of characteristic $0$, $k\subseteq E$ a subfield and $h\in E[X]$ monic with $h^2\in k[X]$. Then
$h\in k[X]$.
\formallabel{ReciprocityTheory.ChartRestriction.L1.sqrt_poly}
\end{lemma}

\begin{proof}
Let $n=\deg h$. We show $h_i\in k$ by descending induction on $i\le n$; $h_n=1$. For $i<n$, the coefficient of
$X^{i+n}$ in $h^2$ is $2h_i+\sum h_ah_b$, the sum over $a+b=i+n$ with $a,b\notin\{i,n\}$; in that sum $i<a,b<n$ (or a
factor vanishes), so all terms lie in $k$ by induction, and $2h_i\in k$.
\end{proof}

\begin{lemma}[Kernel polynomials]\label{pf:reckernel:kerPoly}
Let $\alpha,T$ be as above and $k\subseteq\C$ a subfield containing $g_2(L),g_3(L)$ and $\alpha$. Then $R_T\in k[X]$
and $\prod_{P\in T}\wp(P)\in k$.
\formallabel{ReciprocityTheory.ChartRestriction.L1.kerPoly_mem}
\formallabel{ReciprocityTheory.ChartRestriction.L1.prod_mem_of_rootPoly}
\end{lemma}

\begin{proof}
Apply \cref{pf:reckernel:proj} with $U(z)=z^2\wp(z)$, $V(z)=U(\alpha z)$ (whose coefficients $\alpha^mU_m$ lie in $k$),
$A=V_T$ (of degree $\le2n+1$), $B=R_T^2$ (monic of degree $2n$) and the identity of \cref{pf:reckernel:veluSeries}.
The polynomials $V_T$ and $R_T^2$ are coprime: at a root $x_P$ of $R_T$,
$V_T(x_P)=f(x_P)\prod_{P'\ne P}(x_P-x_{P'})^2\ne0$, since the $x_{P'}$ are distinct and $f(x_P)=\wp'(P)^2\ne0$ ($P$ is
not a $2$-torsion point). Thus $R_T^2\in k[X]$, and $R_T\in k[X]$ by \cref{pf:reckernel:sqrtPoly}. Finally
$\prod_P\wp(P)=(-1)^{|T|}R_T(0)$.
\end{proof}

\begin{lemma}[Ideal kernels]\label{pf:reckernel:kerIdeal}
Let $L$ have complex multiplication by $\omega_d$ with $3\mid d$ and $d\ne3$, and let $\ell\ge5$ be a prime. There are
a prime $p>\ell$ with $p\equiv2\pmod3$ and a finite set $T''\subset\C$ with $2|T''|+1=p$, such that $z\notin L$ and $pz\in
L$ for $z\in T''$, and $\prod_{z\in T''}\wp(z)\in k$ for every subfield $k\subseteq\C$ containing $g_2(L)$, $g_3(L)$ and
$\omega_d$.
\formallabel{ReciprocityTheory.ChartRestriction.L1.kerProd_ideal}
\formallabel{ReciprocityTheory.ChartRestriction.L1.prime_choice}
\formallabel{ReciprocityTheory.ChartRestriction.L1.card_G}
\formallabel{ReciprocityTheory.ChartRestriction.L1.rootPoly_inter_mem}
\formallabel{ReciprocityTheory.ChartRestriction.L1.xsOf_filter}
\end{lemma}

\begin{proof}
\emph{Choice of $p$.} Write $d=3d'$; then $d'>1$ and $d'\equiv1\pmod4$. Choose a prime $q\mid d'$ and a non-square $c$
modulo $q$, and by the Chinese remainder theorem $a\in\N$ with $a\equiv2\pmod3$, $a\equiv c\pmod q$ and
$a\equiv1\pmod{d'/q}$; then the Jacobi symbol $(\frac ad)=(\frac a3)(\frac aq)(\frac a{d'/q})=(-1)(-1)(1)=1$. By
Dirichlet's theorem on primes in arithmetic progressions there is a prime $p\equiv a\pmod d$ with $p>\ell+3+d$; then
$p\equiv2\pmod3$, and quadratic reciprocity gives $(\frac{-d}p)=1$, so $p\mid s^2+d$ for some $s$. Replacing $s$ by $s+p$
if necessary we may take $s$ even, and replacing it by $s+2p$ if $p^2\mid s^2+d$ we obtain an even $r$ with $p\mid
r^2+d$ and $p^2\nmid r^2+d$.

\emph{The kernel.} Let $\beta=r+\sqrt{-d}\in k$, so that $|\beta|^2=r^2+d=pM$ with $p\nmid M$, and $r^2+d$ is odd. Let
$T$ be half representatives for $\alpha=p$ and $T'$ for $\beta$ (\cref{pf:reckernel:halfReps}), and
$T''=\{P\in T:\beta P\in L\}$. The group $\{z:pz\in L,\ \beta z\in L\}/L$ is the kernel of the reduction modulo $p$ of
the integer matrix of $\beta$ on $L\cong\Z^2$; this reduction is singular (its determinant $pM$ is divisible by $p$) and
non-zero (otherwise $p^2$ would divide $pM$), so the kernel has $p$ elements and $2|T''|+1=p$. Moreover
$\wp(T'')=\wp(T)\cap\wp(T')$, and $\prod_{x\in\wp(T)\cap\wp(T')}(X-x)$ is the monic greatest common divisor of $R_T$
and $R_{T'}$, which lie in $k[X]$ by \cref{pf:reckernel:kerPoly}; the greatest common divisor of two polynomials in
$k[X]$ lies in $k[X]$. Hence $\prod_{T''}\wp\in k$ ($\wp$ is injective on $T''$).
\end{proof}

\begin{lemma}[Kernel products]\label{pf:reckernel:kerProd}
Let $L$ have complex multiplication by $\omega_d$, let $\ell\ge5$ be a prime with $\ell\nmid d$, and let $e=2$, or $e=3$
and $d\ne3$. There are $N\ge1$ and a finite set $T\subset\C$ such that $\ell\nmid N$, $e\nmid N$, $e\nmid|T|$,
$z\notin L$ and $Nz\in L$ for all $z\in T$, and $\prod_{z\in T}\wp(z)\in k$ for every subfield $k\subseteq\C$
containing $g_2(L)$, $g_3(L)$ and $\omega_d$.
\formallabel{ReciprocityTheory.ChartRestriction.L1.kerProd}
\formallabel{ReciprocityTheory.ChartRestriction.L1.kerProdAt_of_norm}
\formallabel{ReciprocityTheory.ChartRestriction.L1.kerProd_prin}
\end{lemma}

\begin{proof}
\emph{Principal case.} Suppose $r\in\N$ is such that $N=r^2+d$ is odd, $\ell\nmid N$, $e\nmid N$ and $e\nmid\frac{N-1}2$.
Let $\alpha=r+\sqrt{-d}=r+2\omega_d-1\in\OO_{K_d}$; then $\alpha L\subseteq L$, $|\alpha|^2=N$ and $\alpha\in k$. Take for
$T$ a set of half representatives (\cref{pf:reckernel:halfReps}): $|T|=\frac{N-1}2$, its elements are not in $L$
and are killed by $N$, and $\prod_T\wp\in k$ by \cref{pf:reckernel:kerPoly}. Suitable $r$ exist in the following
cases. If $e=2$, take $r=0$: $N=d$ is odd and prime to $\ell$, and $\frac{d-1}2$ is odd as $d\equiv3\pmod4$. If $e=3$
and $d\equiv2\pmod3$, take $r=0$: then $d\equiv11\pmod{12}$ and $\frac{d-1}2\equiv5\pmod6$. If $e=3$ and
$d\equiv1\pmod 3$, then $d\equiv7\pmod{12}$; take $r\in\{2,4\}$ with $\ell\nmid r^2+d$ (possible because the two
values differ by $12$ and $\ell\ge5$); then $r^2+d\equiv11\pmod{12}$.

\emph{Remaining case: $e=3$ and $3\mid d$.} Then $d\ne3$, and \cref{pf:reckernel:kerIdeal} gives $p$ and $T''$; take
$N=p$ and $T=T''$: $\ell\nmid p$ as $p>\ell$, $3\nmid p$, and $3\nmid\frac{p-1}2$ as $p\equiv2\pmod3$.
\end{proof}

\subsubsection{Units among torsion abscissae}

\begin{lemma}[Scaled torsion points]\label{pf:reckernel:scaledTorsion}
Let $\ell$ be a prime, $\iota\colon\C\to\reckCl$ a field isomorphism, $L$ a lattice, $\lambda\in\reckCl^\times$,
$z\notin L$ and $N\ge1$ with $Nz\in L$. Then $\bigl(\lambda^2\iota(\wp(z)),\tfrac12\lambda^3\iota(\wp'(z))\bigr)$
is a non-singular point of $W_{\lambda^4\iota(g_2(L)),\lambda^6\iota(g_3(L))}$ whose $N$-th multiple is the origin.
\formallabel{Formalism.L2.Th1.E.scaled_torsion_point_of_add}
\formallabel{ReciprocityTheory.ChartRestriction.L1.torPt}
\end{lemma}

\begin{proof}
The map $z\mapsto(\wp(z),\frac12\wp'(z))$ (and $L\mapsto$ origin) is a homomorphism from $\C/L$ to the group of
points of $W_{g_2(L),g_3(L)}$ (\lref{Formalism.L2.Th1.E.PtOfAdd_of}); applying the field isomorphism $\iota$
and the isomorphism $(x,y)\mapsto(\lambda^2x,\lambda^3y)$ of Weierstrass curves preserves non-singularity and the
group law.
\end{proof}

We use two shared facts on Weierstrass curves $W$ with integral coefficients and unit discriminant over a field
$C$ with an ultrametric absolute value, and an integer $n$ with $|n|=1$: a point $P=(x,y)$ with $nP=0$ has
$|x|,|y|\le1$ (\lref{Formalism.ELL.torsionIntegral_of}); and if $P=0$ or $|x(P)|>1$ (we say $P$ \emph{lies in
the formal group}) and $nP=0$, then $P=0$
(\lref{Formalism.ELL.formal_group_points_no_prime_to_p_torsion}); see \cite[Ch.~VII]{Silverman2009}.

\begin{lemma}\label{pf:reckernel:torsUnit}
Let $C$ be a field with an ultrametric absolute value, $|2|=|3|=1$, $g_2,g_3\in C$ with $|g_2|,|g_3|\le1$ and
$|g_2^3-27g_3^2|=1$, $n\ge1$ with $|n|=1$, and $P=(x,y)$ a point of $W_{g_2,g_3}$ with $nP=0$.
\begin{enumerate}
\item If $|g_2|<1$ and $3\nmid n$, then $|x|=1$.
\item If $|g_3|<1$ and $2\nmid n$, then $|x|=1$.
\end{enumerate}
\formallabel{ReciprocityTheory.ChartRestriction.L1.torsUnit3}
\formallabel{ReciprocityTheory.ChartRestriction.L1.torsUnit2}
\end{lemma}

\begin{proof}
The model $y^2=x^3-\frac{g_2}4x-\frac{g_3}4$ is integral with unit discriminant, so $|x|,|y|\le1$; suppose $|x|<1$.

(1) As $|g_2|<1$ and $|g_2^3-27g_3^2|=1$, $|g_3|=1$; then $y^2=-\frac{g_3}4+x^3-\frac{g_2}4x$ gives $|y|=1$. The
tangent slope $\lambda=(3x^2-\frac{g_2}4)/(2y)$ has $|\lambda|<1$, so $2P=(x_2,y_2)$ with $x_2=\lambda^2-2x$,
$y_2=-(\lambda(x_2-x)+y)$ satisfies $|x_2|<1$ and $|y_2+y|<1$. If $2P=-P$, then $3P=0$. Otherwise $x_2\ne x$ (for
$x_2=x$ we would have $y_2=\pm y$, and $y_2=y$ contradicts $|y_2+y|<1=|2y|$), the chord slope
$m=(y_2-y)/(x_2-x)$ has $|y_2-y|=1>|x_2-x|$, hence $|m|>1$ and $|x(3P)|=|m^2-x_2-x|>1$. In both cases $3P$ lies in
the formal group and $n(3P)=3(nP)=0$, so $3P=0$. Since $\gcd(3,n)=1$ and $nP=0$, $P=0$, a contradiction.

(2) As $|g_3|<1$, $|g_2|=1$, and $|y|<1$ from the equation. If $2P\ne0$, the tangent slope
$\lambda=(3x^2-\frac{g_2}4)/(2y)$ has numerator of absolute value $1$ and $|2y|<1$, so $|x(2P)|=|\lambda^2-2x|>1$.
Thus $2P$ lies in the formal group, and $n(2P)=0$ gives $2P=0$. With $n=2k+1$, $P=nP-k(2P)=0$, a contradiction.
\end{proof}

\begin{lemma}\label{pf:reckernel:unitProd}
Let $\ell\ge5$, $\iota$, $L$ as before and $\mu\in\reckCl$ with $\mu^{12}\iota(\Delta(L))=1$,
$|\mu^4\iota(g_2(L))|\le1$, $|\mu^6\iota(g_3(L))|\le1$. Let $e=3$ and $|\mu^4\iota(g_2(L))|<1$, or $e=2$ and
$|\mu^6\iota(g_3(L))|<1$. Let $N\ge1$ with $\ell\nmid N$, $e\nmid N$, and $T$ a finite set of $z\notin L$ with $Nz\in
L$. Then $|\mu|^{2|T|}\,\bigl|\iota\bigl(\prod_{z\in T}\wp(z)\bigr)\bigr|=1$.
\formallabel{ReciprocityTheory.ChartRestriction.L1.unit_prod}
\formallabel{ReciprocityTheory.ChartRestriction.L1.unit_point}
\end{lemma}

\begin{proof}
The curve $W_{\mu^4\iota(g_2),\mu^6\iota(g_3)}$ has discriminant $\mu^{12}\iota(\Delta(L))=1$, and $|2|=|3|=|N|=1$.
For $z\in T$ the point of \cref{pf:reckernel:scaledTorsion} (with $\lambda=\mu$) is killed by $N$, so
$|\mu^2\iota(\wp(z))|=1$ by \cref{pf:reckernel:torsUnit}. Take the product over $T$.
\end{proof}

\subsubsection{Divisibility of valuations}

\begin{lemma}\label{pf:reckernel:equalInv}
Let $L$ have complex multiplication by $\omega_d$.
\begin{enumerate}
\item If $d=3$, then $g_2(L)=0$.
\item If $g_2(L)g_3(L)\ne0$ and $c^2=g_3(L)/g_2(L)$, then $L_0=cL$ has complex multiplication by $\omega_d$,
  $j(L_0)=j(L)$ and $g_2(L_0)=g_3(L_0)=\frac{27\,j(L)}{j(L)-1728}$.
\end{enumerate}
\formallabel{ReciprocityTheory.ChartRestriction.L1.g2_eq_zero_of_hasCMBy_three}
\formallabel{ReciprocityTheory.ChartRestriction.L1.exists_cm_lattice_equal_invariants}
\end{lemma}

\begin{proof}
(1) $\omega_3$ is a primitive sixth root of unity and $\omega_3L=L$, so $\sum'_{w\in L}w^{-4}=\omega_3^{-4}\sum'_{w\in
L}w^{-4}$ with $\omega_3^4\ne1$. (2) $g_2(cL)=c^{-4}g_2=g_2^3/g_3^2$ and $g_3(cL)=c^{-6}g_3=g_2^3/g_3^2$, while
$j-1728=1728\cdot27g_3^2/\Delta$ gives $\frac{27j}{j-1728}=\frac{g_2^3}{g_3^2}$.
\end{proof}

For $\ell\ge5$ the integers $2,3,27,1728$ are units in $\reckCl$, and if $|J|<1$ then $|J-1728|=1$. Every
non-zero $x\in\reckKur$ has $|x|\in\ell^{\Z}$ (\lref{ReciprocityTheory.ChartRestriction.L1.kur_val}).

\begin{lemma}[Valuation of $j$]\label{pf:reckernel:valDiv3}
Let $L$ have complex multiplication by $\omega_d$, $\ell\ge5$ prime, $\ell\nmid d$, and $J=\iota(j(L))$ with
$J\in\reckKur$, $|J|\le1$ and $J\ne0$. Then $|J|=\ell^{3m}$ for some $m\in\Z$.
\formallabel{ReciprocityTheory.ChartRestriction.L1.valDiv3}
\formallabel{ReciprocityTheory.ChartRestriction.L1.three_dvd_of_norm_relations}
\end{lemma}

\begin{proof}
If $|J|=1$ take $m=0$; let $|J|<1$. Then $g_2(L)\ne0$ (as $j\ne0$), hence $d\ne3$ by
\cref{pf:reckernel:equalInv}(1), and $g_3(L)\ne0$ (as $J\ne1728$). Let $L_0$ be as in \cref{pf:reckernel:equalInv}(2)
and $c=27J/(J-1728)=\iota(g_2(L_0))=\iota(g_3(L_0))$; then $c\in\reckKur$, $|c|=|J|$ and
$|\iota(\Delta(L_0))|=|c^3-27c^2|=|c^2\cdot27\cdot1728/(J-1728)|=|J|^2$. Choose $\mu\in\reckCl$ with
$\mu^{12}\iota(\Delta(L_0))=1$; then $|\mu|^{12}|J|^2=1$, $|\mu^4c|^3=|J|<1$ and $|\mu^6c|^2=1$. Let $(N,T)$ be given
by \cref{pf:reckernel:kerProd} for $L_0$ and $e=3$, and apply it to the subfield
$k=\iota^{-1}(\reckKur)$, which contains $g_2(L_0)$, $g_3(L_0)$ and $\omega_d$ (\cref{pf:reckernel:hensel}):
$\Pi=\iota(\prod_T\wp)\in\reckKur$. By \cref{pf:reckernel:unitProd}, $|\mu|^{2|T|}|\Pi|=1$, so $\Pi\ne0$ and
$|\Pi|=\ell^a$, $|J|=\ell^n$ with $a,n\in\Z$. Raising $|\mu|^{12}\ell^{2n}=1$ to the power $|T|$ and
$|\mu|^{2|T|}\ell^a=1$ to the power $6$ gives $\ell^{2n|T|}=\ell^{6a}$, so $3\mid n|T|$, and $3\mid n$ since
$3\nmid|T|$.
\end{proof}

\begin{lemma}[Valuation of $j-1728$]\label{pf:reckernel:valDiv2}
In the situation of \cref{pf:reckernel:valDiv3} with $J\ne1728$ (instead of $J\ne0$), $|J-1728|=\ell^{2m}$ for some
$m\in\Z$.
\formallabel{ReciprocityTheory.ChartRestriction.L1.valDiv2}
\formallabel{ReciprocityTheory.ChartRestriction.L1.valDiv2_even_of_pows}
\end{lemma}

\begin{proof}
$|J-1728|\le1$; if it equals $1$ take $m=0$, so let $x:=|J-1728|<1$, whence $|J|=1$. Then $j(L)\ne0,1728$, so
$g_2(L)g_3(L)\ne0$; let $L_0$ and $c=27J/(J-1728)$ be as before, so $|c|=x^{-1}$ and
$\iota(\Delta(L_0))=c^2\cdot27\cdot1728/(J-1728)$ has absolute value $x^{-3}$. Choose $\mu$ with
$\mu^{12}\iota(\Delta(L_0))=1$ and put $t=|\mu|$: $t^{12}=x^3$, so $t^4=x$, $t<1$, $|\mu^4c|=1$ and
$|\mu^6c|=t^6x^{-1}=t^2<1$. With $(N,T)$ from \cref{pf:reckernel:kerProd} for $e=2$ ($N$ and $|T|$ odd), as in the
previous proof $|\Pi|=\ell^a$ and $t^{2|T|}\ell^a=1$, while $x=\ell^b$ for some $b\in\Z$ by
\lref{ReciprocityTheory.ChartRestriction.L1.kur_val} (non-zero elements of $\reckKur$ have absolute values in
$\ell^{\Z}$). From $t^4=\ell^b$ we get $\ell^{b|T|}=t^{4|T|}=\ell^{-2a}$,
so $b|T|=-2a$ and $b$ is even because $|T|$ is odd.
\end{proof}

\subsubsection{The good model}

\begin{lemma}\label{pf:reckernel:model}
Let $\ell\ge5$, $a,b\in\reckCl$ with $a^3-27b^2\ne0$, $J=1728a^3/(a^3-27b^2)\in\reckKur$, $|J|\le1$, and assume
$|J|\in\ell^{3\Z}$ if $J\ne0$ and $|J-1728|\in\ell^{2\Z}$ if $J\ne1728$. Then there are $g_2',g_3'\in\reckKur$ and
$\lambda\in\reckCl^\times$ with $|g_2'|,|g_3'|\le1$, $|g_2'^3-27g_3'^2|=1$, $g_2'=\lambda^4a$ and $g_3'=\lambda^6b$.
\formallabel{ReciprocityTheory.ChartRestriction.L1.model_of_valDiv}
\formallabel{ReciprocityTheory.ChartRestriction.L1.model_generic}
\end{lemma}

\begin{proof}
If $a=0$, then $b\ne0$; take $(g_2',g_3')=(0,1)$ and $\lambda^6=b^{-1}$ ($|0-27|=1$). If $b=0$, take $(1,0)$ and
$\lambda^4=a^{-1}$. Otherwise $J\ne0,1728$; write $|J|=X^3$, $|J-1728|=Y^2$ with $X=\ell^{m}$, $Y=\ell^{n}$, and
put $c=27J/(J-1728)=a^3/b^2\in\reckKur$ and $t=\ell^{m-n}\in\Q$, so that $|t|=Y/X$, $|c|=X^3/Y^2$ and
$|c-27|=|27\cdot1728/(J-1728)|=Y^{-2}$. Let $g_2'=t^2c$, $g_3'=t^3c$, both in $\reckKur$. Then $|g_2'|^3=|J|\le1$,
$|g_3'|^2=|J-1728|\le1$, and $g_2'^3-27g_3'^2=t^6c^2(c-27)$ has absolute value
$(Y/X)^6\cdot X^6/Y^4\cdot Y^{-2}=1$. Finally choose $\lambda$ with $\lambda^2=ta/b$; then
$\lambda^4a=t^2a^3/b^2=g_2'$ and $\lambda^6b=t^3a^3/b^2=g_3'$.
\end{proof}

\begin{theorem}\label{pf:reckernel:goodRed}
Assume that $|\iota(j(L))|\le1$ for all $d$, $L$ with complex multiplication by $\omega_d$, all primes $\ell$ and all
$\iota$, and that $\iota(j(L))\in\reckKur$ for all $d,L,\ell\ge5,\iota$ as in \textup{(GM)}. Then \textup{(GM)} holds. Consequently, the existence of ideal-class Artin symbols as in
\cref{pf:reckernel:artinSymbol} for all number fields, together with the conclusion of \cref{pf:reckernel:cheb},
implies \textup{(GM)}.
\formallabel{ReciprocityTheory.ChartRestriction.L1.goodReduction_of_integral_unramified}
\formallabel{ReciprocityTheory.ChartRestriction.L1.JUnram}
\formallabel{GlobalTheory.ArtinClassCM.good_model}
\formallabel{GlobalTheory.T.cmLattice_good_model_unramified}
\end{theorem}

\begin{proof}
Let $d,L,\ell,\iota$ be as in \textup{(GM)} and $J=\iota(j(L))$. By hypothesis $J\in\reckKur$ and $|J|\le1$, and by
\cref{pf:reckernel:valDiv3,pf:reckernel:valDiv2} the valuation conditions of \cref{pf:reckernel:model} hold. Apply
\cref{pf:reckernel:model} to $a=\iota(g_2(L))$, $b=\iota(g_3(L))$, for which $a^3-27b^2=\iota(\Delta(L))\ne0$ and
$J=1728a^3/(a^3-27b^2)$; it yields a good model of $L$ at $(\ell,\iota)$. For the second statement, the
hypothesis of the first holds by the integrality of $j(L)$ and by
\cref{pf:reckernel:inertiaFixes,pf:reckernel:inertKur}, as in the proof of \cref{pf:reckernel:jUnram}.
\end{proof}

\begin{proof}[Proof of \cref{pf:reckernel:cm2}]
Global Artin reciprocity gives the ideal-class Artin symbols (\cref{pf:reckernel:artinSymbol}), and
\cref{pf:reckernel:cheb} holds unconditionally; apply \cref{pf:reckernel:goodRed}.
\end{proof}

% ---------------------------------------------------------------------------
\subsection{Torsion points and formal groups}\label{pf:reckernel:sec-cm4}

\subsubsection{The formal group law of a Weierstrass curve}\label{pf:reckernel:sec-fgl}
Let $K$ be a field and $W\colon y^2+a_1xy+a_3y=x^3+a_2x^2+a_4x+a_6$ a Weierstrass curve over $K$. There is a unique
power series $w\in K[[T]]$ with $w(0)=0$ and
\begin{equation}\label{pf:reckernel:eq-w}
w=T^3+a_1Tw+a_2T^2w+a_3w^2+a_4Tw^2+a_6w^3
\end{equation}
(\lref{Formalism.ELL.wSeriesUnique}: existence and uniqueness of this series); comparing coefficients in
\eqref{pf:reckernel:eq-w} shows that its coefficients lie in $\Z[a_1,\dots,a_6]$ and that $w=T^3\upsilon$ with
$\upsilon(0)=1$. For a power series $s$ without constant term in finitely many variables write $w(s)$ for the
substituted series and put, in the fraction field of the power series ring,
\[
x(s)=s/w(s),\qquad y(s)=-1/w(s)
\]
(with the convention $1/0=0$). If $s\ne0$, then $w(s)\ne0$, and multiplying the Weierstrass equation by $w(s)^3$ shows
that it is equivalent to \eqref{pf:reckernel:eq-w} at $s$; so $P(s)=(x(s),y(s))$ is a point of $W$ over the fraction
field, and it is non-singular (\lref{Formalism.ELL.wSeriesNonsingularPoint_holds}: the point is non-singular).
Since $s=-x(s)/y(s)$ for $s\ne0$ and $y(0)=0$, a series without constant term is determined by $x(s)$ and $y(s)$.
We say that a series $r$ satisfies the \emph{chord law} for $(s,t)$ if $x(r)$ and $y(r)$ are the coordinates of the sum
of $P(s)$ and $P(t)$ given by the chord formulas, with slope $(y(s)-y(t))/(x(s)-x(t))$; if $s,t\ne0$ and
$x(s)\ne x(t)$, this means $P(r)=P(s)+P(t)$. We use the following facts, valid over every field.
\begin{enumerate}[label=(W\arabic*)]
\item\label{pf:reckernel:W1} For every $s$ without constant term, $w(s)$ is the only series $v$ without constant term
  with $v=s^3+a_1sv+a_2s^2v+a_3v^2+a_4sv^2+a_6v^3$ (\lref{Formalism.ELL.wSubstUnique}: uniqueness of solutions of
  this equation).
\item\label{pf:reckernel:W2} There is a formal group law $\mathcal F(X,Y)$ over $K$ with coefficients in
  $\Z[a_1,\dots,a_6]$ which satisfies the chord law for $(X,Y)$ (\lref{Formalism.ELL.wFormalGroup}: the formal
  group law of $W$).
\item\label{pf:reckernel:W3} If $r(X,Y)$ has no constant term and satisfies the chord law for $(X,Y)$, and $s,t$ are
  non-zero series without constant terms in finitely many variables with $x(s)\ne x(t)$, then $r(s,t)$ satisfies the chord
  law for $(s,t)$ (\lref{Formalism.ELL.wChordSubst}: the chord law is stable under substitution).
\item\label{pf:reckernel:W4} There is a series $\mathcal F_2(X,Y)$ with coefficients in $\Z[a_1,\dots,a_6]$,
  $\mathcal F_2\equiv X+Y$ modulo terms of degree $\ge2$, which satisfies the chord law for $(X,Y)$
  (\lref{Formalism.ELL.wChordSeries}: existence of the chord series).
\end{enumerate}
Note that $x(X)\ne x(Y)$ and $x(X)\ne x(\mathcal F)$ in the fraction field of $K[[X,Y]]$: substituting $X=0$ maps
$x(X)$ to $0$, and $x(Y)$, $x(\mathcal F(0,Y))=x(Y)$ to the non-zero series $1/(Y^2\upsilon(Y))$.

\begin{lemma}[Formal doubling and inversion]\label{pf:reckernel:fglDblInv}
Let $\mathcal F$ be as in \ref{pf:reckernel:W2}.
\begin{enumerate}
\item If $P(X)\ne-P(X)$, then $\mathcal D=\mathcal F(X,X)\ne0$ and $P(\mathcal D)=2P(X)$.
\item If $\psi\in K[[X,Y]]$ is the inverse of $E=1-a_1X-a_3w(X)$, then $\mathcal F(X,-X\psi)=0$.
\end{enumerate}
\formallabel{Formalism.ELL.FormalGroupPointsHom.dbl_formal}
\formallabel{Formalism.ELL.FormalGroupPointsHom.inv_formal}
\end{lemma}

\begin{proof}
(1) By the chord law, $P(\mathcal F)=P(X)+P(Y)$, and by \ref{pf:reckernel:W3} for $(s,t)=(X,\mathcal F)$,
$P(\mathcal F(X,\mathcal F))=P(X)+P(\mathcal F)=2P(X)+P(Y)$. By associativity
$\mathcal F(\mathcal D,Y)=\mathcal F(X,\mathcal F(X,Y))$. If $\mathcal D=0$, then $\mathcal F(X,\mathcal F)=\mathcal F(0,Y)=Y$,
so $2P(X)+P(Y)=P(Y)$ and $P(X)=-P(X)$, which is excluded. Hence $\mathcal D\ne0$. Substituting $X=0$ maps $\mathcal D$ to
$\mathcal F(0,0)=0$ and fixes $Y$, so $x(\mathcal D)\ne x(Y)$, and \ref{pf:reckernel:W3} for $(s,t)=(\mathcal D,Y)$ gives
$P(\mathcal D)+P(Y)=P(\mathcal F(\mathcal D,Y))=2P(X)+P(Y)$, i.e.\ $P(\mathcal D)=2P(X)$.

(2) Put $J=-X\psi$ and $v=-w(X)\psi$. Multiplying the relation of \ref{pf:reckernel:W1} for $(J,v)$ by $E^3$ and using
$E\psi=1$ and the relation \eqref{pf:reckernel:eq-w} at $X$ one checks that $(J,v)$ satisfies the relation of
\ref{pf:reckernel:W1} ($E\ne0$, and the difference of the two sides is a combination of the two relations); hence
$w(J)=-w(X)\psi$. Then $x(J)=x(X)$ and $y(J)=-1/w(J)=-y(X)-a_1x(X)-a_3$, i.e.\ $P(J)=-P(X)$. By the chord law and
\ref{pf:reckernel:W3} for $(s,t)=(\mathcal F,J)$ (note $x(J)=x(X)\ne x(\mathcal F)$),
$P(\mathcal F(\mathcal F,J))=P(\mathcal F)+P(J)=P(Y)$; hence $\mathcal F(\mathcal F(X,Y),J)=Y$. Substituting $Y=0$, which fixes
$J$, and using $\mathcal F(X,0)=X$ gives $\mathcal F(X,J)=0$.
\end{proof}

Now let $K$ be complete for an ultrametric absolute value, and call $W$ \emph{integral} if $|a_i|\le1$ for all $i$. A
power series is \emph{integral} if all its coefficients have absolute value $\le1$; by \eqref{pf:reckernel:eq-w}, $w$ is
integral for integral $W$. For an integral series $G(X,Y)$ and $|a|,|b|<1$ the series
$G(a,b)=\sum_{i,j}c_{ij}a^ib^j$ converges, and $G\mapsto G(a,b)$ is a ring homomorphism on integral series which is
compatible with substitution of integral series without constant terms (the terms of total degree $n$ have absolute value
at most $\max(|a|,|b|)^n$). For a point $P$ of $W(K)$ put $u(P)=-x(P)/y(P)$ and $u(0)=0$; $P$ \emph{lies in the formal
group} if $P=0$ or $|x(P)|>1$.

\begin{theorem}[Points of the formal group]\label{pf:reckernel:fgHom}
Let $K$ be complete for an ultrametric absolute value, $W$ an integral Weierstrass curve over $K$, $w$ as in
\eqref{pf:reckernel:eq-w}, and $\mathcal F_2(X,Y)$ an integral series without constant term which satisfies the chord law
for $(X,Y)$. Then $u(P+Q)=\mathcal F_2(u(P),u(Q))$ for all points $P,Q$ of $W(K)$ lying in the formal group.
\formallabel{Formalism.ELL.formal_group_points_hom}
\formallabel{Formalism.ELL.formal_group_points_hom_of}
\formallabel{Formalism.ELL.FormalGroupPointsHom.e1hom_main}
\formallabel{Formalism.ELL.FormalGroupPointsHom.case_chord}
\formallabel{Formalism.ELL.FormalGroupPointsHom.case_tan}
\formallabel{Formalism.ELL.FormalGroupPointsHom.case_inv}
\end{theorem}

\begin{proof}
Let $\mathcal F$ be the formal group law of \ref{pf:reckernel:W2}. Both $\mathcal F_2$ and $\mathcal F$ have no constant term and
satisfy the chord law for $(X,Y)$, so $x(\mathcal F_2)=x(\mathcal F)$ and $y(\mathcal F_2)=y(\mathcal F)$, and hence
$\mathcal F_2=\mathcal F$. We use the parametrisation of the formal group
(\lref{Formalism.ELL.formal_group_points_param}): \emph{for $t\in K$ with $0<|t|<1$ one has $w(t)\ne0$, and
$(t/w(t),-1/w(t))$ is a non-singular point with $|x|>1$ and $u=t$; conversely every point $(x,y)$ of $W(K)$ with $|x|>1$
satisfies $0<|t|<1$ for $t=-x/y$, and $x=t/w(t)$, $y=-1/w(t)$.}

Let $a=u(P)$, $b=u(Q)$; then $|a|,|b|<1$. If $P=0$ or $Q=0$, the claim is $\mathcal F(0,b)=b$ or $\mathcal F(a,0)=a$,
which follow by evaluating $\mathcal F(0,Y)=Y$ and $\mathcal F(X,0)=X$. Otherwise $P=(a/w(a),-1/w(a))$ and
$Q=(b/w(b),-1/w(b))$ with $w(a),w(b)\ne0$. We distinguish three cases.

\emph{Chord case, $x(P)\ne x(Q)$.} Then $N=a\,w(b)-b\,w(a)\ne0$. Clearing the denominators $w(X)$, $w(Y)$,
$Xw(Y)-Yw(X)$ and $w(\mathcal F)$ (all non-zero, as their values at $(a,b)$ are $w(a)$, $w(b)$, $N$ and, by
\ref{pf:reckernel:W1}, $w(\mathcal F)\ne0$ because $\mathcal F\ne0$), the chord law for $(X,Y)$ becomes two polynomial identities
$\Phi_x=\Phi_y=0$ with coefficients in $\Z[a_1,\dots,a_6]$ between $X,Y,w(X),w(Y),\mathcal F$ and $w(\mathcal F)$; and by the same
computation in $K$, for $c,v\in K$ with $v\ne0$ the pair $(c/v,-1/v)$ is the chord sum $P+Q$ if and only if
$\Phi_x$ and $\Phi_y$ vanish at $(a,b,w(a),w(b),c,v)$. Evaluating at $(a,b)$ gives these identities for $c=\mathcal F(a,b)$ and
$v=w(\mathcal F)(a,b)$, together with the evaluation $v=c^3+a_1cv+\dots+a_6v^3$ of \ref{pf:reckernel:W1}. If $v=0$, the
latter gives $c=0$, and then $\Phi_y=0$ reduces to $N\,w(a)=0$, a contradiction. Hence $v\ne0$, $P+Q=(c/v,-1/v)$ and
$u(P+Q)=c$.

\emph{Inverse case, $x(P)=x(Q)$ and $y(P)=-y(Q)-a_1x(Q)-a_3$.} Then $Q=-P$, $P+Q=0$, and a direct computation gives
$b=u(-P)=-a/(1-a_1a-a_3w(a))$, where $1-a_1a-a_3w(a)\ne0$. Evaluating \cref{pf:reckernel:fglDblInv}(2) at $(a,a)$ gives
$\mathcal F(a,b)=0=u(P+Q)$.

\emph{Tangent case, $x(P)=x(Q)$ and $y(P)\ne-y(Q)-a_1x(Q)-a_3$.} Then $Q=P$, $b=a$, and the condition says
$-2+a_1a+a_3w(a)\ne0$; evaluating shows that $P(X)\ne-P(X)$ over $K[[X,Y]]$. By \cref{pf:reckernel:fglDblInv}(1), $\mathcal D=\mathcal F(X,X)$
satisfies $P(\mathcal D)=2P(X)$, i.e.\ the tangent formulas, and the argument of the chord case (with the tangent slope,
the polynomial identities of the doubling formulas and the denominator $-2+a_1X+a_3w(X)$) gives
$u(2P)=\mathcal D(a,a)=\mathcal F(a,a)$.
\end{proof}

\begin{theorem}[Addition estimate]\label{pf:reckernel:fgAdd}
Let $K$ be complete for an ultrametric absolute value and $W$ an integral Weierstrass curve over $K$. For all points
$P,Q$ of $W(K)$ lying in the formal group,
\[
|u(P+Q)-u(P)-u(Q)|\le\max(|u(P)|,|u(Q)|)^2 .
\]
\formallabel{ReciprocityTheory.ChartRestriction.L1.formalGroup_add_estimate}
\formallabel{ReciprocityTheory.ChartRestriction.L1.chord_eval_est}
\formallabel{ReciprocityTheory.ChartRestriction.L1.formalGroup_param_norm_lt_one}
\formallabel{ReciprocityTheory.ChartRestriction.L1.chord_norm_u_sq}
\end{theorem}

\begin{proof}
First, $|u(P)|<1$ for $P$ in the formal group. This is clear for $P=0$. For $P=(x,y)$ with $|x|>1$, the right-hand side
of the Weierstrass equation has absolute value $|x|^3$. If $|y|\le|x|$, the left-hand side would have absolute value
$\le|x|^2<|x|^3$; hence $|y|>|x|>1$, so $|y|^2$ dominates $|a_1xy|\le|x||y|$ and $|a_3y|\le|y|$, and $|y|^2=|x|^3$.
Therefore $|u(P)|^2|x|=|x|^3/|y|^2=1$ and $|u(P)|<1$.

Let $\mathcal F_2$ be as in \ref{pf:reckernel:W4}. Its coefficients lie in $\Z[a_1,\dots,a_6]$ and therefore have
absolute value $\le1$ (ultrametric inequality). By \cref{pf:reckernel:fgHom}, $u(P+Q)=\mathcal F_2(a,b)$ with $a=u(P)$,
$b=u(Q)$. Put $r=\max(|a|,|b|)<1$. The series $\mathcal F_2-X-Y$ has no terms of degree $\le1$, so
$\mathcal F_2(a,b)-a-b=\sum_{i+j\ge2}c_{ij}a^ib^j$ with $|c_{ij}|\le1$, and every term has absolute value
$\le r^{i+j}\le r^2$. By the ultrametric inequality, $|\mathcal F_2(a,b)-a-b|\le r^2$.
\end{proof}

\subsubsection{Formal group estimates}
Let $C$ be a complete field with an ultrametric absolute value and $W$ a Weierstrass curve over $C$ with integral
coefficients and $|\Delta(W)|=1$; $u(P)$ and the formal group are as in \cref{pf:reckernel:sec-fgl}.

\begin{lemma}[Formal group facts]\label{pf:reckernel:fgFacts}
Let $C$ and $W$ be as above.
\begin{enumerate}[label=(F\arabic*)]
\item\label{pf:reckernel:F1} The points lying in the formal group form a subgroup, and $|u(P)|<1$ for them.
\item\label{pf:reckernel:F2} For $P,Q$ in the formal group, $|u(P+Q)-u(P)-u(Q)|\le\max(|u(P)|,|u(Q)|)^2$.
\item\label{pf:reckernel:F3} Assume moreover that $W$ is short ($a_1=a_2=a_3=0$). A point not in the formal group is
  non-zero and has $|x|,|y|\le1$; and if $P,P'$ do not lie in the formal group, $P-P'\ne0$ does, and $P+P'$ does
  not, then $|x(P)-x(P')|\le|u(P-P')|$.
\end{enumerate}
\formallabel{ReciprocityTheory.ChartRestriction.L1.integral_of_not_mem_formalGroup}
\end{lemma}

\begin{proof}
The first assertion of \ref{pf:reckernel:F3} holds by definition for $x$, and then $|y|^2=|x^3+a_4x+a_6|\le1$ by
the ultrametric inequality. The inequality $|u(P)|<1$ in \ref{pf:reckernel:F1} and the estimate
\ref{pf:reckernel:F2} are proved in the first paragraph of the proof of \cref{pf:reckernel:fgAdd} and in
\cref{pf:reckernel:fgAdd}. The remaining assertions are classical \cite[Ch.~IV and \S VII.2]{Silverman2009} and are
proved in the development: the points lying in the formal group form a subgroup
(\lref{Formalism.ELL.formal_group_points_subgroup}: $0$ and the points with $|x|>1$ form a subgroup), and the
second assertion of \ref{pf:reckernel:F3} is
\lref{ReciprocityTheory.ChartRestriction.L1.norm_x_sub_le_of_formalGroup} (the stated inequality).
\end{proof}

\begin{lemma}\label{pf:reckernel:fgEst}
Let $C,W$ be as above and $D$ a point in the formal group.
\begin{enumerate}
\item For every $n\in\N$, $nD$ lies in the formal group and $|u(nD)-n\,u(D)|\le|u(D)|^2$.
\item If $M\in\N$ with $|M|<1$, then $|u(M^kD)|\to0$ as $k\to\infty$.
\end{enumerate}
\formallabel{ReciprocityTheory.ChartRestriction.L1.formalGroup_nsmul_estimate}
\formallabel{ReciprocityTheory.ChartRestriction.L1.formalGroup_pow_tendsto_zero}
\end{lemma}

\begin{proof}
(1) Induction on $n$, the case $n=0$ being trivial. Since $|u(D)|<1$ and $|n|\le1$, the induction hypothesis gives
$|u(nD)|\le\max(|u(nD)-nu(D)|,|nu(D)|)\le|u(D)|$, so by \ref{pf:reckernel:F2}
$|u(nD+D)-u(nD)-u(D)|\le|u(D)|^2$, and
$u((n+1)D)-(n+1)u(D)=\bigl(u(nD+D)-u(nD)-u(D)\bigr)+\bigl(u(nD)-nu(D)\bigr)$ has absolute value $\le|u(D)|^2$;
$(n+1)D$ lies in the formal group by \ref{pf:reckernel:F1}.
(2) Let $c=\max(|M|,|u(D)|)<1$. For $Q=M^kD$, (1) gives $|u(Q)|\le|u(D)|$ (by induction) and
$|u(MQ)|\le\max(|u(Q)|^2,|M||u(Q)|)\le c\,|u(Q)|$; hence $|u(M^kD)|\le c^k|u(D)|\to0$.
\end{proof}

\subsubsection{Prime-to-$\ell$ torsion}
Call $(g_2,g_3)\in\reckCl^2$ a \emph{good pair} if $\ell\ge5$, $g_2,g_3\in\reckKur$, $|g_2|,|g_3|\le1$ and
$|g_2^3-27g_3^2|=1$. Then $W=W_{g_2,g_3}$ satisfies the assumptions above over $C=\reckCl$, and its coefficients lie
in $\reckKur$, so that the set of points with coordinates in $\reckKur$ is a subgroup (the addition formulas are
rational with coefficients in $\reckKur$). By the second shared fact before
\cref{pf:reckernel:torsUnit}, a point $P\ne0$ with $NP=0$ and $\ell\nmid N$ does not lie in the formal group, hence
$|x(P)|,|y(P)|\le1$ by \ref{pf:reckernel:F3}.

\begin{lemma}[Lifting the reduction]\label{pf:reckernel:torLift}
Let $(g_2,g_3)$ be a good pair and $P=(x,y)$ a point of $W$ with $|x|\le1$ and $|y|=1$. There is a point $Q$ with
coordinates in $\reckKur$ such that $P-Q$ lies in the formal group.
\formallabel{ReciprocityTheory.ChartRestriction.L1.tor_lift}
\end{lemma}

\begin{proof}
Choose $a\in\reckKur$ ($0$ or a root of unity of order prime to $\ell$) with $|x-a|<1$. Then
$a^3-\frac{g_2}4a-\frac{g_3}4$ is congruent to $y^2$, so it has absolute value $1$ and lies in $\reckKur$; by
\cref{pf:reckernel:hensel} it has a square root $b\in\reckKur$ with $|b|=1$, and $Q=(a,b)$ is a point of $W$. If
neither $P-Q$ nor $P+Q$ lay in the formal group, then $P\pm Q\ne0$ would have integral coordinates
\ref{pf:reckernel:F3}, and the shared lemma \lref{Formalism.L2.Th1.E.UnitDiff_of} (for integral points
$P,Q\ne0$ of an integral model with unit discriminant such that $P\pm Q\ne0$ are integral, $|x(P)-x(Q)|=1$) would
give $|x-a|=1$, a contradiction. So $P-Q$ or $P-(-Q)$ lies in the formal group, and $-Q=(a,-b)$ has coordinates in
$\reckKur$ as well.
\end{proof}

\begin{theorem}[Torsion coordinates]\label{pf:reckernel:torKur}
Let $(g_2,g_3)$ be a good pair, $N\ge1$ with $\ell\nmid N$ and $P\ne0$ a point of $W_{g_2,g_3}$ with $NP=0$. Then
$x(P),y(P)\in\reckKur$.
\formallabel{ReciprocityTheory.ChartRestriction.L1.tor_kur}
\formallabel{ReciprocityTheory.ChartRestriction.L1.tor_x_mem}
\formallabel{ReciprocityTheory.ChartRestriction.L1.tor_x_step}
\formallabel{ReciprocityTheory.ChartRestriction.L1.tor_two}
\formallabel{ReciprocityTheory.ChartRestriction.L1.tor_y_unit}
\formallabel{ReciprocityTheory.ChartRestriction.L1.inSub_nsmul}
\end{theorem}

\begin{proof}
\emph{Case $2P=0$.} Then $y=0$, and $x$ is a root of $X^3-\frac{g_2}4X-\frac{g_3}4$ with $|x|\le1$ and
$|3x^2-\frac{g_2}4|=1$ (\lref{Formalism.L2.Th1.E.UnitKap_of}: for a non-zero $2$-torsion point of an integral
model with unit discriminant, the derivative $3x^2+a_4$ is a unit); so $x\in\reckKur$ by \cref{pf:reckernel:hensel}.

\emph{Case $2P\ne0$.} The points $P$ and $2P$ are non-zero and killed by $N$, so they are integral, and
$|2y|=1$ (\lref{Formalism.L2.Th1.E.UnitPsi2_of}: if $P$ and $2P\ne0$ are integral, the value $2y+a_1x+a_3$ is a
unit), i.e.\ $|y|=1$. Let $Q$ be as in
\cref{pf:reckernel:torLift} and $D=P-Q$. Let $M=\ell^{\varphi(N)}$, so that $|M|<1$ and $n_k:=M^k\equiv1\pmod N$. Then
$n_kP=P$ and $n_kQ=P-n_kD$. If $n_kD=0$, then $x(n_kQ)=x(P)$. Otherwise $n_kD$ lies in the formal group
(\cref{pf:reckernel:fgEst}), $n_kQ=P-n_kD$ does not (else $P$ would), and $P+n_kQ=2P-n_kD$ does not (as $2P$ does not);
by \ref{pf:reckernel:F3}, $|x(P)-x(n_kQ)|\le|u(n_kD)|$. By \cref{pf:reckernel:fgEst}(2) the right-hand side tends
to $0$. Since $n_kQ$ has coordinates in $\reckKur$ and $\reckKur$ is closed, $x(P)\in\reckKur$. Finally
$y^2=x^3-\frac{g_2}4x-\frac{g_3}4\in\reckKur$ and $|y|=1$, so $y\in\reckKur$ by \cref{pf:reckernel:hensel}.
\end{proof}

\begin{theorem}[Torsion integrality]\label{pf:reckernel:cm4}
\textup{(TI)} holds.
\formallabel{ReciprocityTheory.ChartRestriction.L1.cm4}
\formallabel{ReciprocityTheory.ChartRestriction.L1.cmTorsion_of_point}
\end{theorem}

\begin{proof}
Let $\lambda$ be such that $(g_2',g_3')=(\lambda^4\iota(g_2(L)),\lambda^6\iota(g_3(L)))$ is a good model, and let
$z,N$ be as in \textup{(TI)}. By \cref{pf:reckernel:scaledTorsion}, $P=(\lambda^2\iota(\wp(z)),\frac12\lambda^3
\iota(\wp'(z)))$ is a non-zero point of $W_{g_2',g_3'}$ with $NP=0$. By \cref{pf:reckernel:torKur} its coordinates lie
in $\reckKur$, and they are integral since $P$ is not in the formal group. As $|2|=1$, the same holds for
$\lambda^3\iota(\wp'(z))=2y(P)$.
\end{proof}

\begin{remark}
The proof of (TI) uses neither the complex multiplication of $L$ nor the hypotheses on $d$ and
$\ell\nmid d$: it holds for every lattice admitting a good model in the sense of \cref{pf:reckernel:def-CM}.
\end{remark}

\subsubsection{The formal group at inert primes}

Let $K$ be a complete field with an ultrametric absolute value and $\ell$ a prime. A \emph{Frobenius lift} on $K$ is a
ring endomorphism $\sigma$ with $|\sigma x|=|x|$ for all $x$ and $|\sigma x-x^\ell|\le|\ell|$ for $|x|\le1$, where
$|\ell|<1$ and every $x$ with $|x|<1$ has $|x|\le|\ell|$. Then also $|\sigma^2x-x^{\ell^2}|\le|\ell|$ for $|x|\le1$:
$\sigma^2x-x^{\ell^2}=\sigma(\sigma x-x^\ell)+\bigl((\sigma x)^\ell-(x^\ell)^\ell\bigr)$, and
$a^\ell-b^\ell$ is divisible by $a-b$ with integral cofactor. A power series $f\in K[[X]]$ is \emph{integral} if all its
coefficients have absolute value $\le1$, and $f$ satisfies the \emph{functional equation} (for $\ell$ and $\sigma^2$)
if $f(X)-\ell^{-1}(\sigma^2f)(X^{\ell^2})$ is integral, where $\sigma^2f$ is obtained by applying $\sigma^2$ to the
coefficients. The Lubin--Tate logarithm $\lambda_\ell$ satisfies it, since $\lambda_\ell(X)-\ell^{-1}\lambda_\ell(X^{\ell^2})=X$.

\begin{lemma}[Lang]\label{pf:reckernel:lang}
Let $\sigma$ be a Frobenius lift on $K$, and assume that every monic polynomial of positive degree with integral
coefficients has an integral root modulo $\ell$ (i.e.\ an $x$ with $|x|\le1$ and $|P(x)|\le|\ell|$). Then for every
$c\in K$ with $|c|=1$ there is $\eps\in K$ with $|\eps|=1$ and $\eps=c\,\sigma^2(\eps)$.
\formallabel{ReciprocityTheory.ChartRestriction.L1.lang_semilinear}
\formallabel{ReciprocityTheory.ChartRestriction.L1.semilin_root}
\end{lemma}

\begin{proof}
We construct $\eps_n$ with $|\eps_n|=1$, $|\eps_n-c\sigma^2(\eps_n)|\le|\ell|^{n+1}$ and
$|\eps_{n+1}-\eps_n|\le|\ell|^{n+1}$. For $n=0$ choose $x$ integral with $|x^{\ell^2-1}-c^{-1}|\le|\ell|$ (a root modulo
$\ell$ of $X^{\ell^2-1}-c^{-1}$); then $|x|=1$ and
$x-c\sigma^2(x)=-cx\bigl(x^{\ell^2-1}-c^{-1}\bigr)-c\bigl(\sigma^2x-x^{\ell^2}\bigr)$ has absolute value $\le|\ell|$; put
$\eps_0=x$. Given $\eps_n$, write $\eps_n-c\sigma^2(\eps_n)=\ell^{n+1}\rho$ with $|\rho|\le1$, choose $\delta$ integral
with $|\delta^{\ell^2}-c^{-1}\delta-c^{-1}\rho|\le|\ell|$, and put $\eps_{n+1}=\eps_n+\ell^{n+1}\delta$. Then
$|\eps_{n+1}|=1$, $|\eps_{n+1}-\eps_n|\le|\ell|^{n+1}$, and since $\sigma^2(\ell^{n+1})=\ell^{n+1}$,
\[
\eps_{n+1}-c\sigma^2(\eps_{n+1})=\ell^{n+1}\bigl(\rho+\delta-c\sigma^2\delta\bigr)
=\ell^{n+1}\Bigl(-c\bigl(\delta^{\ell^2}-c^{-1}\delta-c^{-1}\rho\bigr)-c\bigl(\sigma^2\delta-\delta^{\ell^2}\bigr)\Bigr),
\]
of absolute value $\le|\ell|^{n+2}$. The sequence $(\eps_n)$ is Cauchy; its limit $\eps$ has $|\eps|=1$ and, $\sigma$
being continuous, $\eps=c\sigma^2(\eps)$.
\end{proof}

\begin{lemma}[From Honda's congruence to a Lubin--Tate isomorphism]\label{pf:reckernel:finIntegral}
Let $\sigma$ be a Frobenius lift on $K$, $\lambda,e\in K[[X]]$ with $\lambda(0)=0$, $\lambda'(0)=1$, $e(0)=0$ and
$e(\lambda(X))=X$, and let $c,\eps\in K$ with $|c|=|\eps|=1$, $\eps=c\,\sigma^2(\eps)$ and
$\lambda(X)-\frac c\ell(\sigma^2\lambda)(X^{\ell^2})$ integral. Then $e(\eps\lambda_\ell(X))$ is integral.
\formallabel{ReciprocityTheory.ChartRestriction.L1.fin_integral}
\end{lemma}

\begin{proof}
Put $f_0=\eps^{-1}\lambda$. Since $\sigma^2(\eps^{-1})=c\,\eps^{-1}$,
$f_0(X)-\ell^{-1}(\sigma^2f_0)(X^{\ell^2})=\eps^{-1}\bigl(\lambda(X)-\frac c\ell(\sigma^2\lambda)(X^{\ell^2})\bigr)$ is
integral, so $f_0$ satisfies the functional equation. Then so does $f_1(X)=f_0(\eps X)=\eps^{-1}\lambda(\eps X)$, by the
substitution part of Hazewinkel's functional equation lemma \cite{Hazewinkel1978}
(\lref{Formalism.LT.funEqSubst}: substituting an integral series without constant term into a series satisfying
the functional equation gives a series satisfying it); and $f_1(0)=0$, $f_1'(0)=1$. By the main part of the functional
equation lemma (\lref{Formalism.LT.funEqHom}: if $f,g$ satisfy the functional equation, $f(0)=g(0)=0$ and
$f'(0)=1$, then $f^{-1}\circ g$ is integral), applied to $f_1$ and $\lambda_\ell$, there is an integral $h$ with
$h(0)=0$ and $f_1(h(X))=\lambda_\ell(X)$. Hence $\lambda(\eps h(X))=\eps f_1(h(X))=\eps\lambda_\ell(X)$, and
$e(\eps\lambda_\ell(X))=\eps h(X)$ is integral.
\end{proof}

\begin{theorem}[Lubin--Tate twist]\label{pf:reckernel:cm3}
\textup{(LT)} holds.
\formallabel{ReciprocityTheory.ChartRestriction.L1.cmIntegralityThree}
\formallabel{ReciprocityTheory.ChartRestriction.L1.cm3}
\formallabel{GlobalTheory.HondaCM3.cm3}
\end{theorem}

The development contains two independent formal proofs of (LT): the first is used for (Mod)
and (Dict), the second for (Ex). Both follow Honda's theory of formal groups \cite{Honda1970}; we describe the first.
Its last two steps are \cref{pf:reckernel:lang,pf:reckernel:finIntegral}; the formal proofs of the other ingredients
(about $3\cdot10^3$ lines) are not converted here.

\begin{proof}[Proof (structure)]
Let $(g_2',g_3')$ be a good model at an inert prime $\ell\ge5$, $K=\reckKur$, and let $\lambda(X)\in K[[X]]$ be the
formal logarithm of $W_{g_2',g_3'}$, i.e.\ the compositional inverse of $u=u_{g_2',g_3'}$; let $\mathcal F$ be the
formal group law with logarithm $\lambda$, which has integral coefficients
(\lref{ReciprocityTheory.ChartRestriction.L1.curve_law}). The ingredients are:
\begin{enumerate}[label=(\alph*)]
\item $K$ is complete, has algebraically closed residue field, and carries a \emph{Frobenius lift} $\sigma$: an
  isometric ring endomorphism with $|\sigma x-x^\ell|\le|\ell|$ for $|x|\le1$
  (\lref{ReciprocityTheory.ChartRestriction.L1.kurFrob} and \lref{ReciprocityTheory.ChartRestriction.L1.kur_resClosed}).
\item $a=\iota(\omega_d)$ lies in $K$, $|a|\le1$, $|1-2a|=1$ and $\sigma(a)=1-a$, because $\ell$ is inert in $K_d$
  (\lref{ReciprocityTheory.ChartRestriction.L1.kur_quad}).
\item The endomorphisms $[a]=u(a\lambda(X))$ and $[\ell]=u(\ell\lambda(X))$ of $\mathcal F$ have integral
  coefficients (\lref{ReciprocityTheory.ChartRestriction.L1.end_integral} and \lref{ReciprocityTheory.ChartRestriction.L1.mulN_integral}),
  and $[\ell]$ has a coefficient of absolute value $1$ in some degree $\le\ell^2$ (height at most $2$;
  \lref{ReciprocityTheory.ChartRestriction.L1.level_le}).
\item (Honda) Under (a)--(c) there is $c\in K$ with $|c|=1$ such that
  $\lambda(X)-\frac c\ell(\sigma^2\lambda)(X^{\ell^2})$ has integral coefficients
  (\lref{ReciprocityTheory.ChartRestriction.L1.hon_main}).
\item (Lang) There is $\eps\in K$ with $|\eps|=1$ and $\eps=c\,\sigma^2(\eps)$ (\cref{pf:reckernel:lang}, whose
  hypothesis on roots modulo $\ell$ holds by (a)).
\item For such $c,\eps$, the series $u(\eps\lambda_\ell(X))$ has integral coefficients
  (\cref{pf:reckernel:finIntegral}, with $e=u$).
\end{enumerate}
The coefficients of $u(\eps\lambda_\ell(X))$ lie in $K$ because $g_2',g_3',\eps\in K$; with (f) this is
(LT). The second proof replaces (c)--(f) by an integrality statement for the Honda logarithm pair of a
short Weierstrass model, the integrality of the multiplication by $\omega_d$ at inert primes, a height bound for good
reduction, and a ``unit multiplier'' theorem for formal groups of height two
(\lref{GlobalTheory.T.isIntegralHondaLogarithmPair_shortWeierstrass}, \lref{GlobalTheory.T.cmMultiplierIntegralSeriesAtInertPrime_proof}, \lref{GlobalTheory.T.formalGroup_height_le_two_good_reduction}, and \lref{GlobalTheory.T.heightTwo_formalGroup_unramified_multiplier}).
\end{proof}

% ---------------------------------------------------------------------------
\subsection{CM theta data}\label{pf:reckernel:sec-theta}

\subsubsection{Lattices of ideals}
Let $D$ be a datum of the construction. Its host provides an imaginary quadratic field $K=\Q(\sqrt{-d})$ with $d$ a
prime, $d\equiv3\pmod 4$, a CM field $E\supseteq K$ with a fixed embedding $\varphi_0\colon E\to\C$, and a finite set
$\mathcal T$ of embeddings of $E$. For an ideal $\mathfrak c\subseteq\OO_K$ and $\delta\in\C$ put
\[
\Lambda(\mathfrak c,\delta)=\delta\cdot\varphi_0(\mathfrak c)\subset\C .
\]
We call $\mathfrak c$ \emph{admissible} if $\mathfrak c\ne0$ and $\Nm\mathfrak c$ is prime to $2d$, and $\delta$ a
\emph{normalising root} for $\mathfrak c$ if $\delta^{12}=\Delta(\Lambda(\mathfrak c,1))$; then $\delta\ne0$,
$\Delta(\Lambda(\mathfrak c,\delta))=1$ (as $\Delta(cL)=c^{-12}\Delta(L)$), and for $\mathfrak c\ne0$, $\delta\ne0$ the group
$\Lambda(\mathfrak c,\delta)$ is a lattice with complex multiplication by $\omega_d$ ($\mathfrak c$ is free of rank $2$
over $\Z$ with discrete image in $\C$, and $\omega_d\mathfrak c\subseteq\mathfrak c$). A \emph{CM base} over $(\mathfrak c,\delta)$
(\lref{Formalism.L2.Ker.Base}: a structure of the development) consists of a period pair for
$\Lambda(\mathfrak c,\delta)$, a point $t_0$ with $2t_0\in\Lambda$, $\omega_dt_0\notin\Lambda$ and
$\bar\omega_dt_0\notin\Lambda$ (a $2$-torsion point not killed by $\omega_d$ or $\bar\omega_d$), an auxiliary isomorphism
$\gamma_0$ of the type of $K$-isomorphisms of the development with a compatible complex-linear automorphism of
$\C^{4\times\mathcal T}$, and a choice attached to $\gamma_0$. A \emph{CM character} is a function $\chi$ on
fractional ideals of $E$, multiplicative on ideals of norm one, with $\chi(z\OO_E)=\prod_{\tau\in\mathcal T}\tau(z)^{-1}$
for $z\in E$ with $z\bar z=1$.

\paragraph{Theta functions.}
For even $N,N'$ and self-dual hermitian $\OO_E$-lattices $L\subset E^N$, $L'\subset E^{N'}$ (such a quadruple is
called \emph{admitted}) let $V=(\C^{N'\times N})^{\mathcal T}$, let $\Lambda_{L,L'}\subset V$ be the group generated
by the vectors $w\cdot(\tau(A))_\tau$ with $w\in\Lambda(\mathfrak c,\delta)$ and $A$ in the $\OO_E$-span $Y(L,L')$ of the
matrices $aA$, $a$ in a fixed ideal of the host and $AL\subseteq L'$, let $H(z,w)=\sum z\bar w/\mathrm{area}(\Lambda)$,
and for $x\in V$ let $(U_xG)(z)=G(z+x)\exp(-\pi H(z,x)-\pi H(x,x)/2)$. A CM base $B$ and a CM character $\chi$ determine
a family of functions $e_\chi=e_\chi(N,N',L,L')$ on $V$ (a $\chi$-factor times a translate by the marked point of the
theta function of $\Lambda_{L,L'}$; \lref{Formalism.L2.Ker.Base.eNat}).

\subsubsection{Curves at inert primes}
Let $p$ be a prime and $\lambda\colon\Qbar\to\C_p$ a field embedding, extended to $\C$ by $\lambda(x)=0$ for
$x\notin\Qbar$. Let $W_p\subset\C_p$ be the closure of the ring generated by the roots of unity of order prime to
$p$ (the ring of integers of $\reckKurp$). For a lattice $\Lambda$ with algebraic invariants let
$W_\lambda(\Lambda)=W_{\lambda(g_2(\Lambda)),\lambda(g_3(\Lambda))}$. A \emph{$p$-adic curve datum} $g$ consists of
$\lambda_g$, a lattice $\Lambda_g$, a change of variables $C_g$ over $\C_p$ and two power series $w_g,x_g$; put
$W_g=C_g\cdot W_{\lambda_g}(\Lambda_g)$ and let $\ell_g$ be its formal logarithm (computed from $w_g$).

\begin{definition}[Curves in position]\label{pf:reckernel:def-position}
Let $D,\mathfrak c,\delta$ be as above and $g$ a $p$-adic curve datum.
\begin{enumerate}
\item $g$ is \emph{in position at $(\mathfrak c,\delta)$} if $\mathfrak c$ is admissible, $\Nm\mathfrak c$ is prime to
  $p$, $\delta$ is a normalising root, $\Lambda_g=\Lambda(\mathfrak c,\delta)$, $W_g$ has integral coefficients and
  unit discriminant, $w_g$ is the invariant-differential series of $W_g$, $x_g(0)=0$ and $x_g\circ\ell_g=X$.
\item $g$ is \emph{defined over the Witt vectors} if the coefficients of $W_g$ lie in $W_p$.
\item $g$ has a \emph{Lubin--Tate isomorphism} if there are $\eps\in W_p^\times$ and $\theta\in W_p[[X]]$ with
  $\theta(0)=0$ and $\theta'(0)\in W_p^\times$ such that $\lambda_p(\theta(X))=\eps\,\ell_g(X)$, where
  $\lambda_p(X)=\sum_{j\ge0}X^{p^{2j}}/p^j$; here $\eps$ and the coefficients of $\theta$ are taken in the ring of
  Witt vectors of $\overline{\F}_p$, which the development embeds into $\C_p$ with image $W_p$.
\end{enumerate}
\formallabel{Formalism.L2.G.Pend.LTIso}
\end{definition}

\begin{definition}[Two existence statements]\label{pf:reckernel:def-Mod}
\begin{enumerate}
\item \textup{(Mod)} For every datum $D$, admissible $\mathfrak c$ with normalising root $\delta$, prime $p\ge5$ inert in
  $K$ and embedding $\lambda\colon\Qbar\to\C_p$, there are an admissible $\mathfrak c'$ with normalising root $\delta'$
  and $\Lambda(\mathfrak c',\delta')=\Lambda(\mathfrak c,\delta)$, and a $p$-adic curve datum $g$ with $\lambda_g=\lambda$
  which is in position at $(\mathfrak c',\delta')$, defined over the Witt vectors and has a Lubin--Tate isomorphism.
\item \textup{(Ex)} For every datum $D$, admissible $\mathfrak c$ with normalising root $\delta$, CM base $B$ over
  $(\mathfrak c,\delta)$, CM character $\chi$, prime $p\ge5$ inert in $K$ with $p\nmid\Nm\mathfrak c$ and embedding
  $\lambda$, there is a $p$-adic curve datum $g$ with a Lubin--Tate isomorphism which is \emph{in position for
  $(B,\chi,\lambda)$}: $g$ is in position at $(\mathfrak c,\delta)$ and defined over the Witt vectors,
  $\lambda_g=\lambda$, and $\Lambda_g$ is the lattice of $B$.
\end{enumerate}
\formallabel{Formalism.L2.KerInert.GenusOneModelExistsLattice}
\formallabel{GlobalTheory.TwV.Lt.CurveEx}
\formallabel{Formalism.L2.KerInert.AdmLoc.CurveIn}
\end{definition}

\begin{lemma}\label{pf:reckernel:g2Alg}
Assume \textup{(GM)}. If $L$ has complex multiplication by $\omega_d$ and $\Delta(L)=1$, then $g_2(L)$ is algebraic.
\formallabel{GlobalTheory.InertGenusOne.isAlgebraic_g2_of_hasCMBy_of_disc_eq_one}
\formallabel{GlobalTheory.InertGenusOne.exists_ringEquiv_transcendental}
\end{lemma}

\begin{proof}
Fix a prime $p\ge5$ with $p\nmid d$. If $g_2(L)$ were transcendental, there would be a field isomorphism
$\iota\colon\C\to\C_p$ with $|\iota(g_2(L))|>1$ ($\C$ and $\C_p$ are algebraically closed of characteristic $0$ and of the same
uncountable transcendence degree, so a transcendental number can be sent to any transcendental element of $\C_p$,
for instance to one of absolute value $>1$). Let $(g_2',g_3')$, $\mu$ be a good
model at $(p,\iota)$. Then $g_2'^3-27g_3'^2=\mu^{12}\iota(\Delta(L))=\mu^{12}$ has absolute value $1$, so $|\mu|=1$ and
$|g_2'|=|\iota(g_2(L))|>1$, a contradiction.
\end{proof}

\begin{lemma}[Curve data for ideals prime to $p$]\label{pf:reckernel:datumCoprime}
Assume \textup{(GM)} and \textup{(LT)}. Let $\mathfrak c$ be admissible with normalising root $\delta$, $p\ge5$ a prime
inert in $K$ with $p\nmid\Nm\mathfrak c$, and $\lambda\colon\Qbar\to\C_p$. Then there is a $p$-adic curve datum $g$ with
$\lambda_g=\lambda$ which is in position at $(\mathfrak c,\delta)$, defined over the Witt vectors and has a Lubin--Tate
isomorphism.
\formallabel{GlobalTheory.InertGenusOne.datum_exists_coprime}
\formallabel{Formalism.L2.KerInert.GenusOneModelExists}
\formallabel{GlobalTheory.InertGenusOne.lamOf_lat}
\end{lemma}

\begin{proof}
Let $L$ be a period pair with lattice $\Lambda(\mathfrak c,\delta)$; it has complex multiplication by $\omega_d$, and
$\Delta(L)=1$ because $\delta$ is a normalising root. As $p$ is inert, $-d$ is not a square modulo $p$ and $p\nmid d$.
By \cref{pf:reckernel:g2Alg}, $g_2(L)$ is algebraic, and so is $g_3(L)$ since $27g_3(L)^2=g_2(L)^3-1$. Extend
$\lambda$ to a field isomorphism $\iota\colon\C\to\C_p$
(\lref{GlobalTheory.PadicComplexIso.exists_ringEquiv_extending}: every embedding of $\Qbar$ into $\C_p$ extends
to an isomorphism $\C\to\C_p$). By \textup{(GM)} there is a good model $(g_2',g_3')$ of $L$ at $(p,\iota)$ with
$g_2'=\mu^4\iota(g_2(L))$, $g_3'=\mu^6\iota(g_3(L))$, and by \textup{(LT)} there is $\eps\in\reckKurp$, $|\eps|=1$, such
that $u_{g_2',g_3'}(\eps\lambda_p(X))$ has integral coefficients in $\reckKurp$. Let $C$ be the change of variables
$(x,y)\mapsto(\mu^2x,\mu^3y)$; then $C\cdot W_\lambda(\Lambda(\mathfrak c,\delta))=W_{g_2',g_3'}$, whose coefficients lie in
$W_p$. From $(g_2',g_3',\mu,\eps)$ one builds the curve datum $g$ with $C_g=C$ and verifies the three properties
(\lref{GlobalTheory.InertGenusOne.datum_of_short}: for a short model with integral coefficients in $W_p$ and unit
discriminant and a unit $\eps$ with integral $u(\eps\lambda_p(X))$, the invariant-differential series, the inverse of
the formal logarithm and a Lubin--Tate isomorphism exist).
\end{proof}

\begin{theorem}\label{pf:reckernel:modEx}
\textup{(GM)} and \textup{(LT)} imply \textup{(Ex)}; \textup{(GM)}, \textup{(LT)} and \textup{(TI)} imply \textup{(Mod)}.
\formallabel{GlobalTheory.InertGenusOne.datum_exists_lattice}
\formallabel{GlobalTheory.TwV.Lt.curveEx_of_cm}
\formallabel{Formalism.L2.KerInert.curve_ideal_prime}
\end{theorem}

\begin{proof}
(Ex) is \cref{pf:reckernel:datumCoprime} applied to $(\mathfrak c,\delta)$, together with the hypotheses of (Ex). For
(Mod) (the hypothesis (TI) is not used), let $p$ be inert; then $p\OO_K$ is prime, and $\mathfrak c=p^k\mathfrak c'$ with $p\nmid\Nm\mathfrak c'$. Put
$\delta'=\delta p^k$. Then $\Lambda(\mathfrak c,\delta)=\Lambda(\mathfrak c',\delta')$, $\mathfrak c'$ is admissible
($\Nm\mathfrak c'$ divides $\Nm\mathfrak c$), and $\delta'^{12}=\delta^{12}p^{12k}=\Delta(\Lambda(\mathfrak c,1))\,p^{12k}
=\Delta(\Lambda(\mathfrak c',1))$ because $\Lambda(\mathfrak c,1)=p^k\Lambda(\mathfrak c',1)$. Now apply
\cref{pf:reckernel:datumCoprime} to $(\mathfrak c',\delta')$.
\end{proof}

\subsubsection{The dictionary, algebraicity and integrality}
The remaining statements concern the theta functions $e_\chi$ and their $p$-adic expansions. We state them in the
form in which they enter, and give their derivations from the CM statements; the proofs of the main lemmas are not
converted here.

\begin{definition}[Three statements on theta data]\label{pf:reckernel:def-three}
\begin{enumerate}
\item \textup{(Alg)} For all $D$, admissible $\mathfrak c$ with normalising root $\delta$, CM bases $B$, CM characters
  $\chi$, admitted $(N,N',L,L')$, all $t\in V$ with $at\in\Lambda_{L,L'}$ for some integer $a\ge1$, all Taylor
  coefficients at $0$ of $U_te_\chi(N,N',L,L')$ are algebraic numbers.
\item \textup{(Dict)} For all $D,\mathfrak c,\delta$, primes $\ell\ge5$ inert in $K$ and $\ell$-adic curve data $g$ in
  position at $(\mathfrak c,\delta)$, defined over the Witt vectors and with a Lubin--Tate isomorphism (unit
  $\eps_g$), there are $g_2,g_3,\eps,s_2\in\reckKur$ such that $g_2,g_3$ are the invariants of $W_g$,
  $\eps=\eps_g^{-1}$, the $\ell$-adic image of the Taylor series of the normalised sigma function of $\Lambda_g$,
  rescaled by the inverse of the scaling unit $u_g$ of $C_g$, is $u_g^{-1}$ times the sigma series of $W_g$ with
  parameter $s_2$, $|u_g^{-1}\eps|=1$, the curve condition of the development holds for $(g_2,g_3,\eps,s_2)$
  relative to the Lubin--Tate isomorphism, and the dictionary integrality condition holds for every datum.
\item \textup{(Int)} For all $D,\mathfrak c,\delta,B,\chi$ as in (Alg), primes $p\ge5$ inert in $K$, embeddings
  $\lambda$ and admitted $(N,N',L,L')$: for every $p$-adic power series $F$ in the coordinates of $V$ that is a
  combination of the expansions of $e_\chi$ at torsion points, every change of variables making
  $W_\lambda(\Lambda(\mathfrak c,\delta))$ integral with unit discriminant, its invariant-differential series $w$, and
  every local orthonormal basis $b$ of $Y(L,L')$ at $p$, the series obtained from $F$ by substituting the formal
  logarithm coordinates (via $w$ and $b$), multiplied by the associated exponential factor, has integral
  coefficients.
\end{enumerate}
\formallabel{Formalism.L2.KerInert.Rep.CurveDict}
\formallabel{Formalism.L2.KerInert.IntInert}
\end{definition}

\begin{theorem}\label{pf:reckernel:threeFromCM}
\textup{(Alg)} holds, and \textup{(GM)}, \textup{(LT)} and \textup{(TI)} imply \textup{(Mod)} and \textup{(Dict)}.
\formallabel{GlobalTheory.CurveThetaInert.torsionTaylorAlgebraic_latticeExists_of_cm}
\formallabel{GlobalTheory.CurveDictionary.curveDict_of_cm}
\end{theorem}

\begin{proof}
(Alg) is proved unconditionally in the development (\lref{GlobalTheory.ThLoc.algTAll}: the statement (Alg)).
(Mod) is \cref{pf:reckernel:modEx}. (Dict) is deduced from (GM) and (TI) (the hypothesis (LT) is not used) together with eight lemmas on CM curves at
inert primes (\lref{GlobalTheory.CurveDictionary.curveDict_of_inputs}): algebraicity of the sigma invariant of
the lattice when its invariants are algebraic, well-formedness of the transformed coefficients of unramified
Weierstrass curves, invertibility of $\ell$ on the coefficients of the model, a coboundary statement for logarithmic
coordinates and its compatibility with the equivalence of series, a comparison of the theta unit function of the short
model with the Weierstrass data and its integrality, and the integrality of the inverse series of a formal-group
isomorphism (\lref{GlobalTheory.T.curveLattice_sigmaTwo_algebraic_of_invariants_algebraic}, \lref{GlobalTheory.T.thetaUnitFunction_shortModel_integral}, and \lref{GlobalTheory.T.inverse_series_formalGroup_iso_integral}).
\end{proof}

\subsubsection{Integrality at inert primes: the reduction}\label{pf:reckernel:sec-intred}
We make the statement (Int) precise and reduce it to three statements on approximation and orthonormal bases. Fix
$D,\mathfrak c,\delta,B,\chi$ as in (Alg), a prime $p\ge5$ inert in $K$ (that is, $\left(\frac{-d}{p}\right)=-1$) and an
embedding $\lambda\colon\Qbar\to\C_p$, extended to $\C$ by $0$ outside $\Qbar$.

\paragraph{Notation.} For $N,N'$ and $\OO_E$-lattices $L\subset E^N$, $L'\subset E^{N'}$ let $\Lambda_{L,L'}\subset V$ be as in
\cref{pf:reckernel:sec-theta}. A point $t\in V$ is \emph{$p'$-torsion} for $(L,L')$ if $nt\in\Lambda_{L,L'}$ for some
$n\ge1$ prime to $p$. We write $L\sim_pL_1$ if $sL\subseteq L_1$ and $sL_1\subseteq L$ for some $s$ prime to $p$. For a
function $G$ on $V$ analytic at $0$, $G_\lambda$ denotes the $p$-adic power series in the variables $X_{\tau,i',i}$
($\tau\in\mathcal T$, $i'\le N'$, $i\le N$) obtained by applying $\lambda$ to the Taylor coefficients of $G$ at $0$. Let
$\theta^0_{N,N'}$ be the theta function of the standard lattices $(\mathrm{std}_N,\mathrm{std}_{N'})$ attached to $B$ by the
development, so that $e_\chi(N,N',\mathrm{std}_N,\mathrm{std}_{N'})$ is built from it. A family $b_1,\dots,b_r$ in an
$\OO_E$-submodule $Y$ of $E^{N'\times N}$ is a \emph{local basis at $p$} of $Y$ if it is linearly independent over $K$, and
$sb_i\in Y$ and $sy\in\sum_i\OO_Kb_i$ for all $y\in Y$ with suitable $s$ prime to $p$; it is \emph{orthonormal} if
$h_Y(b_i,b_j)=\delta_{ij}$ for the hermitian form $h_Y$ of $Y$. A \emph{model} is a change of variables $C$ over $\C_p$ such
that $C\cdot W_\lambda(\Lambda(\mathfrak c,\delta))$ has $p$-integral coefficients and discriminant $1$; for it let $\ell_C$
be the formal logarithm of $C\cdot W_\lambda(\Lambda(\mathfrak c,\delta))$, $u_C$ the scaling factor of $C$, and
\[
\Phi_C(X)=\frac{\sigma_\lambda(u_C^{-1}\ell_C(X))}{u_C^{-1}X},
\]
where $\sigma_\lambda$ is the $\lambda$-image of the Taylor series of the Weierstrass sigma function of
$\Lambda(\mathfrak c,\delta)$. For a family $b=(b_1,\dots,b_r)$ and a series $F$ in the variables $X_{\tau,i',i}$ let
$F^{C,b}\in\C_p[[Y_1,\dots,Y_r]]$ be obtained by substituting
$X_{\tau,i',i}\mapsto\sum_iu_C^{-1}\lambda(\tau(b_i)_{i'i})\,\ell_C(Y_i)$. We call $F$ \emph{$(L,L')$-integral} if for every
model $C$ and every orthonormal local basis $b$ of $Y(L,L')$ at $p$ the series
$F^{C,b}\cdot\prod_{i=1}^r\Phi_C(Y_i)^{-1}$ has $p$-integral coefficients. Finally, $\mathrm{Comb}(L,L')$ is the set of
series $\sum_kr_k\,(U_{t_k}e_\chi(N,N',L_k,L'_k))_\lambda$ (finite sums) with $|r_k|\le1$, $t_k$ $p'$-torsion for
$(L,L')$, and $(N,N',L_k,L'_k)$ admitted with $L_k\sim_pL$, $L'_k\sim_pL'$. In this notation \textup{(Int)} is the
statement: \emph{for all data as above and all admitted $(N,N',L,L')$, every element of $\mathrm{Comb}(L,L')$ is
$(L,L')$-integral.}

A \emph{conjugation} is an isometric ring endomorphism $\sigma$ of $\C_p$ inducing complex conjugation on $\lambda(K)$
(that is, $\sigma(\lambda(x))=\lambda(\bar x)$ for $x\in K$). For a conjugation $\sigma$, a family
$v=(v_1,\dots,v_r)$ of vectors in $\C_p^{\mathcal T\times N'\times N}$ is a \emph{local orthonormal frame} of $Y$ at $p$ if
$v_i=\sum_mA_{im}\lambda(\tau(b_m))$ for a local basis $b$ of $Y$ at $p$ and a matrix $A$ with $p$-integral entries in the
closure of $\lambda(K)$, $|\det A|=1$ and $A\,\lambda(h_Y(b_m,b_n))_{m,n}\,{}^t\sigma(A)=1$. A series $F$ is
\emph{integral in the coordinates $v$} if $F^{C,v}\cdot\prod_i\Phi_C(Y_i)^{-1}$ is integral for every model $C$, where
$F^{C,v}$ is defined as $F^{C,b}$ with $u_C^{-1}v_i(\tau,i',i)$ in place of $u_C^{-1}\lambda(\tau(b_i)_{i'i})$.

\begin{definition}[Approximation and frame statements]\label{pf:reckernel:def-WA}
\begin{enumerate}
\item \textup{(WA)} is the weak approximation statement for hermitian lattices of
  \lref{Formalism.L2.KerInert.WAPlaces}: \emph{if a lattice $L$ in a standard hermitian space over a CM field (of the
  type of the hosts) is
  locally isometric to the standard lattice at every prime above $p$, then a single global isometry carries the standard
  lattice to $L$ locally at all primes above $p$.}
\item \textup{(ONB)}: for every datum $D$ and every inert prime $p\ge5$ there are elements $e_1,\dots,e_m$ of
  $Y(\mathrm{std}_1,\mathrm{std}_1)$ forming a local basis at $p$, pairwise orthogonal for $h_Y$, with
  $h_Y(e_i,e_i)$ a positive rational number whose numerator and denominator are prime to $p$.
\item \textup{(Iso)}: for all $D,\mathfrak c,\delta,B$, inert $p\ge5$, $\lambda$ and field isomorphisms $\iota\colon\C\to\C_p$
  extending $\lambda$, there are a $\C$-linear automorphism $g_1$ of $\C^{4\times\mathcal T}$ and $e\in\C$, a Frobenius lift
  $\sigma$ of $\reckKurp$ (an endomorphism with $|\sigma(a)-a^p|\le p^{-1}$ for $|a|\le1$), a $\sigma$-unitary matrix
  $\gamma$ over $\reckKurp$, a matrix $B'$ over $\C_p$ and a conjugation $\sigma_p$ such that: $g_1$ is an isometry from the
  hermitian form of four copies of the lattice of $B$ to the standard form, the relative indices between the standard
  lattice and the image under $g_1$ of four copies of the lattice of $B$ are, in one order, prime to $p$ and, in the
  other, non-zero, $|\iota(e)|\le1$, the fourfold tensor
  power of the theta function of $B$ is $e$ times the theta kernel attached to $g_1$; $B'\gamma$ is the $\iota$-image of
  $g_1^{-1}$ (entrywise, in the sense of the development); and the rows of $B'$ form a local orthonormal frame of
  $Y(\mathrm{std}_1,\mathrm{std}_1)$ at $p$ for $\sigma_p$.
\end{enumerate}
\formallabel{Formalism.L2.KerInert.OneCopyOrth}
\formallabel{Formalism.L2.KerInert.IsomONB}
\formallabel{Formalism.L2.KerInert.WeakApproxAndOrthonormalBases}
\end{definition}

\begin{lemma}[One copy of the theta function]\label{pf:reckernel:oneCopy}
Assume \textup{(GM)} and \textup{(Iso)}, and assume that the translates $U_s\theta^0_{N,N'}$ by $p'$-torsion points $s$ for
the standard lattices have algebraic Taylor coefficients. Then for all $D,\mathfrak c,\delta,B$, inert $p\ge5$ and
$\lambda$ there are a conjugation $\sigma$ and a local orthonormal frame $v$ of $Y(\mathrm{std}_1,\mathrm{std}_1)$ at $p$
for $\sigma$ such that $(U_s\theta^0_{1,1})_\lambda$ is integral in the coordinates $v$ for every $p'$-torsion point $s$ for
$(\mathrm{std}_1,\mathrm{std}_1)$.
\formallabel{Formalism.L2.KerInert.oneCopyLoc_of_isom}
\formallabel{Formalism.L2.KerInert.oneCopyTh1S_of_isom}
\formallabel{Formalism.L2.KerInert.oneCopyTh1S_of}
\formallabel{Formalism.L2.KerInert.oneCopyLoc_ofS}
\formallabel{Formalism.L2.KerInert.change_of_parameters_const_scale}
\end{lemma}

\begin{proof}
\emph{Step 1.} Extend $\lambda$ to a field isomorphism $\iota\colon\C\to\C_p$
(\lref{Formalism.L2.KerInert.lattice_extension}: $\Qbar$ is countable, and $\C$, $\C_p$ are algebraically closed of
characteristic $0$ and of the same cardinality).

\emph{Step 2 (Lubin--Tate coordinates).} The development provides a non-zero $c\in\C$ and a Lubin--Tate normalised short
model $(g_2,g_3,s_2,\eps)$ over $\reckKurp$ of the scaled lattice $c\Lambda$ ($\Lambda$ the lattice of $B$) with
$|g_2|,|g_3|\le1$, $|g_2^3-27g_3^2|=1$, $|\eps|=1$ and integral theta series, which moreover satisfies a norm relation for
every Frobenius lift (\lref{Formalism.L2.KerInert.scale_lattice_bound}: existence of such a scaled model). Combined
with the data $g_1,e,\sigma,\gamma,B',\sigma_p$ of \textup{(Iso)} for $\iota$, the integrality theorem
\lref{Formalism.L2.KerInert.scale_integrality} (\emph{for such data and every $t$ with $Nt$ in the lattice for some
$N\ge1$ prime to $p$, the series $\exp\bigl(\sum_\tau\phi(X_\tau)\bigr)$ times the $\iota$-image of the Taylor expansion of
the translate of the theta function of $B$ by $t$, read in the coordinates $(\iota c)^{-1}B'$, is integral}, where
$\phi$ is a series without constant term determined by the model) gives the following: with the conjugation $\sigma_p$ and
the local orthonormal frame $v$ formed by the rows of $B'$, for every such $t$ the series
\begin{equation}\label{pf:reckernel:eq-LTint}
\exp\Bigl(\sum_{\tau}\phi(X_\tau)\Bigr)\cdot\mathrm{Read}_{\kappa v}(G_t)\qquad(\kappa=(\iota c)^{-1})
\end{equation}
is integral, where $G_t$ is the $\iota$-image of the Taylor expansion of the translate by $t$ and $\mathrm{Read}_{\kappa v}$
substitutes $X_\tau\mapsto\sum_j\kappa v_j(\tau)\,\lambda_p(Y_j)$ with the Lubin--Tate logarithm
$\lambda_p(X)=\sum_{j\ge0}X^{p^{2j}}/p^j$.

\emph{Step 3 (coordinates).} For a $p'$-torsion point $s$ for $(\mathrm{std}_1,\mathrm{std}_1)$ one has
$(U_s\theta^0_{1,1})_\lambda=G_s$ and $Ns$ lies in the lattice for some $N\ge1$ prime to $p$
(\lref{Formalism.L2.KerInert.coordTie_of_algZero}: the coordinate identification, which uses the algebraicity
hypothesis).

\emph{Step 4 (models).} For every model $C$ there are series $h$ and $U$ with $h(0)=0$, $h$ integral, $|h'(0)|=1$, $U$
integral, such that $u_C^{-1}\ell_C(h(X))=\kappa\lambda_p(X)$ and $\Phi_C(h(X))^{-1}=U(X)\exp(\phi(X))$
(\lref{Formalism.L2.KerInert.MModelS.model_series_algebraic}: comparison of the models with the Lubin--Tate
normalised model).

\emph{Step 5 (change of parameters).} Let $C$ be a model, $h,U$ as in Step 4, and let $S$ be the substitution
$Y_j\mapsto h(Y_j)$. By Step 4, $S$ transforms $G_s^{C,v}$ into $\mathrm{Read}_{\kappa v}(G_s)$ and
$\prod_j\Phi_C(Y_j)^{-1}$ into $\prod_jU(Y_j)\cdot\exp\bigl(\sum_j\phi(Y_j)\bigr)$. Hence the image under $S$ of
$G_s^{C,v}\prod_j\Phi_C(Y_j)^{-1}$ is the product of the integral series $\prod_jU(Y_j)$ and the integral series
\eqref{pf:reckernel:eq-LTint}, so it is integral. Since $h(0)=0$, $h$ is integral and $h'(0)$ is a unit, $h$ has a
compositional inverse $g$ with integral coefficients; the substitution $Y_j\mapsto g(Y_j)$ inverts $S$ and preserves
integrality (it substitutes integral series without constant term). Therefore
$G_s^{C,v}\prod_j\Phi_C(Y_j)^{-1}$ is integral, which is the assertion.
\end{proof}

\begin{theorem}[Reduction of \textup{(Int)}]\label{pf:reckernel:intRed}
Assume \textup{(GM)}, \textup{(WA)}, \textup{(ONB)} and \textup{(Iso)}. Then \textup{(Int)} holds.
\formallabel{Formalism.L2.KerInert.integral_inert_of_weakApprox}
\formallabel{Formalism.L2.KerInert.intInert_of}
\formallabel{Formalism.L2.KerInert.thFourInt_of}
\formallabel{Formalism.L2.KerInert.weakApprox_of}
\end{theorem}

\begin{proof}
We use the following properties of integrality, proved in the development for all data as above.
\begin{enumerate}[label=(I\arabic*)]
\item\label{pf:reckernel:I1} $(L,L')$-integral series are stable under linear combinations with $p$-integral coefficients
  (\lref{Formalism.L2.KerInert.combination_fact}).
\item\label{pf:reckernel:I2} If $L\sim_pL_1$ and $L'\sim_pL_1'$, then local bases of $Y(L,L')$ at $p$ are local bases of
  $Y(L_1,L_1')$ at $p$, and $p'$-torsion points for $(L,L')$ are $p'$-torsion points for $(L_1,L_1')$
  (\lref{Formalism.L2.KerInert.pairing_fact}); consequently $(L_1,L_1')$-integrality implies $(L,L')$-integrality.
\item\label{pf:reckernel:I3} For unitary $\gamma,\gamma'$ and lattices $M,M'$: the pull-back $t^*$ of a $p'$-torsion point $t$
  for $(\gamma M,\gamma'M')$ is $p'$-torsion for $(M,M')$ (\lref{Formalism.L2.KerInert.pullback_torsion}); and if
  $U_{t^*}G$ has algebraic Taylor coefficients and $(U_{t^*}G)_\lambda$ is $(M,M')$-integral, then
  $(U_t\,\gamma^*G)_\lambda$ is $(\gamma M,\gamma'M')$-integral, where $\gamma^*G(z)=G({}^t\bar\gamma'z\gamma)$
  (\lref{Formalism.L2.KerInert.pullback_fact}).
\item\label{pf:reckernel:I4} The translates $U_te_\chi(N,N',M,M')$ for admitted $(M,M')$ with $M\sim_p\mathrm{std}_N$,
  $M'\sim_p\mathrm{std}_{N'}$ and $p'$-torsion $t$, and the translates $U_s\theta^0_{N,N'}$ for $p'$-torsion $s$ of the
  standard lattices, have algebraic Taylor coefficients (\lref{Formalism.L2.KerInert.AlgP.algebraic_fact_kernels} and \lref{Formalism.L2.KerInert.AlgP.zero_value_algebraic}).
\item\label{pf:reckernel:I5} $G\mapsto G_\lambda$ is linear for algebraic coefficients:
  $(\sum_kc_kG_k)_\lambda=\sum_k\lambda(c_k)(G_k)_\lambda$ if the $c_k$ are algebraic and the $G_k$ have algebraic Taylor
  coefficients (\lref{Formalism.L2.KerInert.power_series_linearity}).
\item\label{pf:reckernel:I6} For admitted $(M,M')$ with $M\sim_p\mathrm{std}_N$, $M'\sim_p\mathrm{std}_{N'}$ and $p'$-torsion
  $t$, there are algebraic $c_k$ with $|\lambda(c_k)|\le1$ and $p'$-torsion points $s_k$ for the standard lattices with
  $U_te_\chi(N,N',M,M')=\sum_kc_kU_{s_k}\theta^0_{N,N'}$ (\lref{Formalism.L2.KerInert.pairs_standard_at_p}).
\item\label{pf:reckernel:I7} If a series is integral for some orthonormal local basis $b$ of $Y(L,L')$ at $p$ (that is,
  $F^{C,b}\prod\Phi_C^{-1}$ is integral for every model $C$), then it is $(L,L')$-integral
  (\lref{Formalism.L2.KerInert.basisInv_of_loc}, from the unit substitution theorem
  \lref{Formalism.L2.KerInert.unit_substitution_local}: \emph{an orthonormal change of the coordinates
  $\sum_jM_{ij}\ell_C(Y_j)$ with $p$-integral unitary $M$ is realised by an integral substitution $Y\mapsto S(Y)$, and
  transforms $\prod\Phi_C^{-1}$ by an integral factor}); the same holds for local orthonormal frames
  (\lref{Formalism.L2.KerInert.locInv_of}).
\item\label{pf:reckernel:I8} If for every $p'$-torsion $s$ of the standard lattices of rank $(2,2)$ the function
  $U_s\theta^0_{2,2}$ has algebraic Taylor coefficients and $(U_s\theta^0_{2,2})_\lambda$ is
  $(\mathrm{std}_2,\mathrm{std}_2)$-integral and integral for some orthonormal local basis, then for all
  even $N,N'$ and $p'$-torsion $s$ of the standard lattices, $(U_s\theta^0_{N,N'})_\lambda$ is integral for some orthonormal
  local basis (\lref{Formalism.L2.KerInert.tensor_matrix_units_fact}: the standard theta function is a tensor
  product of copies of rank $(2,2)$, and matrix units give compatible orthonormal bases).
\end{enumerate}

\emph{Rank $(2,2)$.} Let $\sigma$ and $v$ be as in \cref{pf:reckernel:oneCopy} (its algebraicity hypothesis is
\ref{pf:reckernel:I4}). By the four-products identity (\lref{Formalism.L2.KerInert.four_products_from_one}:
\emph{if all $(U_s\theta^0_{1,1})_\lambda$ are algebraic and integral in the coordinates of a local orthonormal frame $v$ of
$Y(\mathrm{std}_1,\mathrm{std}_1)$ for $\sigma$, then there is a local orthonormal frame $v_4$ of
$Y(\mathrm{std}_2,\mathrm{std}_2)$ for $\sigma$ in whose coordinates all $(U_s\theta^0_{2,2})_\lambda$ are integral}) we obtain
$v_4$. By \textup{(ONB)} and \lref{Formalism.L2.KerInert.MMOnb.globONB_of} (\emph{there is an orthonormal local basis
$b$ of $Y(\mathrm{std}_2,\mathrm{std}_2)$ at $p$ whose image $(\lambda(\tau(b_i)))$ is a local orthonormal frame for every
conjugation}) and \ref{pf:reckernel:I7} for frames, every $(U_s\theta^0_{2,2})_\lambda$ is integral for the orthonormal local
basis $b$, hence $(\mathrm{std}_2,\mathrm{std}_2)$-integral by \ref{pf:reckernel:I7}; it is algebraic by \ref{pf:reckernel:I4}.
By \ref{pf:reckernel:I8} and \ref{pf:reckernel:I7}, $(U_s\theta^0_{N,N'})_\lambda$ is
$(\mathrm{std}_N,\mathrm{std}_{N'})$-integral for all even $N,N'$ and all $p'$-torsion $s$ of the standard lattices.

\emph{Standard lattices at $p$.} Let $(M,M')$ be admitted with $M\sim_p\mathrm{std}_N$, $M'\sim_p\mathrm{std}_{N'}$, and $t$
$p'$-torsion for $(M,M')$. By \ref{pf:reckernel:I6}, \ref{pf:reckernel:I4} and \ref{pf:reckernel:I5},
$(U_te_\chi(N,N',M,M'))_\lambda=\sum_k\lambda(c_k)(U_{s_k}\theta^0_{N,N'})_\lambda$ with $|\lambda(c_k)|\le1$. Each term is
$(\mathrm{std}_N,\mathrm{std}_{N'})$-integral, hence $(M,M')$-integral by \ref{pf:reckernel:I2}, and the sum is $(M,M')$-integral
by \ref{pf:reckernel:I1}.

\emph{General admitted lattices.} Let $(N,N',L,L')$ be admitted and $t$ $p'$-torsion for $(L,L')$. By \textup{(WA)}, applied
to $L$ and to $L'$ (\lref{Formalism.L2.KerInert.wa_main}: the form of weak approximation used for one lattice),
there are unitary matrices $\gamma,\gamma'$ over $E$ and admitted $(M,M')$ with $M\sim_p\mathrm{std}_N$,
$M'\sim_p\mathrm{std}_{N'}$, $L=\gamma M$ and $L'=\gamma'M'$. The functions $e_\chi$ transform by
$e_\chi(N,N',\gamma M,\gamma'M')=\gamma^*e_\chi(N,N',M,M')$
(\lref{Formalism.L2.PairingIndexBounds.Eq.SqG.eulerProduct_eq_of_curve_ideal}: the transformation law of the
functions $e_\chi$ under unitary matrices, part of the identities for curve ideals). By \ref{pf:reckernel:I3}, $t^*$ is
$p'$-torsion for $(M,M')$, $U_{t^*}e_\chi(N,N',M,M')$ is algebraic by \ref{pf:reckernel:I4}, and its $\lambda$-image is
$(M,M')$-integral by the previous paragraph; so $(U_te_\chi(N,N',L,L'))_\lambda$ is $(L,L')$-integral.

\emph{Conclusion.} Let $F=\sum_kr_k(U_{t_k}e_\chi(N,N',L_k,L'_k))_\lambda\in\mathrm{Comb}(L,L')$. By \ref{pf:reckernel:I2},
$t_k$ is $p'$-torsion for $(L_k,L'_k)$, so each term is $(L_k,L'_k)$-integral by the previous paragraph, hence
$(L,L')$-integral by \ref{pf:reckernel:I2}, and $F$ is $(L,L')$-integral by \ref{pf:reckernel:I1}.
\end{proof}

\begin{theorem}[Integrality at inert primes]\label{pf:reckernel:intInert}
\textup{(GM)} implies \textup{(Int)}.
\formallabel{GlobalTheory.intInert_of_complexMultiplication_condition}
\formallabel{Formalism.L2.KerInert.kernels_integral_inert}
\formallabel{Formalism.L2.KerInert.integral_inert_of_approx_basis}
\end{theorem}

\begin{proof}
By \cref{pf:reckernel:intRed} it suffices to establish \textup{(WA)}, \textup{(ONB)} and \textup{(Iso)}. \textup{(WA)} is
\lref{Formalism.L2.KerInert.weak_approximation_places} (weak approximation for hermitian lattices at the places above
$p$, deduced from the simultaneous approximation of elements of a number field at finitely many places). \textup{(ONB)} is
\lref{Formalism.L2.KerInert.one_theta_orthonormal} (an orthogonal local basis of $Y(\mathrm{std}_1,\mathrm{std}_1)$ at $p$
with rational norms prime to $p$). \textup{(Iso)} follows from the existence of local orthonormal bases of the hom-lattices
(\lref{Formalism.L2.KerInert.local_orthonormal_basis}) by
\lref{Formalism.L2.KerInert.isomONB_ofLB} (construction of the isomorphism with the standard data and of the frame).
The hypothesis \textup{(GM)} enters through the existence of Lubin--Tate normalised models in Step 2 of the proof of
\cref{pf:reckernel:oneCopy} and of good models over the Witt vectors in the unit substitution theorem used in
\ref{pf:reckernel:I7}.
\end{proof}

\subsubsection{Admissibility below maximal level}
The last statement concerns the theta-series data attached to a \emph{site} of a datum: a word $u$, a tuple $\mu$, a
place $w$ of the coefficient field over a prime $p$, a filling with levels and a chain, kernel data at $\mu$, and a
system of choices; from these the development forms a system of lattices and levels $X$ over the Witt vectors
together with pending entries.

\begin{definition}\label{pf:reckernel:def-adm}
\textup{(Adm)}: for every datum $D$, every site at a place $w$ over $p$, every system of choices for it, every $p$-adic
curve datum $g$ with a Lubin--Tate isomorphism and every CM base and character, if $g$ is in position for the base and
character (\cref{pf:reckernel:def-Mod}(2)), the theta-series data hypotheses below maximal level hold, the hypotheses
on the auxiliary series $\Xi$ and the stabiliser law hold, and the ties hold at every step $j<m$, then the pending
entries of the site are admitted.
\formallabel{GlobalTheory.TwV.Lt.AdmAtEnt2}
\end{definition}

\subsubsection{The first kernel law and admissibility}\label{pf:reckernel:sec-admproof}
We describe the objects of (Adm) in more detail. A site of a datum $D$ consists of a word $u$, a tuple $\mu$ and a place $w$
of the coefficient field over a prime $p$ satisfying the index condition of the construction (which implies that $p\ge5$
and that $p$ is inert in $K$), together with fills, levels and a chain at $(u,w)$ and kernel data at $\mu$. With standard
choices at $p$ (a uniformiser of the local ring, the ramification index, bounds for the $p$-parts of the orders of the
stabilisers, and bases $\omega$) the development forms from a site a \emph{formal-group system} $X$ at $p$: a $p$-adic
field $K_X$ with valuation ring, a set of places $v$ with integers $n'_v$, lattice tuples at the levels $j\le m$, level
groups $\Gamma$ acting on them, spaces of forms and the chain maps $K_j$. Given moreover an embedding $\mathcal E$ of the coefficient
field, a $p$-adic curve datum $g$ with a Lubin--Tate isomorphism, a bijection between the pairs $(v,s)$ ($s<n'_v$) and
$\mathcal T$, frame data (lattices $\ker_N(j,L)$ attached to the lattice tuples, with frames $\varphi$), data
$\mathfrak c,\delta,B,\chi,\lambda$, a finite set of places, an integer $c_l$, auxiliary functions $\Xi_0$ and characters
$\mu_T$, the development defines \emph{pending entries} $Y$: for each step $j<m$, representative lattice tuples $L_0$ at
level $j$, lattice tuples $L'$ at level $j+1$ and frames $\varphi$, a theta series $\Theta_{j}(L_0,L',\varphi)$ over $\C_p$,
obtained from the $p$-adic expansion of a \emph{kernel theta function} (a combination of the functions $e_\chi$ for the
lattice pair $(\ker_N(j,L_0),\ker_N(j+1,L'))$ with coefficients built from $\Xi_0$, the places and $c_l$), multiplied by a
scalar and a power of $p$, transported by the Lubin--Tate isomorphism of $g$ and by the frame substitution, together with
auxiliary series $\beta$, $\Theta_V$ and $\Theta_F$. The pending entries are \emph{admitted} if the following thirteen
clauses hold:
\begin{enumerate}[label=(A\arabic*)]
\item\label{pf:reckernel:A1} the Dwork lemma holds at every prime;
\item\label{pf:reckernel:A2} the local extension condition for the series $\Theta_V$ at every place $v$, every $j<m$, all
  lattice tuples and frames;
\item\label{pf:reckernel:A3} the extension generator condition for the theta series $\Theta_j(L_0,L',\varphi)$ (all
  $j<m$, lattice tuples and frames);
\item\label{pf:reckernel:A4} the standard property of the series $\beta$;
\item\label{pf:reckernel:A5} non-negativity of the valuations of the coefficients of the column projections of the theta
  series;
\item\label{pf:reckernel:A6} the chain maps $K_j$ carry the rational spaces at level $j$ into those at level $j+1$;
\item\label{pf:reckernel:A7} archimedean signature matching (an expansion equality for $j+1<m$ and an expansion
  decomposition for $j+1=m$);
\item\label{pf:reckernel:A8} the explicit formula expressing $\Theta_F$ through the auxiliary maps $\Xi_v$, powers of two
  units $u,u'$ of valuation $0$ and a power of $p$ given by volumes;
\item\label{pf:reckernel:A9} \emph{Minkowski counting}: the kernels are invariant under the level groups, that is,
  $\ker_j(\gamma L,\chi(\gamma)\gamma G)=\ker_j(L,G)$;
\item\label{pf:reckernel:A10} an adjointness identity between the sums over up-edges and over down-edges of the
  representatives, weighted by the inverse orders of the stabilisers;
\item\label{pf:reckernel:A11} lattice inclusions for the kernel operators, applied to series with valuation bounds, at
  every modification of a lattice tuple by edges;
\item\label{pf:reckernel:A12} two selection conditions on the system $X$;
\item\label{pf:reckernel:A13} the lower bound $0$ for the local systems at every step $j<m$ and place $v$.
\end{enumerate}

\begin{theorem}[First kernel law at inert primes]\label{pf:reckernel:firstLaw}
Assume \textup{(GM)}. Let $D$ be a datum, $p\ge5$ a prime, $X$ a formal-group system at $p$ with an embedding $\mathcal E$, finitely
many classes and finite stabilisers at every level, $g$ a $p$-adic curve datum with a Lubin--Tate isomorphism, and let
frame data, $\mathfrak c,\delta,B,\chi,\lambda$, a finite set of places, $c_l$, $\Xi_0$, $\mu_T$ and embeddings $\rho_v$ of
$E$ into the local fields at the places $v$ of $X$ be given. Assume that $g$ is in position at $(\mathfrak c,\delta)$ and
defined over the Witt vectors, $p$ is inert in $K$, $\mathfrak c$ is admissible, $\delta$ is a normalising root, $\chi$ is a
CM character, $\lambda=\lambda_g$, $\Lambda_g$ is the lattice of $B$, and the following \emph{frame hypotheses} hold:
\begin{enumerate}[label=(\roman*)]
\item the pairs $(\ker_N(j,L_0),\ker_N(j+1,L'))$ are admitted for all $j<m$, representatives $L_0$ and lattice tuples $L'$;
\item $c_l>0$ and $p\nmid c_l$;
\item the values of $\Xi_0$ at representatives are integral outside a finite set of primes not containing $p$;
\item the scalars of the kernel theta functions are algebraic, and multiplied by the corresponding power of $p$ their
  $\lambda$-images have absolute value $\le1$;
\item $\lambda\circ\tau_{(v,s)}$ coincides with the twisted embedding attached to $(v,s)$ composed with $\rho_v$ (with complex
  conjugation for odd $s$);
\item the frame matrices are invertible;
\item for every orthonormal local basis $b$ at $p$ of the hom-lattice of the kernel lattices, the matrices $\rho_v(b_i)$ map
  the column spans of the frames into each other in the prescribed way.
\end{enumerate}
Then the pending entries satisfy \ref{pf:reckernel:A2}, \ref{pf:reckernel:A3} and \ref{pf:reckernel:A5}.
\formallabel{Formalism.L2.KerInert.KGLt.first_kernel_law_global_small}
\formallabel{Formalism.L2.KerInert.KGLt.first_kernel_small_of_frame_twist}
\formallabel{Formalism.L2.KerInert.KGLt.first_kernel_small_of_frame}
\formallabel{Formalism.L2.KerInert.KGLt.TranslatesOfInertKernelAtPlace}
\end{theorem}

\begin{proof}
Replacing $\Xi_0$ by its restriction to representatives (and $0$ elsewhere) does not change the theta series at
representatives and preserves (iii); so we may assume that (iii) holds everywhere. By construction, the theta family
is $\Theta=\mathrm{sub}_S(\mathrm{sc}(F))$, where $F_j(L_0,L')=p^{e_j}\,(s_j\,G_j)_\lambda$ with the kernel theta functions $G_j$,
their scalars $s_j$ and exponents $e_j$, $\mathrm{sc}$ is the scaling by the constant of the Lubin--Tate isomorphism and
$\mathrm{sub}_S$ the linear substitution by the frame generators $S$. We verify the hypotheses of the representative form
of the first kernel law (\lref{Formalism.L2.KerInert.RepLt.K1_representatives_LT_of_dictionary}: \emph{if the
standard property of $\beta$, the combination property of the family $F$ and the frame hypotheses for $\mathrm{sc}(F)$
and $S$ hold, then (Dict) and (Int) imply the extension generator condition for $\Theta_j(L_0,L',\varphi)$ at all
representatives}).
\begin{itemize}
\item The series $\beta$ attached to $g$ has the standard property, since $g$ is in position and defined over the Witt
  vectors at an inert $p\ge5$ (\lref{Formalism.L2.KerInert.beProp_ofA}).
\item Combination property: every $(G_j)_\lambda$ is an integral combination of $p$-adic expansions of translates of
  $e_\chi$ by $p'$-torsion points for the kernel lattices, which are admitted by (i)
  (\lref{Formalism.L2.KerInert.RepLt.combAtLt_of_dataRec}, using the algebraicity of the translates for the given
  $B$ and $\chi$, \lref{Formalism.L2.KerInert.AlgTP.algT_of}, which is the statement of \textup{(Alg)} for these data,
  and the data records given by (ii), (iii)); by (iv) this persists after
  multiplication by $p^{e_j}\lambda(s_j)$ (\lref{Formalism.L2.KerInert.RepLt.combAtLt_scalar}).
\item Frame hypotheses: orthonormal local bases of the hom-lattices exist
  (\lref{Formalism.L2.KerInert.RepLt.onbExLt_of}, from the global orthonormal basis of rank four and weak
  approximation, \lref{Formalism.L2.KerInert.orthBasisExists_of}); the trace condition holds by the rank ties of (i)
  and the trace formula over $\mathcal T$ (\lref{Formalism.L2.KerInert.RepLt.lamTraceLt_of}); and with (v), the
  Frobenius twist of the frames (\lref{Formalism.L2.KerInert.RepFrob.frobTw_of}), (vi) and (vii), the frame
  hypotheses follow (\lref{Formalism.L2.KerInert.RepLt.frameHypLt_of_lat}).
\end{itemize}
Under \textup{(GM)}, (Dict) and (Int) hold (\lref{Formalism.L2.KerInert.Rep.curve_dictionary}: the curve
dictionary under (GM); \cref{pf:reckernel:intInert}). Hence the extension generator condition holds at all
representatives, which is equivalent to \ref{pf:reckernel:A3}
(\lref{Formalism.L2.PairingIndexBounds.PowerSeriesNecessity.adm_two_pendant_iff_representative}: the condition at
representatives is equivalent to the condition for all lattice tuples, by the equivariance of the construction).
Then \ref{pf:reckernel:A2} follows from \ref{pf:reckernel:A3}
(\lref{Formalism.L2.KerInert.admissible_second_of_variant}), and \ref{pf:reckernel:A5} follows from
\ref{pf:reckernel:A3} together with \ref{pf:reckernel:A13}, which holds at inert primes for curves in position
(\lref{Formalism.L2.KerInert.admissible_third_of_second_eleventh} and \lref{Formalism.L2.KerInert.admissible_eleventh_inert}).
\end{proof}

\begin{theorem}[Admissibility from the first kernel law]\label{pf:reckernel:admLaw}
Assume that the conclusion of \cref{pf:reckernel:firstLaw} holds for all data satisfying its hypotheses. Then
\textup{(Adm)} holds.
\formallabel{Formalism.L2.KerInert.AdmLoc2.admissible_entries_of_kernel_law}
\formallabel{Formalism.L2.KerInert.AdmLoc2.admissible_place_of_kernel_law}
\formallabel{Formalism.L2.KerInert.AdmLoc2.admissible_general_of_kernel_law}
\formallabel{Formalism.L2.KerInert.AdmLoc2.admissible_stabilizer_of_law}
\end{theorem}

\begin{proof}
Let a site, standard choices, $\mathcal E$, $g$ and the remaining data be as in (Adm), with $g$ in position for the base and the
character, the theta-series data hypotheses below maximal level (which comprise the frame hypotheses (i)--(vii) of
\cref{pf:reckernel:firstLaw}), the hypothesis on $\Xi$, the stabiliser law, and ties at every step $j<m$. We verify
\ref{pf:reckernel:A1}--\ref{pf:reckernel:A13} for the pending entries $Y$.
\begin{itemize}
\item \ref{pf:reckernel:A1}: the Dwork lemma is proved in the development
  (\lref{Formalism.MainNotions.FormalGroupSystem.Pf.positivity_condition_holds}).
\item \ref{pf:reckernel:A2}, \ref{pf:reckernel:A3}, \ref{pf:reckernel:A5}: unpacking "in position" and the
  theta-series data hypotheses gives exactly the hypotheses of \cref{pf:reckernel:firstLaw}.
\item \ref{pf:reckernel:A4}: \lref{Formalism.L2.KerInert.admissible_second_inert} (the standard property of $\beta$
  for curves in position at inert primes).
\item \ref{pf:reckernel:A6}: \lref{Formalism.L2.KerInert.SecondKernelInstances.admissible_third_instance} (rationality
  of the chain maps at the site).
\item \ref{pf:reckernel:A7}: if $m=0$ there is nothing to prove. If $m>0$, the ties at the steps $j<m-1$, together with a
  coupled section at the top step $m-1$ and its compatibility and Hecke condition, which are derived from the tie at
  $m-1$ and from \ref{pf:reckernel:A9}
  (\lref{Formalism.L2.KerInert.SecondKernelCoupled.coupled_section_exists}, \lref{Formalism.L2.KerInert.SecondKernelCoupled.coupled_compatibility_at_place}, and \lref{Formalism.L2.KerInert.SecondKernelCoupled.hecke_coupled_entry}),
  give the signature matching (\lref{Formalism.L2.KerInert.SecondKernelCoupled.admissible_fourth_instance}).
\item \ref{pf:reckernel:A8}: the hypothesis on $\Xi$ says that additive maps with the required properties exist, which is
  equivalent to the formula \ref{pf:reckernel:A8} (\lref{Formalism.L2.G.Pend.invariance_iff_additive_maps_exist}).
\item \ref{pf:reckernel:A9}: the stabiliser law for the theta series implies the corresponding law for the normalised
  geometric series (\lref{Formalism.L2.G.Pend.kapLaw_of_stabLaw}), and that law implies the invariance of the kernels
  under the level groups (\lref{Formalism.L2.G.Pend.ker_law}, for forms in the admissible subspace).
\item \ref{pf:reckernel:A10}: \lref{Formalism.L2.G.Pend.up_down_adjoint} (adjointness of the up- and down-edge
  sums).
\item \ref{pf:reckernel:A11}: \lref{Formalism.L2.PairingIndexBounds.Ent.pending_tuple_admissible} (lattice
  inclusions for pending tuples, proved from the shape of the lattices at the middle levels).
\item \ref{pf:reckernel:A12}: \lref{Formalism.L2.KerInert.admissible_all_levels} (the selection conditions hold for
  the system of every site).
\item \ref{pf:reckernel:A13}: \lref{Formalism.L2.KerInert.admissible_eleventh_inert} (the lower bound at inert primes
  for curves in position).
\end{itemize}
Hence $Y$ is admitted.
\end{proof}

\begin{theorem}[Admissibility]\label{pf:reckernel:adm}
\textup{(GM)} implies \textup{(Adm)}.
\formallabel{GlobalTheory.TwV.Lt.inert_theta_admissibility_of_complex_multiplication}
\formallabel{Formalism.L2.KerInert.AdmLoc2.admissible_entries_small_chain}
\end{theorem}

\begin{proof}
By \cref{pf:reckernel:firstLaw}, \textup{(GM)} implies the first kernel law, and by \cref{pf:reckernel:admLaw} the first
kernel law implies \textup{(Adm)}.
\end{proof}

\begin{theorem}\label{pf:reckernel:curveEx}
\textup{(GM)} implies \textup{(Ex)}.
\formallabel{GlobalTheory.TwV.Lt.exists_genusOne_formal_iso_of_unramified_good_reduction}
\end{theorem}

\begin{proof}
By \cref{pf:reckernel:cm3}, \textup{(LT)} holds; apply \cref{pf:reckernel:modEx}.
\end{proof}

% ---------------------------------------------------------------------------
\subsection{The twisted-system row and the kernel-condition row}\label{pf:reckernel:sec-rows}

\begin{definition}[Local models at a constructed point]\label{pf:reckernel:def-modelsAt}
For a datum $D$ and an object $O$ over $D$, $\mathrm{LM}(D,O)$ means: if the lattice data of $O$ are tied and
$2r_G\le k$ (where $r_G$ is the rank parameter of the size data of $O$ and $k$ the rank of $D$), then the local systems
of $O$ admit models in the sense of the development (the twisted forms at level $0$ agree with those at $\mu_0$, and at
every used tuple, word and place outside the twisting set there are levels and a chain whose site is admissible and
compared with the lattice components).
\formallabel{GlobalTheory.Svd.ModelsAt}
\end{definition}

\subsubsection{Joints and admitted systems}\label{pf:reckernel:sec-joints}
Let $D$ be a datum and $O$ an object over $D$, with extra data $\mathrm{ex}_O$ (twisted forms, lattice comparison, kernel
data), and let $t$ be a tied kernel entry over $O$. For a word $u$, a tuple $\mu$ and a place $w$ over $p$ satisfying the
index condition and lying outside the twisting set of $\mu$, a \emph{joint} of $t$ at $(u,\mu,w)$ consists of levels for
the fills of the lattice instance at $(u,w)$, a comparison input and a comparison hypothesis for the twisted chain at
$\mu$, labels, and a compatible system of chain charts, adelic data, pins and ties along the twisted chain, identified at
level $0$ with the member chain of $u$, together with the local identifications at $w$ (the group compatibility of the
lattice comparison, the sharpness of the forms, the coprimality of $\Nm\mathfrak c$ and $p$, the splitting and lattice
conditions of the local charts, and the coefficient-field condition), all in the sense of the development. The joint
determines a chain $\mathrm{ch}(H)$ for the levels. We say that $t$ has \emph{joints at $\mu$} if it has a joint at
$(u,\mu,w)$ for every $u$ and every $w$ as above. For levels $\mathrm{Lv}$ and a chain $\mathrm{Ch}$ at $(u,w)$:
\begin{itemize}
\item $\mathrm{Ch}$ is \emph{compared} with the lattice components at $(u,\mu,w)$ if there are semilinear maps $c_j$ from the
  lattices of the twisted forms at $\mu$ to the level spaces which agree at level $0$ with the lattice comparison, are
  compatible with the chain maps, and satisfy the rationality conditions below level $m$ and a valuation condition at
  level $m$;
\item the site $(u,w,\mathrm{Lv},\mathrm{Ch})$ is \emph{admissible} if for some standard choices at $p$ and some bounds
  $d_G$ with $d_G(0)=0$ and $d_G(j)\le v_p(D_\Gamma(u))$ the formal-group system $X$ of the site carries admitted pending
  entries (\ref{pf:reckernel:A1}--\ref{pf:reckernel:A13}).
\end{itemize}
In these terms, $\mathrm{LM}(D,O)$ (\cref{pf:reckernel:def-modelsAt}) says: if the lattice data of $O$ are tied and
$2r_G\le k$, then the twisted forms at level $0$ agree with those at $\mu_0$, and there is a choice of levels and chains
at all $(\mu,u,w)$ such that for used $\mu$ and $w$ outside the twisting set of $\mu$ the chain is compared with the
lattice components and the site is admissible.

\begin{theorem}[Joints carry admitted systems]\label{pf:reckernel:jointSite}
Assume \textup{(Adm)} and \textup{(Ex)}. Let $t$ be a tied kernel entry over $O$, and let $J$ be a joint of $t$ at
$(u,\mu,w)$ with levels $\mathrm{Lv}$ and comparison hypothesis $H$, where $\mu$ is used. Assume that the stabilisers at
level $0$ are trivial and $2r_G\le k$. Then the chain $\mathrm{ch}(H)$ is compared with the lattice components at
$(u,\mu,w)$, and the site $(u,w,\mathrm{Lv},\mathrm{ch}(H))$ is admissible.
\formallabel{GlobalTheory.T.jointSite_chain_compared_carries_admittedSystem}
\formallabel{GlobalTheory.TwV.site_cmpLite_siteAdm}
\formallabel{GlobalTheory.TwV.admissible_at_each_site}
\formallabel{GlobalTheory.T.jointData_at_p_and_adelicDatum_nonempty}
\formallabel{GlobalTheory.TwV.joint_at_nonempty}
\end{theorem}

\begin{proof}
\emph{Comparison.} The kernel data at $\mu$ attached to the labels and the comparison hypothesis of $J$ consist of
comparison maps with the four required properties (agreement at level $0$, compatibility with the chain maps,
rationality, valuations), by the construction of the kernel data at $\mu$ from comparison inputs in the development.
Hence $\mathrm{ch}(H)$ is compared with the lattice components.

\emph{Standard choices.} Since the lattice data are tied and $2r_G\le k$, the fills of the lattice instance satisfy the
standard conditions at every place $v$ of the site (\lref{GlobalTheory.TwU.std_at_served}). We take the uniformiser,
ramification index and bases of the development, and the bounds $d_G=d_{G,T}$ attached to the twisted chain; then
$d_G(0)=0$ and $d_G(j)\le v_p(D_\Gamma(u))$ for all $j$ (\lref{GlobalTheory.TwV.Lt.dGT_spec} and \lref{GlobalTheory.TwV.Lt.dGT_zero}).
Let $X$ be the resulting formal-group system.

\emph{Data at $p$.} The joint yields an embedding $\mathcal E$ of the coefficient field into the local field of the site and an
embedding $\lambda\colon\Qbar\to\C_p$ compatible with $\mathcal E$ on the coefficient field
(\lref{GlobalTheory.TwV.exists_cEmbC_lam}). The index condition gives $p\ge5$ and $p$ inert in $K$
(\lref{GlobalTheory.TwU.XdInert.gidx_facts}), and the joint contains the coprimality of $\Nm\mathfrak c$ and $p$; so
\textup{(Ex)}, applied to the CM base and the character of $t$, provides a $p$-adic curve datum $g$ with a Lubin--Tate
isomorphism which is in position for $(B,\chi,\lambda)$. Auxiliary functions $\Xi$ and characters $\mu_T$ satisfying the
adelic character laws at the joint exist (\lref{GlobalTheory.T.exists_adelicDatum_with_laws_at_joint}: the adelic
chain datum of the joint kernel satisfies the character laws), and the level groups of $X$ have finitely many classes and
finite stabilisers. These data form a \emph{site datum} at $p$, which determines pending entries $Y$ for $X$.

\emph{Admissibility.} We apply \textup{(Adm)} to the site, the standard choices, the site datum and the kernel data at
$\mu$; its hypotheses are verified by the following properties of site data, all proved in the development.
\begin{enumerate}
\item The lattice tuples of the joint satisfy the curve dictionary at $p$
  (\lref{GlobalTheory.T.latticeTuple_joint_dictionary}).
\item Given the dictionary: the pairs of kernel lattices are admitted
  (\lref{GlobalTheory.T.kernelSite_latticeTuplePairs_admissible}); the place $w$ induces the position embeddings
  (\lref{GlobalTheory.T.site_place_induces_position_embeddings}); the arithmetic group moves the frames
  (\lref{GlobalTheory.T.kernelSite_arithmeticGroup_moves_frames}); the theta series do not depend on the data away
  from the place over $p$ (\lref{GlobalTheory.T.kernelSite_series_independent_away_from_place_over_p}); and the
  series scalars, multiplied by the appropriate power of $p$, are integral
  (\lref{GlobalTheory.T.kernelSite_seriesScalar_pow_integral}).
\item From these, the theta-series data hypotheses below maximal level hold
  (\lref{GlobalTheory.T.allSites_k1_of_dict_pair_pos_scal}), and the kernel law, the stabiliser law in fixed
  coordinates and the hypothesis on $\Xi$ hold
  (\lref{GlobalTheory.T.kernelSite_law_and_fixedCoordinateStabiliserLaw}).
\item The ties at every step $j<m$: extend $\lambda$ to an isomorphism $\iota\colon\C\to\C_p$
  (\lref{GlobalTheory.ThLoc.place_extends_to_iso}); since the coefficient field consists of algebraic numbers,
  $\iota$ agrees with $\lambda$, hence with $\mathcal E$, on it. The stabilisers of the kernel classes are tied because the stabilisers at level $0$ are trivial
  and $\mu$ is used (\lref{Formalism.MainNotions.G.KAat.tieStab_kaAtAny}); the series reading of the entries and the
  series compatibility hold (\lref{Formalism.L2.KerInert.SecondKernelProjection.series_reading_entries} and \lref{Formalism.L2.KerInert.SecondKernelSeries.series_compatibility_of});
  the sharpness of the forms is part of $J$; and the frame values at step $j$ are given by a scalar $c_K$ with
  $|c_K|=1$ (\lref{GlobalTheory.T.kernelSite_frameValues_step_zero} and \lref{GlobalTheory.T.kernelSite_frameValues_positive_steps}),
  so that $|c_K|\le1$, $|c_K^{-1}|\le1$ and $c_K\ne0$.
\end{enumerate}
Hence \textup{(Adm)} shows that the pending entries $Y$ are admitted, and together with the standard choices and the bounds
$d_G$ this is the admissibility of the site.
\end{proof}

\begin{theorem}[Theta local models at joints]\label{pf:reckernel:modelsJoints}
Assume \textup{(Adm)} and \textup{(Ex)}. Let $t$ be a tied kernel entry over $O$ with joints at every used tuple, assume
that the twisted forms at level $0$ agree with those at $\mu_0$, that the stabilisers at level $0$ are trivial, and that
$2r_G\le k$. Then the local systems of $O$ admit models, i.e.\ the conclusion of $\mathrm{LM}(D,O)$ holds.
\formallabel{GlobalTheory.T.thetaLocalModels_of_splitUnramified_joints}
\formallabel{GlobalTheory.TwV.modelsOK_of_joints}
\end{theorem}

\begin{proof}
The first condition is a hypothesis. For each $(\mu,u,w)$ choose levels and a chain as follows: if $\mu$ is used and $w$
lies outside the twisting set of $\mu$, take the levels of a joint $J$ of $t$ at $(u,\mu,w)$ and the chain
$\mathrm{ch}(H)$; otherwise take the levels and the chain of the lattice instance. For used $\mu$ and $w$ outside the
twisting set, \cref{pf:reckernel:jointSite} applied to $J$ shows that the chosen chain is compared with the lattice
components and that the site is admissible.
\end{proof}

\begin{theorem}[Joints for standard kernels]\label{pf:reckernel:jointsStd}
Let $(w,Z,O)$ be a column choice, a later stage and an object over its datum such that $(w,Z,O)$ has a part, and assume
that the lattice data of $O$ are tied. Then there is a tied kernel entry over $O$ with joints at every tuple, and the
twisted forms of $O$ at level $0$ agree with those at $\mu_0$.
\formallabel{GlobalTheory.Svd.jointsAt_all_std}
\formallabel{GlobalTheory.T.kernel_relation_true_columns_joints_all_tuples}
\formallabel{GlobalTheory.Pdv.oneTie_jointsAt_all}
\formallabel{GlobalTheory.TwU.AllTuples.joints_all_true}
\formallabel{GlobalTheory.T.trueColumn_joints_hold_all_tuples}
\end{theorem}

\begin{proof}
The part of $(w,Z,O)$ contains a tie of $O$ with the true columns of the column choice: the extra data of $O$ are those of
the true-column construction (its twisted forms are the twisted forms $\mathrm{tw}$ of the construction, and the
parameter $\pi_T$ of its error data is that of the construction), together with a standard kernel entry $t$, the tied symmetric columns, their full
standard data, the group compatibility of the lattice comparison and a tuple tie; and the values of the character
$\chi_D$ are values of the coefficient field (\lref{GlobalTheory.Mem.ChiDVal.chiD_values_and_cc_top}). A standard
kernel entry can be replaced by one which is moreover standard-read and marked at the pinned classes
(\lref{GlobalTheory.Pdv.stdKer_stdRd_mark}). For such $t$ we show that $t$ has joints at every tuple $a$; the second
assertion holds because the twisted forms are those of the construction
(\lref{GlobalTheory.T.tw_e0_mu0_eq_eMu0_zero}: at the reference tuple the level-$0$ forms coincide).

Fix a tuple $a$. By \lref{GlobalTheory.T.trueColumn_exists_chainCharts_kernelData_forms}, for every word $u$ there
is a \emph{joint input} at $(u,a)$: chain charts at every stage of the twisted chain of $a$ which are good for the true
column choice (\lref{GlobalTheory.T.trueColumnChoice_goodCharts_every_stage}), adelic kernel data compatible with
the step charts, with pins and ties (\lref{GlobalTheory.T.trueColumn_kernelData_compatible_stepCharts}), the
identification of the base lattices with the standard lattices
(\lref{GlobalTheory.T.companionLattice_trueColumn_baseLattice_eq_standard}), the identifications at level $0$, the
rationality of the tuple kernels and the coefficient-field condition. From a joint input at $(u,a)$ for every $u$ and the
tie of the lattice data, joints are built at every place $w$ outside the twisting set
(\lref{GlobalTheory.T.jointsAtSplitUnramifiedPlaces_of_chainCharts}): one constructs comparison site data over
the chain charts (\lref{GlobalTheory.T.exists_comparisonSiteData_over_chainCharts}), verifies the comparison
statements at the site, which give the comparison hypothesis
(\lref{GlobalTheory.T.comparison_data_statements_at_site}), the statements at level $0$
(\lref{GlobalTheory.T.step_zero_level_statements_of_joint}), labels whose level-$0$ class map is the given one
(\lref{GlobalTheory.T.exists_labIn_c0_cls_eq_levelZeroClass}), the equality of the charts and twists of the
comparison site with those of the tuple chain (\lref{GlobalTheory.T.comparisonSite_charts_twists_eq_tupleChain}),
the local identifications at $w$ (\lref{GlobalTheory.T.site_localIdentifications_eq_stepZero}) and the sharpness of
the forms (\lref{GlobalTheory.FormsSharpTie.formsSharpAt_of_pins}). These are the components of a joint.
\end{proof}

\begin{theorem}[Local models]\label{pf:reckernel:localModels}
Assume \textup{(Adm)} and \textup{(Ex)}. Then $\mathrm{LM}(D_\beta,O_\beta)$ holds at every bundle point $\beta$ whose core is a
position with a part.
\formallabel{GlobalTheory.T.tuplesHaveLocalModels_at_constructedPoint_over_trueColumns}
\formallabel{GlobalTheory.Svd.modelsAt_bpt}
\formallabel{GlobalTheory.Svd.modelsAt_of_part}
\formallabel{GlobalTheory.T.tuplesHaveLocalModels_of_standardColumnObjects}
\end{theorem}

\begin{proof}
A position with a part comes from a column choice $w$ and a later stage $Z$ with a part; it suffices to prove the
statement for the object of the column choice and then to transport it to the object of the bundle point, which
differs only in division lattices and heights
(\lref{GlobalTheory.T.tuplesLocalModels_of_otherDivisionLatticesHeights}). For the column choice, the standard
kernel entry has \emph{joints} at all used tuples (\cref{pf:reckernel:jointsStd}), the stabilisers at
level $0$ are trivial over true columns (\lref{GlobalTheory.Pdv.stepZero_trivial_stabilisers_trueColumns}), and
the twisted forms at level $0$ agree with those at $\mu_0$. Given these, the theta local models at split unramified
joints (\cref{pf:reckernel:modelsJoints}) provide, at every used tuple $\mu$, word
$u$ and place $w$ outside the twisting set, levels and a chain whose site carries an admitted system compared with the
lattice components (\cref{pf:reckernel:jointSite}); this is where
(Adm) and (Ex) are used.
\end{proof}

\subsubsection{Lower bounds at the second real place}\label{pf:reckernel:sec-second}
Let $D$ be a datum with size data, an instance and a lattice instance, constants, error data and a choice of levels and
chains at all $(\mu,u,w)$; replacing the levels and chains of the lattice instance by the chosen ones at $\mu$ gives the
\emph{twisted lattice instances}. For a site at $(u,w)$ over $p$ write $\mathcal O_1(w)$ for the set of places of the
site which are images of the representatives of the first orbits, and for a place $v$ and $j\le m$ let $\Lambda_j(v,g)$
be the statement that all coefficients of the column projections at $v$ of the theta series of the pending entries at
step $j$, twisted by the exponential of the phase at $(j,v)$ and multiplied by an arbitrary integral series, have
valuation $\ge g$. The \emph{low term} of $(u,w)$ is
\[
\mathrm{low}(u,w)=\sum_{j}\bigl(F_j(w)-\mathrm{ES}(u,j,w)\bigr)+v_{D\Gamma}(u,w),
\]
the sum over the lifts $j$ of the size data, with the size functions $F_j$, the error sums $\mathrm{ES}$ and the
$p$-adic valuation $v_{D\Gamma}(u,w)$ of the stabiliser constant of $u$. The \emph{twisted system bound} (with constant
$1$) is the statement: if the error definitions hold, then for every used tuple $\mu$, every tie of the twisted lattice
instance at $\mu$, every word $u$ and every place $w$ with the index condition outside the twisting set of $\mu$, the
formal-group system of the site carries admitted pending entries for some standard choices and bounds $d_G$, and there are
real numbers $\phi(v,j)$ for $v\notin\mathcal O_1(w)$ with $\Lambda_j(v,\phi(v,j+1))$ for $j<m$ and
\begin{equation}\label{pf:reckernel:eq-low}
\mathrm{low}(u,w)\le\sum_{j=1}^{m}\ \sum_{v\notin\mathcal O_1(w)}\phi(v,j)-\sum_{j<m}d_G(j).
\end{equation}

\begin{theorem}[Lower bounds at the second real place]\label{pf:reckernel:secondPlace}
Assume that the lattice data are tied, the fills are complete, the size data are defined, the orbit letter bounds with
constants $(11,1)$ hold, the entry has the shape property, the constants are large for $(11,1)$, and that for every used
$\mu$, word $u$ and place $w$ with the index condition outside the twisting set of $\mu$ the chosen site is admissible.
Then the twisted system bound with constant $1$ holds.
\formallabel{GlobalTheory.T.local_systems_lower_bounds_second_real_place}
\formallabel{GlobalTheory.TwU.SiteSys.sysTw_sites}
\end{theorem}

\begin{proof}
Let $\mu,u,w$ be as in the twisted system bound, over the prime $p$. By hypothesis the site is admissible: there are
standard choices, bounds $d_G$ with $d_G(0)=0$ and $d_G(j)\le v_p(D_\Gamma(u))$, and admitted pending entries $Y$; we use
them. The places of the site off the first orbits have $n'_v=n_P(p)$, empty target sets, and positions over the second
real embedding: there is a map $r$ from places to the orbit index set such that the positions of $v\notin\mathcal O_1(w)$
are $(\mathrm{true},r(v))$ and its shifts, and such that for every function $f$ on the index set which is invariant under
the shift $t\mapsto t+h_P$ one has $\sum_{v\notin\mathcal O_1(w)}f(r(v))=\sum_{t}f(t)$, the sum over the representatives $t$
(\lref{GlobalTheory.SecondOrbits.places_off_first_orbits}). Put
\[
\phi(v,j)=\begin{cases}\mathrm{FOrb}(j,p,(\mathrm{true},r(v)))-\mathrm{err}(j,p,(\mathrm{true},r(v)))&\text{if }p\le X''(n),\\
0&\text{otherwise},\end{cases}
\]
with the orbit function $\mathrm{FOrb}$, the error function $\mathrm{err}$ and the threshold $X''(n)$ of the size data.

\emph{The bounds $\Lambda_j$.} Let $v\notin\mathcal O_1(w)$ and $j<m$. If $p\le X''(n)$, then
$\Lambda_j(v,\phi(v,j+1))$ is the local model bound at places over the second embedding
(\lref{GlobalTheory.T.local_model_bound_place_over_second_embedding}: \emph{under the hypotheses of the theorem and
$p\le X''(n)$, for every admitted system at the site whose series $\beta$ has the standard property, and every place $v$ with
$n'_v=n_P(p)$ whose positions are $(\mathrm{true},t)$ and its shifts, the bound $\Lambda_j(v,\mathrm{FOrb}-\mathrm{err})$ holds
for $j<m$}), applied with the standard property \ref{pf:reckernel:A4} of $Y$. If $p>X''(n)$, then $\Lambda_j(v,0)$ is
clause \ref{pf:reckernel:A13} of $Y$.

\emph{The inequality \eqref{pf:reckernel:eq-low}.} By \lref{GlobalTheory.T.low_term_le_second_real_place_bounds}
(\emph{if the size data are defined, $d_G(0)=0$ and $d_G(j)\le v_p(D_\Gamma(u))$, then $\mathrm{low}(u,w)$ is at most
$\sum_{j=1}^m\sum_tf_j(t)-\sum_{j<m}d_G(j)$, where $f_j(t)$ is $\mathrm{FOrb}(j,p,(\mathrm{true},t))-\mathrm{err}(j,p,(\mathrm{true},t))$
for $p\le X''(n)$ and $0$ otherwise, and $t$ runs over the representatives}). Since $w$ satisfies the index condition, $p$ is
not a signal prime, so the shift of the position $(\mathrm{true},t)$ is $(\mathrm{true},t+h_P)$; and $\mathrm{FOrb}$ and $\mathrm{err}$ are
invariant under shifts of positions (\lref{Formalism.BasicNotions.FOrb_step} and \lref{GlobalTheory.TwU.SiteSys.errF_posPl}).
Hence $f_j$ is shift invariant, and the reindexing property of $r$ gives
$\sum_tf_j(t)=\sum_{v\notin\mathcal O_1(w)}f_j(r(v))=\sum_{v\notin\mathcal O_1(w)}\phi(v,j)$, which proves
\eqref{pf:reckernel:eq-low}.
\end{proof}

\begin{theorem}[The twisted-system row]\label{pf:reckernel:rowTwU}
For every bundle type, slot $Q$ and linkage $\reckLk$ with
$\reckRowE\bigl(Q,\reckLk,(D,O,c,b)\mapsto\mathrm{LM}(D,O)\bigr)$, the row
$\reckRow\bigl(Q,\reckLk,(\mathcal A,\sigma)\mapsto\reckClosed(\mathcal A,\mathrm{TwU})\bigr)$ holds, where $\mathrm{TwU}$ is
the twisted-system condition of \cref{pf:reckernel:def-servH}.
\formallabel{GlobalTheory.Svd.rowTwU}
\formallabel{GlobalTheory.Proof.UnitaryWeightModels.rowTwU_of_models}
\end{theorem}

\begin{proof}
Take $E=\mathrm{LM}$ as the auxiliary property of the row. Let $\sigma$ be an $E$-selection for $\mathcal A$; we show
$\reckClosed(\mathcal A,\mathrm{TwU})$ with $\eps_1=1$ and the trivial threshold for every $\eps$. Let $F\in\mathcal A$. Its
selected bundles are built for a parameter linked to $F$ and satisfy $\mathrm{LM}$ at every member. The lattice data of the
members are tied, and $2r_G\le k$ holds for all large $x$
(\lref{GlobalTheory.TwU.two_rG_le_k_large}: a size inequality valid along every family); hence the local systems
of the members have models for large $x$. Moreover, for large $x$ the members carry an entry with the shape property
(\lref{GlobalTheory.ShapeRows.shape}), and the orbit letter bounds with constants $(11,1)$ hold
(\lref{Formalism.MainNotions.DIV.Bridge.DworkPowerSeriesWeights.split_bound_eleven_one}). From these data, the
models give the Weil--Jacquet--Shalika compatibility of the lattice components
(\lref{GlobalTheory.Proof.ColumnWeightComparison.cmpV_of_cmpLite}) and the lower bounds for the local systems at the
second real place yield the twisted system bound
(\cref{pf:reckernel:secondPlace}). Taking $x_0$ the maximum of the two
thresholds proves $\mathrm{TwU}(F)$.
\end{proof}

\subsubsection{The kernel conditions at true columns}\label{pf:reckernel:sec-truecols}
We use the notation of \cref{pf:reckernel:def-servH}(2). For a kernel entry $E$ and auxiliary data $(\epsilon,\beta,s)$
(a scalar, power series and signs at each place) the development forms the \emph{derived entry}
$E^\ast=E^\ast(\epsilon,\beta,s)$, whose theta series $\mathrm{Th}_{E^\ast}(u,\mu,w,i,c,c')$ are those of $E$. For a prime $p$
let $h_P(p)$ be the element of the group of the host given by the Frobenius at $p$, and $n_P(p)$ its order (twice the
order of $2h_P(p)$ if $p$ is a signal prime); $p$ is of \emph{residue class zero} if it lies in the inert class of the host,
and of \emph{class $C_s$} if moreover $h_P(p)\ne0$ has odd order. A place $w$ is \emph{level-free} for $(u,\mu)$ if it lies
outside the levels of the division data.

\begin{theorem}[Normalised theta and CM-curve inputs]\label{pf:reckernel:inputs}
\textup{(Alg)} implies the following statements, which together form the \emph{inputs} of the development:
\begin{enumerate}
\item the $j$-invariant of a lattice with complex multiplication by $\omega_d$ ($d$ squarefree, $d\equiv3\pmod 4$) is an
  algebraic integer;
\item (good ordinary models at $2$) for $d$ squarefree with $d\equiv7\pmod 8$, a lattice $L$ with complex multiplication
  by $\omega_d$ and an isomorphism $\iota\colon\C\to\C_2$, some non-zero multiple $cL$ is uniformised over $\C_2$ by a
  Weierstrass curve with integral coefficients and $|a_1|=1$;
\item (jets at odd primes) for self-dual lattices with complex multiplication, theta functions $\theta$ of the given
  normalisation, odd primes $p$ not dividing the discriminant, isomorphisms $\iota$ with integral curves and locally
  normalised $\theta$: all linear combinations, with integral coefficients, of translates of $\theta$ by torsion points of
  order prime to $p$, read in integral frames of multiplier vectors, have integral jets, and pair integrally with integral
  polynomials;
\item (jets at $2$) the same at $p=2$ for good ordinary models with $|a_1|=1$;
\item (jets at $3$) for $d\equiv7\pmod8$ and $d\equiv2\pmod3$ the analogous integrality of jets at $p=3$ for the theta
  functions translated by half-periods not killed by $\omega_d$, $\bar\omega_d$;
\item (shape at $d$) assuming the separation statement for embeddings of Galois number fields, at primes $p$ dividing the
  discriminant the jet of order $2$ of the logarithm is bounded by $|d|_p^{-1/2}$, and the linear combinations of
  item 3, read in eigenframes, factor as $\exp(-\mathrm{jet}\cdot\beta)$ times a series with integral jets;
\item \textup{(Alg)} itself; and
\item every embedding $\lambda\colon\Qbar\to\C_p$ extends to a field isomorphism $\C\to\C_p$.
\end{enumerate}
\formallabel{GlobalTheory.T.normalisedThetaCMCurveTheorems_of_algebraicTranslates}
\formallabel{GlobalTheory.ThLoc.inputs_of_algT}
\end{theorem}

\begin{proof}
(1) is \lref{GlobalTheory.ThLoc.cm1}, whose proof is as follows. Write $L=a(\Z\tau+\Z)$ with $\tau$ in the
upper half plane; then $j(L)=1728g_2(L)^3/\Delta(L)$ equals the modular function $j(\tau)$, and
$\omega_d(\Z\tau+\Z)\subseteq\Z\tau+\Z$ with $\omega_d\notin\Z$. A value $j(\tau)$ at a point $\tau$ with a non-integral
multiplier of $\Z\tau+\Z$ is an algebraic integer (\lref{GlobalTheory.isIntegral_modularJ_of_mul_mem}: the classical
integrality of singular moduli \cite[Ch.~II, \S6]{Silverman1994}). (2) is
\lref{GlobalTheory.ThLoc.goodOrdinaryModelAtTwo_cmLattice_seven_mod_eight} (good ordinary models at $2$ for
$d\equiv7\pmod8$). (3), (4) and (6) are the three components of \lref{GlobalTheory.ThLoc.locJets_all}; each is
deduced from four ingredients of the development (a statement on the fourth power of the standard theta function, a
transfer of Gram matrices, the composition of frames and the linearity of jets)
(\lref{GlobalTheory.ThLoc.LocJets.jetsOdd}, \lref{GlobalTheory.ThLoc.LocJets.jetsTwo}, and \lref{GlobalTheory.ThLoc.LocJets.shape}),
the case $p=2$ resting on the integrality of the theta operators for orthogonal Gauss models at $2$. (5) is
\lref{GlobalTheory.ThLoc.thetaZJetsThree}, deduced from an integrality statement for exponentials of translates.
(7) is the hypothesis, and (8) is \lref{GlobalTheory.ThLoc.place_extends_to_iso}.
\end{proof}

\begin{theorem}[Kernel class sums carry the forms]\label{pf:reckernel:kFi}
Let $t$ be a tied kernel entry over an object $O$ which has the column-side properties and whose twisted forms are pinned
(for every used $\mu$ there is stabiliser data $G$ such that the forms of the twisted chain at $\mu$ are cut out by $G$ in the
sense below, and $G$ is given by the stabilisers of the pins). Then the chain maps $K_i$ ($i<m$) of the member chains and of
the twisted chains at used tuples carry the filtration step $F_i$ of the forms into $F_{i+1}$.
\formallabel{GlobalTheory.T.kernelClassSums_carry_chainComplexForms}
\formallabel{GlobalTheory.KFiClose.kFi_of_tie}
\end{theorem}

\begin{proof}
Stabiliser data $G$ for a twisted chain assign to each level $i$ and class $c$ a set $G(i,c)$ of pairs $(g,k)$ with
$g\in\GL_{N_i}(E)$ and $k\in\C$. The forms are \emph{cut out by $G$} if every form in $F_i$ ($i\le m$) satisfies the
valuation condition for the frame of the lift object at level $i$, and for $1\le i\le m$ a form $\varphi$ lies in $F_i$
if and only if every $\varphi(c)$ satisfies the valuation condition and the transformation law of the lift object for all
$(g,k)\in G(i,c)$. The \emph{term law} for $(g,k)$ at $(i,c,c')$ says: $\det g\ne0$, and there is $\kappa\in\C$ with
\[
\mathrm{Th}_i(c,c')\bigl((\tau(g)z_\tau)_\tau\bigr)=\kappa\,\mathrm{Th}_i(c,c')(z)\quad\text{for all }z,\qquad
k=\kappa\prod_\tau\tau(\det g)^{-s_\tau(i)},
\]
where $\mathrm{Th}_i(c,c')$ are the kernel theta functions of the lift object and $s_\tau(i)$ its divisor exponents.

Let $\mu$ be used and let $G$ be as in the hypothesis. The kernel theta functions of $t$ are entire
(\lref{GlobalTheory.KerTieEntire.kerTie_entire}); the unfolding conditions at the levels $i<m$ are part of $t$; and
the term law holds for all $(g,k)\in G(i+1,c')$, $i<m$ (\lref{GlobalTheory.T.classSum_term_law_stabiliser_target}:
\emph{for kernel entries with the column-side properties and stabiliser data given by the pins, every element of the
stabiliser sets satisfies the term law}). Under these conditions the kernel operator maps stabiliser-equivariant forms to
stabiliser-equivariant forms (\lref{GlobalTheory.T.kernelOperator_maps_stabiliserEquivariant_forms}: \emph{if the
forms are cut out by $G$, the unfolding conditions hold, the kernel theta functions are entire and satisfy the term law
for all elements of $G(i+1,c')$, then $K_i(F_i)\subseteq F_{i+1}$ for all $i<m$}); here $K_i$ is the class sum of the
kernel theta functions, and the term law makes each term transform under $G(i+1,c')$ as required. This proves the
assertion for the twisted chains. The reference tuple $\mu_0$ is used, and the statement for the member chain of $u$
follows from the statement for the twisted chain at $\mu_0$
(\lref{GlobalTheory.KFiMu0.memberForms_stable_of_reference_tuple}).
\end{proof}

\begin{theorem}[Orbit bound at inert places]\label{pf:reckernel:orbitBound}
Let $E^\ast$ be a derived entry with the shape property such that $\mathrm{row}(u,i,(\mathrm{false},t))_0-\mathrm{row}(u,i,(\mathrm{false},t))_1\le2a$ for
all $u$, $i<m$ and top positions $t$, let the constants satisfy $\mathrm{dl}_1\ge5$, $\mathrm{dl}_2\ge4$, assume that
$|\epsilon(w)|=1$ and that the series $\beta(w,\cdot)$ are standard at every place $w$ of residue class zero with
$h_P$ of even order or of signal class, and that all signs are $-1$. Then for every place $w$ over $p$ of residue class zero
with $h_P(p)\ne0$ and $p\le X''(n)$, every $u$, $i<m$ and representative $t_0$: whenever the local orbit series $\mathrm{Th}$ at
$(u,i,(\mathrm{false},t_0))$ has an orbit theta expansion with $\epsilon(w)$ and sign $-1$ (its $\epsilon(w)$-twist is an
integral series times the exponential of $-\sum\beta$ composed with the jet logarithm), all coefficients of its column projections have $p$-adic valuation at
least $\max(g_0(i+1,t_0),0)$, where $g_0$ is the orbit threshold of the size data.
\formallabel{GlobalTheory.T.orbitBound_standardSeries_inertPlaces}
\formallabel{GlobalTheory.OrbitBoundStdRec.orbit_bound_rec}
\end{theorem}

\begin{proof}
Let $w$, $u$, $i$, $t_0$ and $\mathrm{Th}$ be as in the statement. Since $p$ is of residue class zero, it is not a signal
prime, so $n_P(p)$ is the order of $h_P(p)$; the orbit data at $w$ give that this order is even. Hence $|\epsilon(w)|=1$ and
$\beta(w,\cdot)$ is standard, and the sign is $-1$. With these:
\begin{enumerate}
\item all coefficients of the column projections of $\mathrm{Th}$ have valuation $\ge0$
  (\lref{GlobalTheory.T.standardOrbitSeries_valuationBound_zero});
\item they have valuation $\ge\gamma$, where $\gamma$ is the \emph{orbit gain} of the local orbit datum
  (\lref{GlobalTheory.T.kvalGe_orbitGainWithoutOperator}: the orbit gain bound without operator for standard
  series);
\item $g_0(i+1,t_0)\le\gamma$, by the shape property, the row bound, $\mathrm{dl}_1\ge5$, $\mathrm{dl}_2\ge4$ and $p\le X''(n)$
  (\lref{GlobalTheory.KerStd.WordFree.wordFreeNumber_le_orbitBound}: the word-free number is at most the orbit
  gain).
\end{enumerate}
If $g_0(i+1,t_0)\le0$, the bound $\max(g_0,0)=0$ is (1). Otherwise $\max(g_0,0)=g_0\le\gamma$ by (3), and a lower bound
$\gamma$ for all valuations implies the lower bound $g_0$; so the claim follows from (2).
\end{proof}

\begin{theorem}[Frobenius conjugation conditions]\label{pf:reckernel:frobConj}
Assume \textup{(Alg)} and the host statement at the primes $p\ge5$ of residue class zero or of signal class (for all
admissible $\mathfrak c$, normalising roots $\delta$, CM bases $B$, CM characters $\chi$ and embeddings $\lambda$). Let $t$
be a tied kernel entry which satisfies the kernel-entry hypotheses and has adapted bases, and assume that the split bound
of the host is at least $3$. Then for all auxiliary data $(\epsilon,\beta,s)$ the derived entry satisfies:
\begin{enumerate}
\item at level-free places $w$ over primes $p$ of residue class zero with $h_P(p)=0$, for $i<m$ with $N_iN_{i+1}$ even,
  the theta kernels satisfy the \emph{theta expansion condition at the real places};
\item at level-free places $w$ over primes $p$ of class $C_s$, for $i<m$ with $N_iN_{i+1}$ even, the theta kernels satisfy
  the \emph{even conditions at all real places}.
\end{enumerate}
\formallabel{GlobalTheory.T.frobeniusConjugation_conditions_adaptedKernelDatum}
\formallabel{GlobalTheory.Ggeo.frobeniusConjugation_conditions_adaptedKernelDatum_alt}
\end{theorem}

\begin{proof}
\emph{Reduction to standard series.} Let $w$ be level-free for $(u,\mu)$ over a prime $p$ of residue class zero, $i<m$ and
$c,c'$ classes. Since the split bound is at least $3$, every prime of residue class zero is $\ge5$
(\lref{Formalism.MainNotions.DIV.Inh.orbit_admissible_of_split_bound}). The adapted bases of $t$ at $w$ give, for the
lattices $L=L_i(c)$ and $L'=L_{i+1}(c')$ of the lift object, matrices $P,P'$ which are \emph{standard bases} of $L$, $L'$ at
$\lambda=\lambda_w$ (\lref{GlobalTheory.T.isStandardBasesLatticeAtPlace_of_adapted_bases}: bases adapted to all
$\tau$ and dual under the conjugation pairing are standard). By \textup{(Alg)}, $(N_i,N_{i+1},L,L')$ is admitted and the
$p$-adic expansion $G$ of the kernel theta function $\mathrm{Th}_i(c,c')$ lies in $\mathrm{Comb}(L,L')$
(\lref{GlobalTheory.Ggeo.KerComb.comb_of_tied}). The host statement holds at $(B,\chi,p,\lambda)$, and the theta series
of the derived entry is the standard series of $G$ with respect to $P,P'$
(\lref{GlobalTheory.T.atLin_adaptedCoef_eq_standardSeries}). Hence every property of $p$-adic series which holds
for all standard series $\mathrm{std}_{P,P'}(G)$ of elements $G\in\mathrm{Comb}(L,L')$, for all admitted $(L,L')$ with standard
bases, whenever the host statement holds and $p\ge5$, holds for $\mathrm{Th}_{E^\ast}(u,\mu,w,i,c,c')$.

(1) For $h_P(p)=0$, the theta expansion condition at the real places holds for such standard series
(\lref{GlobalTheory.T.thetaExpansionRealPlaces_stdSer_of_adm'}).
(2) A prime of class $C_s$ is of residue class zero, hence not a signal prime, so $n_P(p)$ is the order of $h_P(p)$, which is
odd; the even conditions at all real places hold for standard series when $n_P(p)$ is odd
(\lref{UpperTheory.lattice_pairs_odd_degree_isometric}: lattice pairs of odd degree are isometric, which yields the
even conditions).
\end{proof}

\begin{theorem}[Kernel conditions at true columns]\label{pf:reckernel:atTrue}
Assume \textup{(Alg)}, \textup{(Mod)}, \textup{(Dict)} and \textup{(Int)}. Let $I$ be a true-column construction with its
pins, $l$ a law, $Z$ a later stage, $l_0$ a law at level $0$, $\mathrm{ex}$ extra data over the datum of $(l,Z)$ and $t$ a tied
kernel entry over $\mathrm{ex}$. Assume that the columns are normalised and have their jets, that the twisted forms of
$\mathrm{ex}$ are those of the construction, that $t$ is a standard kernel of the construction whose entry is of column type,
that $t$ is adapted and top-pinned, that the split bound of the host is at least $3$, that $\mathrm{dl}_1\ge5$ and
$\mathrm{dl}_2\ge4$, and that every derived entry has the shape property and the row bound
$\mathrm{row}_0-\mathrm{row}_1\le2a$ at top positions. Then the kernel conditions $\mathrm{KC}(\mathrm{ex},E_t)$ hold for the entry $E_t$ of
$t$.
\formallabel{GlobalTheory.PdvClose.at_of_true_columns}
\end{theorem}

\begin{proof}
\emph{Duties and hypotheses.} The duties of the standard kernel comparison hold for $t$
(\lref{GlobalTheory.T.conditions_at_standardKernelComparison_trueColumns}): $t$ has the column-side properties, its
entry is of column type, the rank condition holds, and there is a tied kernel entry $t'$ with the column-side properties
and pinned forms. From the column-side properties and the top pin, $t$ satisfies the kernel-entry hypotheses
(\lref{GlobalTheory.T.kernelEntries_hyp_of_columnSideProperties}). By \cref{pf:reckernel:inputs}, \textup{(Alg)} gives
the inputs; and \textup{(Mod)}, \textup{(Dict)}, \textup{(Int)} give the host statement at the primes $p\ge5$ of residue
class zero or of signal class (\lref{GlobalTheory.KerStd.HostInert.hostAt_inert}).

\emph{Auxiliary data.} From \textup{(Alg)}, the host statement, the hypotheses and the split bound we obtain $\epsilon$ and
$\beta$ such that $|\epsilon(w)|=1$ and the series $\beta(w,\cdot)$ are standard at the relevant places, and such that the
even conditions at the inert places hold for the derived entry with these data and signs $-1$
(\lref{GlobalTheory.T.admissibleKernelSeries_inertPrime_nonConjugateFrobenius}: admissible kernel series at inert
primes with non-conjugate Frobenius). We take the derived entry $E^\ast=E_t^\ast(\epsilon,\beta,-1)$.

\emph{The components of $\mathrm{KC}$.} The pair scalars are units at the level-free places
(\lref{GlobalTheory.T.kernelPairScalars_isUnit_trueColumns}); the transition maps are rational, by the column type
of the entry and the rank condition (\lref{GlobalTheory.PdvKm.kmapRat_kerEntC'}); the chain maps preserve the
filtrations of the forms, by \cref{pf:reckernel:kFi} applied to $t'$; and the stabilisers at level $0$ are trivial, by
the level data of the construction (\lref{GlobalTheory.T.stab0_rec7_of_levelData}). For $E^\ast$: the even conditions
at the inert places were obtained above; the even conditions at the places of residue class zero with $h_P=0$ and of class
$C_s$ are \cref{pf:reckernel:frobConj}; the orbit bound is \cref{pf:reckernel:orbitBound} (with the shape property and the
row bound of the hypothesis, and the properties of $\epsilon$, $\beta$); jet integrality and the shape condition at $d$
follow from the inputs and the kernel-entry hypotheses
(\lref{GlobalTheory.T.jetIntegralityBis_ofKer_of_hyp} and \lref{GlobalTheory.T.kernelDatum_thetaSeries_shape_at_d_adaptedBases});
and the shape inequalities are part of the hypothesis.
\end{proof}

\begin{theorem}[The kernel-condition row]\label{pf:reckernel:pdv}
Assume \textup{(Alg)}, \textup{(Mod)}, \textup{(Dict)} and \textup{(Int)}. For every bundle type, slot $Q$, linkage
$\reckLk$ and assignment $t$ of tied kernel entries with
$\reckRowE\bigl(Q,\reckLk,(D,O,c,b)\mapsto\mathrm{Std}(D,O,t(b))\bigr)$, the row
$\reckRow\bigl(Q,\reckLk,\reckAtAll(b\mapsto\mathrm{KC}(\mathrm{ex}_O,E_{t(b)}))\bigr)$ holds.
\formallabel{GlobalTheory.Pdv.rowAt_of_stdTie}
\formallabel{GlobalTheory.PdvTie.at_of_stdTie}
\formallabel{GlobalTheory.T.kernelDatumConditions_at_standardKernelComparison}
\end{theorem}

\begin{proof}
There is a function $x_0(\eps)$ such that for every family of exponent $\eps$, every $x\ge x_0(\eps)$ and every member
at $x$, all entries of the member satisfy the shape inequalities and a row bound, and the split bound of the host is at
least $3$ (\lref{GlobalTheory.ShapeUniform.shape_uniform}). Take as auxiliary property of the row
$E(\kappa,D,O,c,b)$: ``$x_0(\eps_\kappa)$ is at most the scale of $O$, and $t(b)$ is standard-tied'', where
$\eps_\kappa$ is the exponent of $\kappa$. Its filter statement is the conjunction of a trivial one and the hypothesis.
Let $\sigma$ be an $E$-selection, $F$ a member of exponent $\eps$ and $x>x_1(F)$; the linked parameter has exponent
$\eps$, and the scale of the member at $x$ is $x$, so $x\ge x_0(\eps)$. It remains to prove the kernel conditions at a
standard-tied kernel entry satisfying these bounds and the constant inequalities of the family. A standard tie of
$(D,O,t)$ consists of a true-column construction $I$ with its pins, a law $l$, a later stage $Z$ and a law $l_0$ at level
$0$ such that $D$ is the datum of $(l,Z)$, the construction is twisted and has its jets, the twisted forms of $O$ are those
of the construction, $t$ is a standard kernel of the construction whose entry is of column type, and $t$ is adapted and
top-pinned; the columns of a twisted construction are normalised
(\lref{GlobalTheory.Pdv.colNorm_of_tw}). Hence \cref{pf:reckernel:atTrue} gives the kernel conditions.
\end{proof}

% ---------------------------------------------------------------------------
\subsection{Remarks on the formal statements}\label{pf:reckernel:sec-remarks}

\begin{enumerate}
\item \emph{Inputs.} The only literature input of \cref{pf:reckernel:main} is global Artin reciprocity
  (\cref{pf:reckernel:artinLit}); the hypothesis $\sfname{RecTExists}$ only serves to form the type of parameter sets.
  Chebotarev's theorem, the facts on modular forms \ref{pf:reckernel:M1}--\ref{pf:reckernel:M6} and the formal-group
  facts \ref{pf:reckernel:F1}--\ref{pf:reckernel:F3} are proved in the development; Dirichlet's theorem on primes in
  arithmetic progressions is taken from Mathlib \cite{Mathlib2020}.
\item \emph{Normalisation of reciprocity.} The formal statement of global Artin reciprocity uses the geometric
  normalisation (uniformisers correspond to inverse Frobenius elements); the ideal-class Artin symbol of
  \cref{pf:reckernel:artinSymbol} is accordingly defined as the inverse of the map induced by the idele class map.
\item \emph{The field $\reckKur$.} All statements (GM), (LT), (TI) refer to the completion $\reckKur$ of the maximal
  unramified extension (\cref{pf:reckernel:rem-kur}). Classically (GM) holds over $\Q_\ell^{\mathrm{ur}}$ itself; the
  formal statement is the (formally weaker) statement over the completion.
\item \emph{Unused hypotheses.} The proof of (TI) does not use the complex multiplication of $L$, nor $\ell\nmid d$. The
  hypothesis of \cref{pf:reckernel:cheb} that the inertia elements above places outside $S$ lie in $H$ is not used
  in its proof. The deduction of (Dict) in \cref{pf:reckernel:threeFromCM} takes (GM), (LT) and (TI) as hypotheses
  but does not use (LT), and the deduction of (Mod) in \cref{pf:reckernel:modEx} does not use (TI).
\item \emph{Two proofs of (LT).} The development contains two independent formal proofs of (LT)
  (\cref{pf:reckernel:cm3}); one is used for (Mod) and (Dict), the other for (Ex).
\item \emph{Algebraicity of $g_2$.} In \cref{pf:reckernel:g2Alg} the algebraicity of $g_2(L)$ for a CM lattice with
  $\Delta(L)=1$ is deduced from (GM); it also follows directly from the integrality of $j(L)$, since
  $1728\,g_2(L)^3=j(L)\Delta(L)=j(L)$.
\item \emph{The CM fields.} The statements (GM), (LT), (TI) are stated for every squarefree $d\equiv3\pmod 4$, whereas
  the host fields of the construction have $d$ prime with $d\equiv3\pmod4$.
\item \emph{Closedness in the service row.} In \cref{pf:reckernel:rowTwU} the closed statement is proved with
  $\eps_1=1$ and the trivial threshold: for an $\mathrm{LM}$-selection the twisted-system condition holds for every
  member of the class, eventually in $x$, without any largeness of the family constants. Similarly, the kernel
  conditions in the service row hold at all selected bundles (\cref{pf:reckernel:pdv}); the closed form used in the
  package is obtained by reusing the thresholds of the first component (\cref{pf:reckernel:kerServ}).
\item \emph{Unused hypotheses in the rows.} The error definitions and the tie of the twisted lattice instance, which
  are hypotheses of the twisted system bound, are not used in the proof of \cref{pf:reckernel:secondPlace}; there the tie
  of the original lattice instance is used. In \cref{pf:reckernel:frobConj} the evenness of $N_iN_{i+1}$, a hypothesis
  of both conditions, is not used.
\item \emph{Proved and deferred parts.} The reduction of the service row to the three CM statements and the six
  statements on theta data (\cref{pf:reckernel:sec-row}), the statement (GM) (\cref{pf:reckernel:sec-cm2}), the statement
  (TI) (with the comparison of the chord-tangent law and the formal group law in \cref{pf:reckernel:sec-fgl}; the
  parametrisation of the formal group and the existence of the formal group law are cited), and (Mod), (Ex) from (GM)
  and (LT) are proved here. For (Int), (Adm), the local models, the lower bounds at the second real place and the kernel
  conditions at true columns the proofs are given at the level of their main steps; the lemmas they rest on, which
  concern the specific theta kernels, the Lubin--Tate normalised models and the formal-group systems of the
  construction, are cited by their formal names. For (LT) and (Dict) the structure of the proof is given.
\end{enumerate}
